\documentclass[10.5pt,twoside,openright]{book}
\usepackage[a5paper,inner=19mm,outer=15mm,top=20mm,bottom=22mm]{geometry}
\usepackage{fontspec}
\setmainfont{texgyrepagella}[Extension=.otf,UprightFont=*-regular,BoldFont=*-bold,ItalicFont=*-italic,BoldItalicFont=*-bolditalic,Ligatures=TeX]
\usepackage{xcolor}\usepackage{fancyhdr}\usepackage{titlesec}\usepackage{tcolorbox}\usepackage{setspace}
\definecolor{accent}{RGB}{154,52,18}\definecolor{gris}{gray}{0.45}
\setstretch{1.08}\setlength{\parindent}{1.3em}\setlength{\parskip}{0pt}\widowpenalty=10000\clubpenalty=10000\emergencystretch=2em
\pagestyle{fancy}\fancyhf{}\renewcommand{\headrulewidth}{0pt}
\fancyhead[LE]{\small\itshape Tu n’es pas nul}\fancyhead[RO]{\small\itshape\rightmark}\fancyfoot[C]{\small\thepage}
\fancypagestyle{plain}{\fancyhf{}\fancyfoot[C]{\small\thepage}\renewcommand{\headrulewidth}{0pt}}
\newcommand{\asterisme}{\par\vspace{0.7em}\centerline{\color{gris}*\quad*\quad*}\vspace{0.7em}\par\noindent}
\newcommand{\marqueur}[1]{\begingroup\color{accent}\fbox{\parbox{0.92\linewidth}{\small\sffamily\bfseries [#1]}}\endgroup}
\newcommand{\soustitre}[1]{\par\vspace{1.1em}\noindent{\bfseries\color{accent}#1}\par\vspace{0.4em}\noindent}
\newenvironment{pourcahier}[1]{\par\vspace{0.9em}\begin{tcolorbox}[colback=black!4,colframe=accent,boxrule=0pt,leftrule=2.2pt,arc=0pt,left=8pt,right=8pt,top=5pt,bottom=5pt]\small\textbf{#1}\par\smallskip}{\end{tcolorbox}\par}
\newenvironment{avertissement}{\par\small\itshape}{\par}
\newcommand{\partie}[2]{\cleardoublepage\thispagestyle{empty}\vspace*{\fill}\begin{center}{\large\color{gris}Partie #2}\\[1.2em]{\Huge\bfseries #1}\end{center}\vspace*{\fill}\cleardoublepage\markboth{}{}}
\newcommand{\chapitre}[2]{\cleardoublepage\markright{#2}\thispagestyle{plain}\vspace*{2.2cm}{\noindent\color{accent}\fontsize{54}{54}\selectfont\bfseries #1}\par\vspace{0.5em}{\noindent\LARGE\bfseries #2\par}\vspace{1.6em}\noindent\ignorespaces}
\newcommand{\chapitresn}[1]{\cleardoublepage\markright{#1}\thispagestyle{plain}\vspace*{2.8cm}{\noindent\LARGE\bfseries #1\par}\vspace{1.6em}\noindent\ignorespaces}
\begin{document}
\frontmatter

\begin{titlepage}\centering\vspace*{3cm}
{\Huge\bfseries Tu n’es pas nul\par}\vspace{0.6cm}
{\large\itshape Ce que tes notes ne prouvent pas\par}\vspace{2.2cm}
{\Large Ludovic A.\par}\vfill{\small manuscrit — version de travail\par}\end{titlepage}
\thispagestyle{empty}\cleardoublepage
\thispagestyle{empty}\vspace*{\fill}\begin{center}\itshape On a écrit un mot sur toi.\\Ce n’était pas ton nom.\end{center}\vspace*{\fill}\cleardoublepage

\mainmatter
%% 01-prologue-le-mot.md
\chapitresn{Prologue — Le mot}

À Ebolowa, en saison sèche, la chaleur de l’après-midi ne s’en va pas avec le soleil. Elle reste sous la tôle de la classe de cinquième, avec l’odeur de la craie et de la poussière rouge que les sandales ont rapportée de la cour. On s’assoit à trois par banc. Celui du milieu n’a jamais de place pour ses coudes ; il écrit en tenant son cahier un peu de biais, comme s’il voulait le cacher à quelqu’un. Ce jour-là, celui du milieu, c’est Rodrigue Essomba. Il a douze ans, une chemise d’uniforme repassée le matin même, et des mains qui ne savent pas où se mettre.

Le professeur de mathématiques rend les copies de la séquence. Il lui reste deux classes à voir avant le soir, et cela s’entend : il ne distribue pas, il lit. Le nom, la note, la copie posée sur le premier banc, d’où elle passe de main en main vers le fond. Les noms tombent dans l’ordre de la pile et non dans celui de l’alphabet, si bien que personne ne sait quand viendra le sien. « Abena : onze. Mbarga : neuf. Ndzié : quatorze, c’est bien. » Dans un coin du tableau noir, l’exercice de la veille survit à moitié, effacé par une manche.

« Essomba : quatre. » La copie claque sur le banc. Puis, sans lever les yeux, déjà tourné vers la suivante : « Toujours nul. »

Dans une classe, entre un mot et le rire qui le suit, il s’écoule un temps très court. Rodrigue connaît ce temps. Il rit le premier. Pas fort, juste assez pour que le rire des autres, quand il arrive, tombe sur un rire déjà commencé, et qu’on puisse croire qu’ils rient avec lui. Ça marche à moitié. Son voisin de gauche lui donne un coup de coude qui se veut amical. Le professeur dit « Mengue : douze », et la classe passe à autre chose, parce qu’une classe passe toujours à autre chose.

La copie arrive dans ses mains. Le quatre est entouré de rouge ; il ne le regarde pas longtemps. Il plie la feuille en quatre, la glisse dans son cahier de brouillon et, jusqu’à la sonnerie, recopie ce qui est au tableau avec une application d’élève modèle. À la récréation, il sort acheter des beignets. Le cahier reste sur le banc.

Le soir, à la maison, il le sort pour faire ses exercices. Sur la couverture, l’étiquette blanche à liseré bleu, celle qu’on achète par planches à la rentrée, porte ses trois lignes : « Nom : », « Classe : », « Matière : ». À la ligne du nom, quelqu’un a barré « Rodrigue Essomba » d’un trait appliqué, un seul, bien droit, et a écrit au-dessus, au stylo à bille, en capitales : NUL. Il ne sait pas qui. Ils étaient nombreux à pouvoir le faire, et il n’a pas envie de compter.

Il pourrait arracher l’étiquette. Elle se décolle mal, mais elle se décolle ; il resterait un rectangle gris de colle, un peu sale, et demain personne ne saurait. Il passe l’ongle sous le bord. Il s’arrête. Il ne sait pas très bien pourquoi il s’arrête, et je ne le sais pas mieux que lui. Il garde l’étiquette, comme on garde une cicatrice qu’on n’a pas choisie : parce qu’elle est à lui, maintenant, et qu’en l’arrachant il faudrait admettre qu’elle comptait.

\asterisme

Ce mot, tout le monde l’a entendu un jour, ou l’a dit. Dans une classe, dans une cour, à table, dans la bouche d’un parent épuisé ou d’un camarade qui voulait faire rire, et parfois dans la sienne, sur quelqu’un d’autre, sans y penser. Je ne crois pas qu’on y échappe tout à fait, ni d’un côté ni de l’autre.

En quatrième, j’ai raté une factorisation toute simple. Je n’étais pas un mauvais élève : je n’avais pas compris la leçon, et c’est tout ce que ce raté disait de moi. Il ne parlait pas de moi, il parlait d’une marche que je n’avais pas montée.

Qu’on l’ait reçu ou prononcé, la question reste la même : ce que ce mot fait, une fois dit, à celui qui l’emporte chez lui.

« J’ai eu quatre » et « je suis nul » ne disent pas la même chose, même quand ils parlent de la même copie. La première phrase raconte un événement. Elle a une date, un après-midi de saison sèche, et un sujet : quelques exercices sur les nombres relatifs. Elle a aussi ses conditions, la chaleur sous la tôle, ce qu’on avait compris ou non trois semaines plus tôt. Elle parle d’une heure dans la vie de quelqu’un. La seconde phrase décrit une personne entière, dans toutes les matières, pour toujours, et elle se passe de date et de sujet puisqu’elle prétend valoir partout.

Le mot fait passer de l’une à l’autre en une seconde. C’est là sa violence propre, et elle ne tient pas au volume de la voix : le professeur ne criait pas. Il a simplement ajouté un adjectif à un chiffre, et le chiffre, qui parlait d’une copie, s’est mis à parler d’un garçon. « Toujours » a fait le reste, en changeant un après-midi en habitude. Dans la bouche du professeur, le mot a duré moins longtemps qu’il n’en faut pour tourner une page. Chez Rodrigue, il a trouvé une chaise, et il s’est assis.

Le mot a un autre effet, plus discret, et je le crois plus grave : il clôt l’enquête avant qu’elle commence. Quand un élève a été déclaré nul, plus personne ne demande ce qui lui manque, puisque la réponse est déjà donnée, et que cette réponse, c’est lui. Or ce qui manque porte presque toujours un nom plus petit. Une leçon de l’an dernier qui n’est jamais passée. Des nuits trop courtes. Quelqu’un pour expliquer autrement, une seule fois. Un manque, on peut le repérer et le nommer, et un manque se comble, lentement parfois, en partie souvent. Une sentence, on ne peut que la subir. Le mot s’accorde, nul ou nulle, et la blessure s’accorde avec lui.

Rodrigue, ce soir-là, ne pense rien de tout cela. Il a douze ans, un quatre plié dans un cahier, et trois lettres en capitales à la place de son nom.

\asterisme

Il serait commode de faire de ce professeur un monstre. Sa journée ressemble plutôt à celle de beaucoup d’autres : des classes où l’on tient à trois par banc, une pile de copies à rendre avant la fin de l’heure, une voix qui fatigue dès le milieu de l’après-midi. Il a dit « toujours nul » comme on dit « au suivant », sans y penser. Le soir même, il l’aura oublié. Il ne se souviendra peut-être jamais de l’avoir dit.

Celui qui dit et celui qui reçoit ne vivent pas la même phrase. Pour l’un, c’est un mot parmi des centaines dans une journée chargée ; pour l’autre, c’est le mot de la journée, et parfois celui de l’année. Je laisse de côté l’excuse comme le procès. Je constate une différence de poids entre les deux bouts de la phrase, et elle explique qu’un mot oublié le soir même par celui qui l’a lancé puisse accompagner l’autre pendant des années. Bien des enseignants le savent mieux que moi et trouvent, dans des conditions qui ne les y aident pas, des mots plus justes. Ce livre leur rendra ce qu’il leur doit.

Toi, en revanche, je te dois quelque chose dès maintenant, puisque tu as ouvert un livre dont le titre ressemble à une consolation. Prends-le plutôt pour une hypothèse : il faudra la vérifier. Ce livre ne te promet pas la réussite. Personne ne peut te la promettre honnêtement, et ceux qui le font vendent autre chose. Il promet de ne pas te mentir. Il te dira ce qui marche, avec ses limites ; ce qui ne dépend pas de toi, et qui pèse lourd ; ce qui en dépend un peu, et ce peu compte quand on part d’un quatre. Il ne te demandera pas de croire en toi. Il te demandera des choses plus modestes et plus fatigantes, comme refaire un exercice sans regarder le corrigé, ou aller dormir à l’heure.

Il te demandera aussi de garder une phrase sous la main, parce qu’elle reviendra : \emph{ce que tu viens de vivre ne prouve pas ce que tu crois.} La note basse a eu lieu. Le rang aussi, et le mot lancé devant la classe, et le rire. Tout cela compte, et je ne te dirai jamais le contraire. Mais ce que ces choses prouvent est beaucoup plus étroit que ce qu’elles t’ont fait croire.

\soustitre{Ce soir, essaie ceci}

Prends un cahier de brouillon neuf. Si tu n’en as pas, prends les dernières pages d’un vieux cahier, celles qu’on garde pour plus tard et qui restent blanches jusqu’à la fin de l’année. Ce cahier-là sera le compagnon de ce livre. Tu y écriras un peu à la fin de chaque chapitre, toujours à la main.

Sur la première page, en haut, écris la date d’aujourd’hui. En dessous, écris une seule phrase : une phrase que quelqu’un a dite sur toi et que tu entends encore. Pas forcément la plus dure. Celle qui revient, qui arrive toute seule quand tu reçois une copie ou qu’on prononce ton nom en classe. Écris-la mot pour mot, avec les mots de celui ou de celle qui l’a dite, même mal choisis, même dans une autre langue que le français. Tu n’es pas obligé d’écrire de qui elle vient.

Ensuite, ne la commente pas. N’ajoute ni « c’est faux » ni « il avait raison ». Ne la souligne pas, et ne mets pas de point d’interrogation au bout. Elle est sur la page au lieu d’être dans ta tête, et ce déplacement-là suffit pour ce soir.

Ferme le cahier. Range-le avec tes affaires de classe, pas au fond d’un tiroir où l’on oublie les choses exprès.

On reviendra à cette page. Pas tout de suite : à la toute fin du livre, quand tu auras de quoi la relire autrement. D’ici là, elle peut attendre. Elle a l’habitude.

\asterisme

Dans la cuisine des Essomba, l’ampoule pend au bout de son fil et attire les moucherons de la saison sèche. Dans la pièce d’à côté, sa mère écoute la radio, le son bas. Rodrigue pose le cahier de brouillon sur la table, couverture vers le haut. Il regarde encore un moment l’étiquette, le trait droit sur son nom, les trois lettres au-dessus. Puis il retourne le cahier, l’ouvre par la fin, à la dernière page, et commence ses exercices là, comme si de rien n’était. Moins trois plus cinq. Moins sept moins deux. Il se trompe au deuxième, efface du doigt mouillé, recommence. La copie pliée glisse sur la table ; il la remet entre deux pages sans la déplier. Il ne finit pas tout.

Quand il referme le cahier, l’étiquette se retrouve en haut, sous l’ampoule.

L’étiquette est restée. Pour l’instant, c’est tout ce qu’on peut dire d’elle.


%% 02-le-rang.md
\partie{Le verdict}{I}

\chapitre{1}{Le rang}

Monsieur Hilaire Kamga a poussé les assiettes au bout de la table de la salle à manger, étalé de vieux journaux pour ne pas tacher le bois, et installé ce qu’il appelle son poste : la calculatrice aux touches usées, le registre de notes de la Seconde C, deux piles de copies de la dernière séquence, une boîte de stylos rouges déjà entamée. Au-dessus, la lampe de bureau. Elle s’allume à condition qu’on cale le fil contre le pied de la table, dans un angle que lui seul connaît. Si l’on bouge la chaise trop vite, elle s’éteint. Il bouge donc lentement. Dehors, Bafoussam est entrée dans décembre : il fait frais dès la tombée du jour, et l’harmattan a posé sur la ville une poussière fine qui voile les lampadaires et se dépose sur tout, jusque sur les copies rangées au fond d’un sac.

Soixante-quatorze noms. Il les connaît tous, ou presque ; il en reste deux ou trois, au fond de la classe, auxquels il associe encore une voix plutôt qu’un visage. Pour chacun, il reprend les notes des deux séquences du trimestre, les coefficients, les moyennes que les collègues des autres matières lui ont fait passer sur des feuilles volantes, et il calcule. La calculatrice fait un petit bruit sec à chaque touche. Dehors, un taxi klaxonne deux fois pour quelqu’un qui ne sort pas. Kamga est le professeur principal de cette seconde du lycée bilingue, et cette soirée-là lui revient chaque trimestre. Quand il réfléchit, il remonte ses lunettes sur son front ; ce soir, elles y restent longtemps.

Vers 22 h, il range les moyennes dans l’ordre pour établir les rangs. La calculatrice ne fait pas ce travail à sa place : il le fait au crayon, sur une feuille de brouillon, en remontant et en redescendant la colonne du doigt. Au milieu, il s’arrête. La quarante-sixième et le quarante-septième sont séparés de deux centièmes. Deux centièmes, c’est un demi-point perdu sur une interrogation de géographie en octobre, une virgule oubliée dans une dictée, une absence un jour de devoir rattrapée tant bien que mal, un exercice de physique rendu avec de la fièvre. Sur le bulletin, l’une sera quarante-sixième et l’autre quarante-septième, et rien n’indiquera que ces deux places se touchent presque. Il le sait. Il écrit les chiffres quand même, parce qu’il faut bien les écrire.

Plus haut dans la colonne, il bute sur Divine Nfor. Seize en mathématiques. Ce seize, c’est lui qui l’a mis, et sans hésiter : deux devoirs propres, une démonstration qu’il a relue deux fois pour être sûr de ne pas se tromper en sa faveur. Mais Divine est arrivée de Bamenda il y a trois ans, elle pense encore en anglais, et ses copies de français reviennent couvertes de rouge. Le calcul fait son travail sans état d’âme. Seize ici, sept là, des notes moyennes ailleurs, des coefficients qui tirent dans un sens puis dans l’autre. Résultat : trente et unième. Kamga regarde ce nombre un long moment. Le rang ne voit pas le seize. Il ne voit pas davantage le sept, d’ailleurs ; il voit une place dans une file, entre un trentième et un trente-deuxième qui n’ont presque rien en commun avec elle.

Alors il fait ce qu’il fait chaque trimestre depuis des années, et que personne ne lui a jamais demandé. Il ouvre un petit cahier à couverture verte, corné aux angles, qu’il appelle, quand il en parle, ce qui est rare, son cahier des noms. Une page par classe, une ligne par élève. Pour chacun, il écrit une chose réussie pendant le trimestre, un fait plutôt qu’une appréciation. Un exercice, une question posée au bon moment, une erreur qu’un élève a trouvée tout seul dans sa propre copie, un tableau effacé sans qu’on le lui demande. Pour certains, il cherche longtemps, en feuilletant le registre et sa mémoire. Il finit toujours par trouver, et il se méfie un peu de lui-même les fois où il trouve trop vite.

\asterisme

Une note est une mesure, et tout ce qui suit part de là. Le jour de la séquence, sur ces trois exercices-là, dans cette salle où il faisait chaud, avec le temps qu’on t’a donné et la forme que tu avais ce matin-là, tu as produit une copie, et quelqu’un l’a lue avec une grille en tête. Ce quelqu’un a mis un nombre. Ce nombre dit une chose vraie : ce jour-là, sur cette tâche, voilà ce qui est passé de ta tête au papier. Il ne dit pas pourquoi le reste n’est pas passé, ni ce que tu sauras dans un mois. Et il ignore tout de ce que tu sais faire et qu’on ne t’a pas demandé.

Au marché, la balance de la vendeuse pèse les tomates ; elle ne dit pas si elles sont bonnes. Le thermomètre de l’infirmerie dit que tu as trente-huit et demi ; il ne dit pas ce que tu as. Personne ne jette la balance ni le thermomètre pour autant. On s’en sert, en sachant ce qu’ils mesurent et ce qu’ils laissent de côté. La note mérite le même traitement : qu’on la prenne au sérieux, et qu’on soit précis sur ce qu’elle mesure.

On le dit rarement aux élèves, sans doute de peur qu’ils en tirent la mauvaise conclusion. Dans les années 1930, en France, le psychologue Henri Piéron et ses collaborateurs ont fait une expérience d’une simplicité presque insolente : ils ont confié les mêmes copies d’examen à plusieurs correcteurs expérimentés, sérieux, de bonne foi. Une même copie pouvait recevoir des notes très éloignées les unes des autres. Piéron avait déjà donné un nom à ce champ d’étude, la docimologie, qui désigne l’étude scientifique des examens et de la notation. Depuis, de nombreuses études ont confirmé l’idée générale : une note dépend aussi de celui qui corrige, de la place de la copie dans la pile, et même de l’écriture de l’élève.

Kamga connaît cela sans avoir besoin de lire les études. Il sait que sa soixante-dixième copie, à minuit, est lue par un homme plus fatigué que sa première. Il a pris l’habitude de mélanger les piles, de corriger exercice par exercice plutôt que copie par copie quand il en a le courage, et de relire au matin ce qu’il a noté trop tard. Le problème demeure. Il diminue un peu, et Kamga préfère un peu à rien.

Inutile, donc, d’aller expliquer à ton professeur que ses notes ne valent rien : ce serait injuste, elles valent quelque chose, et la plupart des correcteurs font de leur mieux avec le temps dont ils disposent. Je te raconte Piéron pour une raison plus étroite. Si une même copie peut recevoir des notes différentes selon la main qui la corrige, alors une note seule ne peut pas porter tout le poids qu’on met parfois dessus. Elle peut dire qu’un exercice n’est pas encore maîtrisé ; pour dire « tu es comme ça », elle n’a pas les épaules. Ce que la docimologie rend imprudent, ce n’est pas la note, c’est le verdict.

\asterisme

Le rang est encore autre chose. La note, au moins, parle de ta copie. Le rang parle de ta copie comparée à toutes les autres, et il dépend autant de ta classe que de toi.

Fais l’expérience en pensée. Prends Divine, avec exactement ses notes, et pose-la dans une autre seconde, dans un autre lycée, où les élèves seraient plus faibles en mathématiques ou plus forts en français. Elle n’a pas changé d’une virgule. Son rang, si. Elle pourrait être douzième ou cinquantième sans avoir rien appris ni rien oublié. Un rang ressemble à une place dans une file d’attente : il dit qui est devant toi et qui est derrière, sans rien dire de la distance qui te sépare du guichet.

La moyenne fait une chose plus discrète, et plus trompeuse : elle écrase. Seize en mathématiques et sept en français donnent, à coefficients égaux, onze et demi. Or onze et demi ne ressemble à rien de ce que Divine a fait. Elle n’a jamais rendu une copie à onze et demi. Elle a rendu des copies très bonnes dans une matière et difficiles dans une autre, et la moyenne en a fait une élève « passable » partout, ce qui est le seul portrait sûrement faux. Le seize disait : il y a là quelqu’un qui raisonne. Le sept disait : il manque à cette élève des outils précis en français écrit, qu’on peut nommer un par un. Ces deux informations servent. Leur moyenne ne sert presque à rien.

À la table de Bafoussam, vers 23 h, Kamga a besoin de la gomme, qui a roulé de l’autre côté du registre. Il tend le bras, la chaise recule d’un centimètre, et la lampe s’éteint. Il soupire, se penche sous la table dans le noir, retrouve le fil à tâtons et le recale dans son angle. La lumière revient en clignant. Il en profite pour se lever, boire un verre d’eau dans la cuisine, regarder par la fenêtre la poussière qui tourne sous le lampadaire de la rue. Ces stylos rouges, il les achète par boîtes au début de l’année, au même libraire, et il sait à peu près combien de copies tient un stylo. Il se dit parfois qu’il pourrait mesurer son année en boîtes plutôt qu’en trimestres. Puis il revient s’asseoir, lentement, et reprend la colonne là où il l’avait laissée.

La moyenne écrase aussi le temps. Dans la colonne de Kamga, un garçon assis près de la fenêtre avait quatre en mathématiques à la première séquence et huit à la deuxième. Huit reste une mauvaise note ; personne ne le niera, et lui moins que les autres. Mais entre quatre et huit, il a doublé. Sur son bulletin, il aura six de moyenne en mathématiques et un rang dans les soixante, et rien n’indiquera dans quel sens va la courbe. Kamga, en recopiant le six, a hésité, puis il a ajouté au crayon, dans la marge du registre, une petite flèche qui monte. Elle ne figurera nulle part ailleurs.

Un trimestre se termine, et tu reçois une feuille où tout cela a été fondu en deux nombres : une moyenne générale, un rang. Si tu la lis comme un jugement, tu n’y trouves qu’une phrase, toujours la même, et elle parle de toi. Si tu la lis comme une carte, elle te montre des reliefs : ici c’est haut, là c’est bas, cette matière monte depuis octobre, celle-là descend. Un jugement se subit. Une carte sert à décider où aller la semaine prochaine. Ton rang de décembre ne prouve pas ce que tu crois. Il dit où tu te trouvais dans une file, à une date donnée, parmi des gens donnés ; le reste, ce qui t’est utile, est déjà sur la feuille, à condition de la lire ligne par ligne au lieu de chercher seulement le chiffre en bas.

\asterisme

Le regard, lui, entre dans la note avant même qu’on l’écrive. Au milieu des années 1960, un psychologue américain, Robert Rosenthal, et une directrice d’école, Lenore Jacobson, ont fait passer un test aux élèves d’une école primaire, puis ont annoncé aux enseignants que certains enfants allaient faire de grands progrès dans l’année. Ces enfants avaient en réalité été tirés au hasard. À la fin de l’année, une partie d’entre eux, surtout parmi les plus jeunes, avaient davantage progressé que leurs camarades. Rosenthal et Jacobson ont raconté cette étude dans un livre devenu célèbre, publié en 1968 et traduit en français sous le titre \emph{Pygmalion à l’école}.

Elle a été très discutée depuis, et pour de bonnes raisons : la méthode a été critiquée, et les reprises n’ont pas toutes donné le même résultat. Les travaux qui ont suivi suggèrent une conclusion plus modeste que la légende : les attentes des enseignants ont un effet réel sur les élèves, mais en général limité. J’en retiens une leçon qui va dans les deux sens. Le regard d’un adulte pèse un peu sur ce que devient un élève, et c’est une raison pour que les adultes fassent attention à ce qu’ils regardent. Mais ce poids reste modeste, et c’est une raison pour que tu ne remettes pas ton année entière entre les mains de quelqu’un qui t’aurait mal regardé. Il t’a peut-être classé trop vite. Ce que tu feras, il ne l’a pas décidé.

Kamga a entendu parler de ces travaux pendant sa formation, il y a longtemps, et il ne s’en souvient qu’à moitié. Il se souvient mieux d’une observation faite sur lui-même : quand il a décidé qu’un élève était faible, il le voit faible plus facilement. Il prend sa bonne réponse pour un hasard et sa mauvaise pour une confirmation. Le cahier vert est né de là, ou à peu près ; il ne saurait pas dire quand. C’est un garde-fou qu’il a fabriqué contre lui-même. Une ligne par élève l’oblige, une fois par trimestre, à chercher ce qu’il ne cherchait pas.

Il ne faut pas en faire un saint pour autant. Kamga est fatigué. Il enseigne depuis plus de vingt ans, ses classes dépassent soixante-dix élèves, et le samedi matin il donne des répétitions dans le salon d’une famille du quartier pour que le mois se termine à la fin du mois. Il se trompe parfois dans une moyenne et doit corriger un bulletin à la main, en grommelant contre la calculatrice, qui n’y est pour rien. Il lui arrive de s’agacer en classe contre un élève qui n’a rien fait de plus que les autres. Il rit rarement, mais quand il rit, c’est franchement, et la classe entière se retourne pour voir ce qui a bien pu arriver. Il sait que les notes mesurent mal, et il note quand même, parce qu’il faut bien que l’année se décide sur quelque chose. Comme beaucoup d’enseignants, il garde la lucidité et fait le travail, les deux à la fois. Il tient.

Les rangs de décembre ne sont d’ailleurs qu’un début. À la fin de l’année viendra une autre feuille, la liste des admis, affichée au portail du lycée, devant laquelle des élèves se presseront et liront de haut en bas, puis de bas en haut. Le rang du premier trimestre en est la répétition générale : il apprend à des adolescents de quinze ans à chercher leur nom dans une colonne et à en conclure une vérité sur eux-mêmes. Autant apprendre tôt à lire ces colonnes correctement.

\soustitre{Ce soir, essaie ceci}

Reprends ton dernier bulletin. Si tu ne l’as plus, ou si le trimestre n’est pas terminé, prends tes dernières notes, celles qu’on t’a rendues ou celles que tu as recopiées quelque part. Ouvre ton cahier de brouillon à une page neuve et écris la date en haut.

À côté de chaque moyenne, écris combien de devoirs elle résume : deux, parfois un seul. Puis écris de quel type ils étaient : un exercice d’application, une rédaction, un oral, une interrogation surprise un lundi matin, un devoir sur table de deux heures. Tu vas sans doute découvrir que certaines moyennes qui te pèsent depuis des semaines reposent sur une seule copie, ou sur deux copies qui ne se ressemblent pas du tout. Rien n’est excusé pour autant, mais tout devient plus net, et cette netteté te dira par où reprendre.

Ensuite, choisis la matière où ta note est la plus basse. Cache cette note avec ta main. Essaie de retrouver ce que tu pensais en sortant du devoir, le sujet encore dans la tête, et écris la note que tu croyais avoir. Sois honnête : personne ne lira cette page. Puis écris la vraie à côté.

Regarde l’écart. Si tu croyais avoir douze et que tu as eu six, l’information est de taille : tu pensais savoir, et tu ne savais pas encore, ou pas sous la forme qu’on te demandait. Le travail à faire porte alors autant sur ta façon de vérifier ce que tu sais que sur la leçon elle-même. Si tu croyais avoir six et que tu as eu six, la nouvelle est presque bonne : tu sais où tu en es, et tu possèdes déjà un instrument précieux, un œil juste sur ton propre travail. Si tu croyais avoir quatre et que tu as eu neuf, tu te sous-estimes, et cela aussi se corrige.

Cet écart t’en apprend plus que ton rang, qui te situe parmi les autres, alors que l’écart te dit si tu sais toi-même ce que tu sais. Demain, au premier cours de cette matière, garde-le en tête. Dans trois jours, quand tu referas un exercice, devine ta note avant de vérifier tes réponses, et note les deux. Au bout de quelques fois, regarde si l’écart se resserre.

Tu n’as besoin de personne pour faire cela. Il te faut un stylo et un quart d’heure d’honnêteté, qui est la partie la plus coûteuse.

\asterisme

Il est presque minuit quand Kamga referme le registre. Les rangs sont recopiés au propre sur les fiches que le secrétariat imprimera demain matin, soixante-quatorze places du premier au dernier, et parmi elles deux élèves que deux centièmes séparent et qui ne le sauront jamais. Il range les stylos rouges dans leur boîte, glisse les copies dans le sac, s’étire, et la chaise grince. Il tend la main vers l’interrupteur, mais avant qu’il l’atteigne le fil se décroche tout seul du pied de la table, et la lumière s’en va. Il reste un instant assis dans le noir. Puis il rit tout seul, un rire bref, pour personne.

Les bulletins partiront demain avec leurs rangs imprimés. Dans le cahier vert resté ouvert sur la table, à la ligne de Divine Nfor, il a écrit de sa petite écriture penchée : « Démonstration du 14 novembre : propre, juste, et elle a trouvé une autre méthode que la mienne. » Ce cahier, aucun bulletin ne le reproduit.

\begin{pourcahier}{Pour le cahier}
Écris deux phrases : « Ma note de \_\_\_ dit que… » et « Elle ne dit pas que… ». Fais-le pour une seule matière, mais fais-le honnêtement : la première phrase compte autant que la seconde.
\end{pourcahier}


%% 03-les-marches-manquantes.md
\chapitre{2}{Les marches manquantes}

À Ebolowa, en novembre, la nuit tombe d’un seul coup. Les motos se font plus rares au carrefour, les radios montent d’un cran, et derrière la maison quelqu’un pile encore dans un mortier, avec la régularité d’une horloge. Chez les Essomba, on travaille à la table de la cuisine, sous l’ampoule qui pend au bout de son fil. La nappe en plastique colle un peu aux avant-bras. Bella, la petite cousine, dort déjà ; son cartable est resté appuyé contre un pied de la table, la fermeture à moitié ouverte.

Rodrigue a quatorze ans, il est en troisième, et il est devant l’exercice 4 depuis vingt minutes. L’énoncé tient sur une ligne : \emph{Réduire l’expression A = 2/3 x + 1/4 x − 5.} Il a fait ce qu’il fallait faire au début. Il a mis le x en facteur, comme le professeur l’a montré au tableau le matin même, et il a obtenu une ligne propre : A = (2/3 + 1/4) x − 5. Puis tout s’est arrêté. Au milieu de la ligne, entre les parenthèses, une petite addition le regarde : 2/3 + 1/4. Il a écrit 3/7 au crayon, puis il l’a gommé. Il sait que 3/7 est faux. Il se souvient d’une croix rouge, en sixième ou en cinquième, posée sur ce genre de réponse. Il ne se souvient pas de la raison.

Il sort le petit tournevis de la poche de sa chemise et le fait tourner entre ses doigts. C’est avec lui que, le samedi, il ouvre des téléphones que d’autres ont déclarés morts. Ici, il n’ouvre rien. Il sait qu’il devrait savoir, et le tournevis tourne plus vite.

En face de lui, à la même table, sa mère fait les comptes du mois. Mama Odile est trésorière de sa tontine, et ce soir est le soir des comptes : demain, au marché, entre les régimes de plantain et les seaux d’arachides, on viendra lui demander qui a cotisé, qui doit encore, combien revient à qui. Elle a devant elle un cahier d’écolier à spirale où les noms des membres sont écrits en colonne, avec, à côté, de petites croix, des ronds, des traits. Pas une opération. Elle murmure. « Mama Ngono, c’est fini. Tantine Rose a donné la moitié, l’autre moitié samedi. Les deux sœurs Mvondo ont une part pour elles deux. » Elle ne compte pas sur ses doigts et ne lève pas les yeux. Elle ne se trompe pas non plus, et Rodrigue le sait pour une raison simple : à la tontine, on conteste tout, jusqu’au goût du jus servi chez celle qui « bouffe » ce mois-ci, c’est-à-dire qui reçoit la cagnotte, et personne n’a jamais contesté les comptes de sa mère.

Par la fenêtre ouverte, à cause de la chaleur, on voit la maison du voisin. Elle est en chantier depuis si longtemps que plus personne, dans le quartier, ne dit « le chantier » ; on dit « chez le voisin », et tout le monde voit de quoi il s’agit. Sur le côté, un escalier de béton nu monte vers un étage qui n’a pas encore de toit. Il n’a pas de rampe. À mi-hauteur, il lui manque deux marches que personne n’a encore coulées : un vide net, assez large pour qu’on hésite, entre la partie basse et la partie haute. Le jour, les enfants du quartier montent jusqu’au trou, regardent de l’autre côté, et redescendent. Ce soir, l’escalier se réduit à une découpe grise contre le ciel, avec son vide au milieu.

\asterisme

On pourrait dire de Rodrigue qu’il est faible en mathématiques. Ses bulletins le disent, en chiffres, chaque trimestre. Regarde pourtant l’exercice d’un peu plus près. Le travail de troisième, il l’a fait : mettre en facteur, c’était la leçon du jour, et il l’a réussie. Ce qui le bloque date d’avant. L’addition de fractions, on la suppose acquise depuis longtemps, rangée quelque part entre le CM2 et la sixième, et le professeur s’en sert comme on se sert d’une marche, en posant le pied dessus sans la regarder, pour aller plus haut. Ce soir, le pied de Rodrigue cherche la marche et trouve le vide.

Le savoir scolaire est fait ainsi, surtout en mathématiques, en physique et dans les langues : chaque leçon s’appuie sur celles d’avant. Le calcul littéral suppose les fractions, qui supposent la division, qui suppose les tables. Tant que les marches sont là, on monte sans y penser, au point de croire que chaque leçon tient toute seule. Le jour où une marche manque, la suivante devient hors d’atteinte, et l’élève, qui ne voit pas le trou sous son pied, prend pour une incapacité ce qui est une absence. Il se dit : je n’y arrive pas. La phrase juste serait plutôt : il me manque une marche, plus bas.

Je me méfie du mot « difficile », qu’on emploie pour tout. Une leçon difficile demande un effort à quelqu’un qui a sous la main ce qu’il faut pour la faire. Une leçon posée sur une marche absente appelle un autre mot, et surtout un autre remède. Sur ce point, de nombreuses études convergent : parmi tout ce qui permet de prévoir ce qu’un élève saura à la fin d’une leçon, ce qu’il sait déjà du sujet compte parmi les indices les plus solides, même si cela ne dit pas tout de ce qu’il peut encore gagner. Le reste compte aussi, bien sûr. Mais devant un élève qui bute, avant de mesurer sa tête, il faudrait compter les marches qu’il a sous les pieds.

Le professeur de troisième, lui, ne voit pas le trou. Il a devant lui une classe pleine, un programme à finir avant le BEPC, et une leçon qui, sur le papier, commence là où la sixième s’est arrêtée. Il ne peut pas redescendre l’escalier avec chacun ; il fait monter tout le monde à la même heure, par les mêmes marches, en espérant qu’elles soient là. Quand un élève reste en bas, il voit un élève qui ne suit pas. Pour voir le reste, il lui faudrait un temps qu’il n’a presque jamais. Ce manque-là n’accuse personne ; il vaut mieux le connaître, pour ne pas en tirer de mauvaises conclusions sur soi.

Il y a deux ans, en cinquième, on a rendu à Rodrigue une copie avec un quatre et deux mots. Peut-être qu’on y trouvait déjà des fractions mal additionnées. Peut-être pas. Une chose est sûre : ce jour-là, personne n’a cherché quelle marche manquait. Un verdict constate, puis il passe à la copie suivante.

\asterisme

Reviens à la table de la cuisine. Ce que Rodrigue doit tenir dans sa tête, à cet instant, pèse plus lourd qu’il n’y paraît. Le x qu’il a mis en facteur, d’abord, et qu’il ne doit pas perdre en route. Le moins cinq, à droite, qui attend son tour. La consigne, « réduire », dont il ne sait pas très bien ce qu’elle exige. Et au milieu, cette addition qu’il ne sait pas faire sans y penser, et qu’il essaie donc de reconstruire en direct : est-ce qu’on additionne en haut et en bas, est-ce qu’il faut multiplier, et quoi, et par quoi. Pendant qu’il cherche, le x lui glisse des doigts.

Les psychologues appellent mémoire de travail cette part de l’esprit qui retient et manipule, consciemment, ce qu’on est en train de traiter. On ne peut y tenir que quelques éléments à la fois. Elle est étroite chez tout le monde, chez le premier de la classe comme chez le dernier, et les différences de l’un à l’autre sont bien plus petites que celles que fait ce qu’on a déjà automatisé. Le reste du savoir dort ailleurs, dans une mémoire à long terme beaucoup plus vaste, et n’en sort que pour servir. Toute la différence se joue là. Ce qu’on a automatisé ne coûte presque rien : celui qui sait que trois fois quatre font douze ne « pense » pas la multiplication, le résultat arrive tout seul et ne prend pas de place. Celui qui doit la reconstruire, en revanche, l’installe dans l’espace étroit, où elle prend la place de tout le reste.

Certains soirs, Mama Odile rentre du marché avec la bassine des invendus sur la tête et un sac dans chaque main. Arrivée devant la porte, elle n’a plus de main pour l’ouvrir. Il faut que quelqu’un vienne, ou qu’elle pose un sac par terre. Un élève à qui il manque une base se trouve dans la même situation : la base non acquise lui occupe une main entière, et il ne lui en reste aucune pour la leçon du jour. C’est pour cela que tant d’élèves disent « je ne comprends rien » à des leçons qu’ils comprendraient sans peine s’ils avaient les mains libres. À ce moment-là, ils manquent de place.

À la table, Rodrigue a essayé une deuxième fois. Il a multiplié en croix, parce qu’on lui avait appris un jour à multiplier en croix, et il a obtenu des nombres qui ne ressemblent à rien. Il a regardé l’heure sur le téléphone de sa mère, posé à côté du cahier à spirale. Puis il a fait ce que font presque tous les élèves dans sa situation : il a cessé de penser à la fraction et s’est mis à penser à lui. À sa tête, qui ne marche pas. Au professeur de cinquième. Au devoir de vendredi. L’exercice 4 a disparu derrière une question plus ancienne, qui n’a rien de mathématique.

Les marches manquantes, d’ailleurs, ne restent pas sagement à leur place. Le psychologue Keith Stanovich a décrit en 1986, à propos de la lecture, ce qu’il a appelé, en reprenant une expression du sociologue Robert Merton, un « effet Matthieu », d’après le verset de l’Évangile selon Matthieu où il est dit, en substance, qu’on donnera encore à celui qui a déjà. L’enfant qui lit bien au début lit davantage, parce que c’est agréable ; en lisant davantage, il apprend des mots, et ces mots l’aident à mieux lire. L’enfant qui lit mal évite de lire, parce que c’est pénible, et l’écart tend à se creuser d’année en année. Les recherches qui ont suivi n’ont pas toujours retrouvé ce creusement avec la même netteté : c’est une tendance documentée, qui souffre des exceptions. Rodrigue, lui, la connaît de l’intérieur sans en connaître le nom. Chaque leçon qu’il n’a pas pu suivre, faute de fractions, est devenue à son tour une marche absente pour la leçon d’après. Le trou de sixième a fait des petits.

\asterisme

Chaque matière a son escalier, et chaque élève a le sien. La dissertation de français repose sur la conjugaison et les accords ; quand ces marches-là manquent, on a des idées et on ne sait pas les poser sur la page. En physique, la marche cachée s’appelle souvent proportionnalité : celui qui ne la maîtrise pas croit détester les forces et les circuits, alors que la règle de trois lui échappe depuis le collège. L’anglais demande d’abord du vocabulaire, ces centaines de mots simples sans lesquels aucune règle de grammaire ne sert à grand-chose. Et l’histoire suppose la lecture d’un document, qu’on croit savoir faire parce qu’on sait lire. Dans chaque cas, l’élève situe sa difficulté à l’étage où elle se montre. Elle est presque toujours née plus bas.

Comment manque-t-on une marche ? De beaucoup de façons, et la plupart sont banales. Une maladie au mauvais moment de l’année. Un déménagement entre deux écoles qui n’en étaient pas au même chapitre. Un professeur muté en cours de trimestre et remplacé tard. Une leçon faite le jour où tu étais absent, et que personne n’a refaite pour toi. Une année entière passée à recopier le tableau sans comprendre, en se promettant d’y revenir, sans jamais y revenir. Personne n’est coupable de tout cela, et tout le monde en porte un peu. De toute façon, chercher un coupable n’a jamais coulé une seule marche.

La honte, je crois, bloque plus que tout le reste. Réapprendre en troisième ce qu’on aurait dû savoir en CM2 a quelque chose d’humiliant, ou du moins on le vit ainsi. On a quatorze ans, une voix qui a mué, un tournevis dans la poche de chemise ; on ne va pas ouvrir un livre d’enfant. Alors on fait semblant. On recopie la ligne du voisin, on attend que le professeur passe à autre chose, on rit le premier quand il faut rire. Et la marche reste vide, mieux gardée par la honte que par n’importe quel cadenas.

Ce soir, le cartable de Bella est à quelques centimètres du genou de Rodrigue. Dedans, il le sait, se trouve le livre de mathématiques de CM2, avec ses disques coupés en quatre et ses tablettes de chocolat partagées entre trois enfants. Il pourrait se pencher. Il ne se penche pas. Qu’une cousine de dix ans découvre, demain matin, que son grand cousin de troisième a fouillé dans son livre pour comprendre une addition : il y a des risques qu’on ne prend pas.

\asterisme

Pendant ce temps, Mama Odile a fini une colonne et en commence une autre. Rodrigue l’écoute sans en avoir l’air. « Les Mvondo, ce mois, c’est leur tour de bouffer. Chacune prend la moitié. Mais la petite doit encore le quart de la dernière fois, alors on retire ça de sa moitié. » Elle s’arrête, compte en silence, pince les lèvres, coche. Rodrigue ne suit pas le détail. Il entend des moitiés et des quarts, et quelque part un tiers : celui d’une femme entrée dans la tontine en cours d’année, qui ne paie qu’une partie de la part. Il entend surtout qu’aucun de ces morceaux ne gêne sa mère.

Il serait facile d’en tirer une scène édifiante, la vendeuse de vivres qui en remontrerait à tous les professeurs. La table de la cuisine ne s’y prête pas. Mama Odile ne « connaît pas les fractions ». Elle connaît les parts. Elle sait ce qu’est une moitié de cotisation, ce que pèse le quart d’un retard, comment répartir une somme entre deux sœurs sans qu’aucune se sente volée. C’est un savoir mathématique réel, juste, éprouvé chaque mois devant des femmes qui ne laissent rien passer.

Ce savoir, elle l’a construit au marché, sans manuel, à force de vendre des arachides au tas et au demi-tas, de rendre la monnaie à des clientes pressées, de partager un sac de macabo avec une voisine de table quand aucune des deux n’avait de quoi l’acheter entier. Il est solide parce qu’il a servi tous les jours, des années durant, sous le regard de gens qui comptaient aussi. Il est limité parce qu’il n’a jamais eu besoin de s’écrire. Il vit sous d’autres noms que ceux de l’école, dans d’autres gestes, et il ne se traduit pas tout seul. Si Rodrigue lui tendait son cahier, elle regarderait les barres et les petits chiffres posés les uns sur les autres, et elle dirait, sans honte particulière : « Ça, c’est votre affaire de l’école. Moi, je ne saurais pas te faire ça. » Elle aurait raison. Elle ne saurait pas.

Entre ce qu’elle sait et ce que son fils cherche s’ouvre donc un écart, une traduction que personne n’a faite, deux ou trois marches que personne n’a coulées, et rien qui ressemble à un mur. Un manque se comble, lentement, par le bas, avec du temps et un peu d’aide, mais il se comble. Le mot écrit en capitales sur une étiquette, il y a deux ans, n’en savait rien.

\soustitre{Ce soir, essaie ceci}

Prends un exercice que tu as raté récemment, dans la matière que tu veux, et ouvre ton cahier de brouillon à une page neuve. Ne refais pas l’exercice. Cherche seulement la première étape où tu t’es arrêté, ou celle où ton professeur a mis sa première croix, et recopie-la en haut de la page. Dessous, écris une phrase qui commence ainsi : « Pour faire ça, il fallait savoir… », et termine-la honnêtement. Pour Rodrigue, ce serait : « Pour réduire A, il fallait savoir additionner 2/3 et 1/4. »

Puis prends ce que tu viens d’écrire et recommence, une ligne plus bas : « Et pour savoir ça, il fallait savoir… ». Rodrigue écrirait sans doute qu’il fallait savoir mettre deux fractions au même dénominateur. Et, plus bas encore, qu’il fallait savoir ce qu’est un dénominateur, et trouver un nombre qui se trouve à la fois dans la table de trois et dans celle de quatre. Descends ainsi, ligne après ligne, sans te presser et sans tricher, jusqu’à tomber sur une chose dont tu es sûr. Pas à peu près sûr : sûr au point de pouvoir l’expliquer à quelqu’un. Pour Rodrigue, ce seraient sans doute les tables de multiplication, qu’il connaît bien, parce qu’on les lui a fait réciter debout pendant des années.

La marche juste au-dessus de cette chose sûre est la tienne. C’est là que ton escalier s’interrompt. Entoure-la. Et maintenant, le plus dur : ne la répare pas ce soir. Ne file pas chercher un cours, ne te lance pas dans vingt exercices. Note-la, avec la date, et ferme le cahier. Tu as le droit de t’arrêter là, et même le devoir : on répare mal ce qu’on n’a pas fini de regarder.

Le diagnostic est déjà un travail. Tu viens de remplacer une phrase vague, « je suis mauvais en maths », par une phrase précise, du genre « je ne sais pas mettre deux fractions au même dénominateur ». La première te laisse les bras ballants. La seconde te dit où poser le pied.

Si tu descends sans trouver de fond, si chaque ligne en appelle une autre et que tu as l’impression de tomber, arrête-toi après cinq ou six lignes. Cela arrive. Il te faudra simplement quelqu’un pour descendre avec toi. Garde la page, elle servira.

À la table, l’exercice 4 n’a pas bougé. La ligne s’arrête toujours au même endroit, sur la même petite addition, et Rodrigue a cessé de faire semblant de chercher. Il regarde sa mère. Elle vient de refermer le cahier à spirale et le glisse sous la nappe en plastique, là où elle range les papiers importants. Les comptes sont justes ; demain, au marché, personne ne contestera.

Il pose la question sans l’avoir préparée, presque malgré lui. « Mama, tu fais comment pour compter tout ça sans rien écrire ? » Elle le regarde comme on regarde quelqu’un qui demande comment on fait pour marcher. Elle hausse les épaules. « Je partage. » Puis elle se lève, vérifie que la porte de derrière est bien fermée, et lui dit de ne pas se coucher trop tard, que l’ampoule ne se paie pas toute seule.

Rodrigue reste devant sa page. Tout en bas, sous la ligne qui bloque, il écrit au crayon un seul mot : « partager ». Il ne sait pas pourquoi. Il ne le gomme pas. Il remet le tournevis dans la poche de sa chemise et éteint l’ampoule.

Dehors, le mortier s’est tu. Dans le noir, le chantier du voisin garde ses deux marches manquantes.

\begin{pourcahier}{Pour le cahier}
Dessine un escalier de cinq marches pour la matière qui te fait le plus peur. Sur chaque marche, écris une notion, de la plus ancienne en bas à celle d’aujourd’hui en haut. Entoure la marche où tu n’es plus sûr.
\end{pourcahier}


%% 04-la-main.md
\chapitre{3}{La main}

À 6 h 15, le pont sur le Wouri respire encore. Grâce Ewane connaît ses embouteillages comme d’autres connaissent leurs leçons : celui qui se forme à la sortie de Bonabéri vers 6 h 45, celui du vendredi soir qui ne se défait jamais tout à fait, celui des soirs de match, et celui des jours de pluie, qui n’obéit à personne. Alors elle part tôt. Du taxi, quand le fleuve est gris, on ne sait plus très bien où il s’arrête et où le ciel commence. Elle a seize ans, elle est en Première A4 espagnol dans un lycée de l’autre rive, et ce mardi d’octobre commence par deux heures de mathématiques.

La salle est pleine, comme toujours. On s’assoit à trois par banc, les sacs coincés entre les pieds, et l’air entre comme il peut par les lames de verre des fenêtres. Le professeur de mathématiques est un homme pressé, et il a ses raisons : la deuxième séquence approche, le programme a pris du retard, et il ne voit cette classe que trois fois par semaine. Il écrit vite, il parle en écrivant, la craie dans la main droite et le chiffon dans la gauche, si bien qu’une ligne disparaît parfois avant qu’on ait fini de la recopier. Il n’a rien contre personne. Il est en retard sur octobre.

Vers 8 h, la pluie commence. À Douala, en cette saison, elle ne s’annonce pas. Quelques gouttes sur la tôle, isolées, presque polies, puis d’un coup le toit entier qui gronde, comme si quelqu’un y vidait des sacs de gravier. Au tableau, le professeur a écrit 3(x − 2) + x = 10. Il se tourne à demi vers la classe, dit une phrase en passant à la ligne suivante, et c’est à ce moment-là, pile, que la pluie couvre tout. Grâce voit sa bouche bouger. Elle voit la craie descendre et écrire 3x − 6 + x = 10. Les parenthèses ont disparu, et le moins est passé au 6. Elle n’a pas entendu pourquoi.

Sa main droite se lève. Pas au-dessus du banc : en dessous, à hauteur du genou, comme une main qui demanderait la permission de demander. Elle monte de quelques centimètres, hésite, puis poursuit son mouvement et va remettre derrière l’oreille une mèche qui s’échappe toujours du même côté de ses tresses, et qui n’avait rien demandé.

Le calcul a duré moins d’une seconde. Il était d’une grande précision. Si elle lève la main, le professeur s’arrête ; s’il s’arrête, il se retourne ; s’il se retourne, quarante têtes se tournent avec lui et découvrent que Grâce Ewane n’a pas compris une ligne que tout le monde, sans doute, a comprise. Il faudra parler fort, à cause de la pluie. Il faudra peut-être répéter. Alors elle copie la ligne. Elle la copie avec soin, le moins bien collé à son 6, comme on recopie une phrase dans une langue qu’on ne parle pas. Et elle sait ce que cela fait, parce que l’espagnol, elle le parle : elle connaît la différence entre une phrase qu’on comprend et une phrase qu’on dessine.

Sur le cahier, la ligne est juste. Dans sa tête, il manque un pas, de la taille d’un signe.

\asterisme

Cette main, on la juge mal. On dit d’une élève comme Grâce qu’elle est timide, et on le dit comme on parlerait d’une mauvaise vue : un défaut de fabrication, qu’il faudrait corriger. Or le dimanche, Grâce chante dans la chorale de sa paroisse. Elle tient la partie d’alto devant une église pleine, et sa voix ne tremble pas. Ce qui lui fait peur, ce n’est donc ni la foule ni le son de sa propre voix ; c’est d’être regardée au moment précis où elle ne sait pas. Une peur pareille dépend de la situation bien plus que du caractère, et une situation, cela s’examine.

Le calcul de Grâce repose sur une hypothèse qu’elle n’a jamais vérifiée : quarante têtes se tournent, quarante paires d’yeux voient, quarante mémoires retiennent. Des psychologues se sont penchés sur ce genre d’hypothèse. En 2000, Thomas Gilovich et ses collègues ont demandé à des étudiants d’entrer, avec un peu de retard, dans une salle où d’autres étaient déjà assis, vêtus d’un tee-shirt qu’ils trouvaient gênant, à l’effigie d’un chanteur de variétés que personne de leur âge n’aurait porté de son plein gré. Après quelques instants, on les faisait ressortir et on leur demandait combien de personnes avaient remarqué le tee-shirt. Ils en surestimaient nettement le nombre : bien moins de gens l’avaient vu qu’ils ne le croyaient. Les chercheurs ont appelé cela l’effet projecteur. Nous nous croyons pris dans une lumière que les autres, occupés d’eux-mêmes, braquent rarement sur nous. L’expérience a été beaucoup citée et discutée ; l’idée générale, que nous exagérons l’attention qu’on nous porte, a été retrouvée dans d’autres situations.

Dans la salle, sous la pluie, que se serait-il passé si Grâce avait levé la main ? Quelques têtes se seraient tournées, oui. La plupart seraient restées sur le tableau, parce qu’une question offre d’abord une pause, et qu’une pause, dans un cours qui va trop vite, se met à profit : on finit de recopier, on rattrape la ligne que le chiffon menaçait. Une ou deux personnes auraient pensé \emph{moi non plus, je n’ai pas entendu}. Un rire serait peut-être parti du fond, deux secondes, comme il en part pour tout et pour rien. À la récréation, presque personne ne se serait souvenu de qui avait demandé quoi. Grâce, elle, s’en serait souvenue pendant trois jours. Toute la distance est là, entre le projecteur qu’on imagine et la salle qui existe.

Certains regardent, bien sûr. Mais tu surestimes, comme tout le monde, la durée du regard des autres et la place qu’il prend dans leur journée. Le regard qui dure, c’est le tien, celui que tu portes sur toi-même à l’instant de lever la main. Et c’est lui qui la fait redescendre.

\asterisme

Grâce ignore autre chose encore, et depuis sa place elle ne peut pas le savoir : combien d’autres ont perdu la même ligne. La pluie n’a pas couvert la voix du professeur pour elle seule. Elle l’a couverte pour toute la salle, au même instant, sur la même phrase. Parmi les têtes penchées sur les cahiers, combien recopiaient 3x − 6 sans savoir d’où sortait ce moins ? Personne ne fera jamais le compte. Je parierais pourtant qu’elle n’était pas seule, et je ne prends pas un gros risque. La question qu’on croit bête est très souvent celle d’une bonne partie de la classe, que chacun garde pour soi en se croyant le seul à se la poser.

Celle qui la pose rend donc service. Aux autres d’abord, qui reçoivent la réponse sans avoir eu à lever la main. Au professeur ensuite, qui ne peut pas deviner. Vues de l’estrade, quarante têtes qui recopient se ressemblent toutes : celle qui comprend et celle qui dessine ont la même nuque. Un enseignant qui avance sans recevoir de questions croit, de bonne foi, que tout le monde suit. Une question, même tardive, même maladroite, lui montre où son explication s’est perdue. Des informations de ce genre, il n’en reçoit pas tant que cela.

Tu connais ce moment. Le professeur se retourne, la craie encore levée, et lance : « C’est compris ? » Personne ne répond, et ce silence passe pour un oui. Chacun jette un œil sur ses voisins, voit des visages tranquilles, et en conclut qu’il est le seul à être perdu. Or ces visages tranquilles appartiennent à des gens qui font, au même instant, le même calcul que lui. Une salle entière peut ainsi se persuader, élève par élève, qu’elle a compris ce que personne n’a compris. Dans une classe, le silence est une réponse très peu fiable.

Ce matin-là, il s’est passé une deuxième chose, plus discrète et, à mes yeux, plus grave. Après la ligne des parenthèses, le professeur a enchaîné sur un exercice semblable. Grâce a essayé de suivre. Mais une partie de sa tête était occupée ailleurs : à la main qu’elle n’avait pas levée, à une petite phrase qui tournait en boucle, \emph{tout le monde a compris sauf moi}, et à la crainte sourde d’être envoyée au tableau pour un calcul qu’elle ne savait pas faire. Pendant ce temps, l’exercice avançait sans elle.

Pour comprendre ce qui lui arrivait, il faut revenir à la mémoire de travail, cet espace étroit où l’on pose ce qu’on est en train de penser, le temps de s’en servir. On n’y tient que quelques éléments à la fois. Les travaux de Mark Ashcraft sur ce qu’on appelle l’anxiété mathématique suggèrent que l’inquiétude s’installe précisément là : elle occupe une partie de cet espace, et il en reste moins pour raisonner. Dans certaines de ses expériences, les personnes anxieuses devant les mathématiques peinaient surtout sur les calculs qui obligent à garder plusieurs choses en tête, comme une retenue à reporter, alors que les opérations simples étaient moins touchées. Sian Beilock et ses collègues, qui ont étudié la performance sous pression, ont observé un phénomène proche, avec un détail inattendu : dans certaines de leurs études, c’étaient les personnes dotées de la plus grande mémoire de travail qui perdaient le plus quand la pression montait, parce qu’elles comptaient justement sur elle. Ces résultats ne disent pas tout de l’échec en mathématiques. Ils disent au moins ceci : la peur coûte à la pensée, et ce coût ne mesure en rien l’intelligence de celui qui le paie.

Fais l’essai, si tu veux. Calcule de tête dix-sept fois six pendant que quelqu’un te raconte une histoire à l’oreille. Tu iras plus lentement, et tu te tromperas plus facilement. Maintenant, imagine que l’histoire, c’est toi qui te la racontes, qu’elle parle de ta honte, et que tu ne peux pas lui demander de se taire. Voilà ce que fait la peur pendant un cours de mathématiques. Ton intelligence reste entière ; elle travaille avec quelqu’un qui lui parle à l’oreille.

À force, un cercle se forme. Tu comprends moins bien parce que tu as peur ; tu as peur parce que tu comprends moins bien ; et à force de tourner, tu finis par conclure que tu es nul en maths, quand tu es surtout devenu quelqu’un qui fait des maths en ayant peur. Les deux se ressemblent de loin, et ils se soignent différemment. Contre une incapacité, on ne pourrait pas grand-chose. Contre une peur nourrie d’incompréhension, on peut au moins ceci : retirer un peu d’incompréhension, une ligne à la fois, et la peur perd un peu de son appui.

\asterisme

Le soir, le pont est bouché, évidemment : il a plu toute la matinée et Douala ne s’en est pas remise. Dans le taxi immobile, coincée entre une dame qui ramène du marché un sac de poisson fumé et une portière qui ne ferme qu’à moitié, Grâce regarde un moment le fleuve, couleur de sauce d’arachide, où des pirogues tiennent quand même contre le courant. Le chauffeur a mis la radio, qui annonce un camion en panne à l’entrée du pont, ce que tout le monde dans la voiture avait déjà compris. Elle ouvre son cahier de mathématiques sur ses genoux. 3x − 6. Elle relit la ligne d’avant, puis celle-ci. Elle essaie de retrouver le pourquoi toute seule, en remplaçant x par 5 : neuf d’un côté, neuf de l’autre. Ça tombe juste, mais elle ne sait pas si c’est une raison ou un hasard, et elle s’y perd. Le taxi avance de deux mètres. Elle referme le cahier. \emph{Je verrai plus tard.}

En mathématiques, plus tard a une forme très précise. La semaine suivante, la leçon portera sur des équations plus longues, avec des parenthèses à développer à chaque ligne et des moins cachés devant elles. Dans deux semaines, il y aura la séquence. La ligne recopiée sans être comprise ne restera pas sagement au fond du cahier : elle reviendra, au milieu d’autres exercices, sous d’autres habits, et chaque fois elle fera trébucher. Tu as vu, au chapitre précédent, l’escalier auquel il manque des marches. Voilà comment une marche vient à manquer. Elle ne disparaît pas d’un coup ; elle s’en va un mardi d’octobre, sous la pluie, avec une question qu’on n’a pas posée parce que quarante têtes allaient peut-être se tourner. Et le prix monte avec le temps. Poser la question en octobre coûte trente secondes. La poser en mai, quand tout le reste a été construit dessus, coûte des semaines.

Lever la main, du reste, n’est qu’une façon de demander parmi d’autres : la plus visible et, pour certains, à certaines heures, la plus chère. Il y a la fin du cours : pendant que le professeur efface le tableau, tu peux t’approcher du bureau ; il est moins pressé à 10 h 02 qu’à 9 h 50, et il n’y a plus quarante têtes, seulement la sienne. Il y a aussi le papier qu’on laisse, la question écrite qui attend son lecteur. Le délégué de classe peut porter une question au nom de plusieurs, c’est même un peu son rôle. Un camarade peut te répondre à la récréation, à condition de te souvenir qu’un camarade peut se tromper avec beaucoup d’assurance. Et le professeur se croise ailleurs qu’en classe, dans la cour ou au portail, où une phrase suffit : « Monsieur, juste une question sur la ligne de ce matin. » Tous les enseignants ne répondent pas avec la même patience ; certains diront « pas maintenant », et il faudra alors demander quand. Mais un enseignant, presque toujours, préfère une question tardive à une copie vide.

\asterisme

Dans un cas, rien de ce qui précède ne suffit. Quand le regard des autres devient une moquerie qui revient tous les jours, un surnom qui humilie, des insultes, une mise à l’écart, des messages qui te suivent jusque chez toi sur le téléphone, il ne s’agit plus de timidité ni d’effet projecteur. Cela porte un nom : le harcèlement. Il ne se règle pas en apprenant à lever la main, et il ne se règle pas seul. Il se dit à un adulte de confiance, un parent, un enseignant, le surveillant général, le conseiller d’orientation. Si le premier n’écoute pas, va voir le deuxième. En parler, c’est encore demander, et de toutes les demandes, c’est celle qui compte le plus.

\soustitre{Ce soir, essaie ceci}

Avant le prochain cours de la matière qui te fait peur, prends un petit papier. Pas une feuille entière : un bout, un coin de page, n’importe quoi qui tienne dans une poche. Ouvre ton cahier à la dernière leçon et cherche l’endroit où tu as décroché. Pour le trouver, repère la dernière ligne que tu comprends, puis la suivante, celle que tu as recopiée sans la suivre. Ta question se trouve entre les deux. Écris-la de la façon la plus précise possible : « À la ligne 3, pourquoi le signe change ? » vaut beaucoup mieux que « je n’ai pas compris ». Si tu n’arrives pas à nommer ce qui se passe entre ces deux lignes, recopie-les sur le papier et écris dessous : « Que s’est-il passé entre celle-ci et celle-là ? » C’est déjà une question précise.

Tu n’es pas obligé de la poser en classe. Au prochain cours, donne-la au professeur à la fin, pendant qu’il range ses affaires, ou pose-la-lui à la sortie. Si cela te paraît encore trop, laisse le papier sur le bureau en passant. Ce qui compte cette semaine, c’est que la question quitte ta tête et arrive quelque part.

La semaine suivante, pose-en une à voix haute, pendant le cours. Une seule. Choisis le moment qui coûte le moins, souvent celui où le professeur demande lui-même s’il y a des questions et où le silence s’installe. Tu as le droit de lire ton papier. Personne ne te demande d’improviser.

Cet exercice entraîne la précision. Le courage viendra peut-être avec, ou peut-être pas ; on n’en a pas besoin pour commencer. Une question précise fait moins peur qu’une question vague, pour une raison simple : elle prouve que tu as travaillé. « Je n’ai rien compris » oblige le professeur à tout reprendre, et tu le sens avant même d’avoir fini ta phrase. « Pourquoi le signe change à la ligne 3 » se règle en trente secondes, et plus d’une tête attendait la réponse. Si cette réponse ne t’éclaire pas, ne fais pas semblant d’avoir compris pour libérer le professeur : dis-lui jusqu’où tu l’as suivi et à quel mot tu t’arrêtes encore. Voilà une deuxième question précise, qui vaut la première. Le soir, note dans ton cahier de brouillon la date, la question, et ce qui s’est passé quand tu l’as posée, même si la réponse ne t’a pas éclairé tout de suite.

\asterisme

Ce soir-là, chez elle, pendant que sa mère termine un ourlet à la machine, Grâce sort de son sac un rouleau de tickets de caisse. Elle en a toujours un. Sa tante, qui tient un kiosque de transfert d’argent et de crédit téléphonique, lui donne les rouleaux entamés, et Grâce écrit au dos tout ce qui ne mérite pas une page entière : des listes de courses, des mots d’espagnol, l’ordre des chants du dimanche. Elle déroule une longueur de la largeur d’une main, l’arrache, la retourne. Elle écrit d’abord : « Je n’ai pas compris l’équation. » Elle barre. En dessous, d’une écriture plus petite : « De la ligne 1 à la ligne 2, quand on enlève les parenthèses, pourquoi le moins passe au 6 ? » Elle plie le ticket en deux et le glisse dans son cahier de mathématiques, à la bonne page.

Au cours suivant, le jeudi, elle ne lève pas la main. À la sonnerie, quand tout le monde pousse vers la porte, elle passe devant le bureau et pose le ticket à côté de la boîte de craie, sans un mot, sans s’arrêter. Le professeur, de dos, efface le tableau.

Le lendemain, il entre, pose son cartable, prend la craie et, avant même d’écrire la date, dit : « Quelqu’un m’a demandé hier pourquoi le moins passe au 6 quand on développe. Je reprends. » Il écrit 3(x − 2). Ce matin, il ne pleut pas, et chaque mot arrive jusqu’au fond de la salle. Ici et là, des têtes hochent, une d’abord, puis plusieurs, avec cet air soulagé qu’on prend quand quelqu’un a demandé à voix haute ce qu’on gardait pour soi. Grâce ne lève pas les yeux de son cahier. Du coin de l’œil, sans tourner la tête, elle compte les têtes.

\begin{pourcahier}{Pour le cahier}
Écris la dernière question que tu n’as pas posée. Puis réécris-la de façon plus précise, en nommant la ligne, le mot ou l’étape. Garde-la pour le prochain cours.
\end{pourcahier}


%% 05-le-cahier-du-voisin.md
\chapitre{4}{Le cahier du voisin}

Le sujet est écrit à la craie, en haut du tableau, dans l’angle que le professeur de français essuie toujours en dernier : « La lecture peut-elle remplacer l’expérience ? » Devoir de deux heures, une dissertation, pour la Première F3 du lycée technique de Ngaoundéré. En novembre, ici, le matin arrive frais et repart poussiéreux. Le vent sec qui commence à descendre du nord dépose partout une couche fine, couleur d’ocre : sur les rebords des fenêtres, sur les bancs, sur le sujet lui-même.

Ibrahim Saïdou a recopié le sujet sur sa copie double, puis il a tiré la marge à la règle, en mesurant deux fois, comme il fait tout. Il a écrit une introduction de quatre lignes. Il a commencé un premier paragraphe. Ensuite, la pointe du stylo s’est posée sur la feuille et n’en est plus repartie. Les idées ne lui manquent pas. Il en a une, très précise, qui lui vient des mains : il avait lu dix fois le schéma d’un moteur de ventilateur avant d’en ouvrir un, et le jour où il l’a ouvert, rien ne ressemblait au schéma, ni la couleur des fils ni l’odeur de vernis chauffé. Voilà sa réponse au sujet. Mais il ne sait pas si l’on a le droit de parler de ventilateurs dans une dissertation, ni comment on passe d’un exemple à une idée. Alors il attend que la phrase suivante arrive d’elle-même. Elle n’arrive pas.

À côté de lui, sur le même banc, son voisin en est à la deuxième copie double. Ibrahim ne lit pas ce qu’il écrit ; il n’en voit que la forme. Dans la salle, on n’entend que les stylos et, de loin en loin, le professeur qui tousse dans la poussière en corrigeant les cahiers d’une autre classe. Une écriture serrée, régulière, qui remplit les lignes jusqu’au bord. Des paragraphes séparés par un blanc net, avec un retrait au début, et qui ont l’air de savoir où ils vont. De temps en temps, le voisin tourne une page, et ce petit bruit sec, Ibrahim l’entend mieux que le vent dans les persiennes. Il fait le compte sans le vouloir : une demi-page d’un côté, bientôt six de l’autre. Un rapport de un à douze. C’est un réflexe d’électrotechnicien, il calcule les rapports comme d’autres se rongent les ongles. Le calcul ne l’aide en rien, mais il est juste, et c’est peut-être ce qui le rend si lourd.

Le soir, la famille dîne chez l’oncle, comme souvent quand le père d’Ibrahim est sur la route avec son camion. Dans un coin du salon, depuis la fin de la saison des pluies, un ventilateur sur pied attend qu’on se décide à le jeter. Tout le monde l’a déclaré mort. Même à Ngaoundéré, où les soirées de novembre fraîchissent, l’oncle prétend qu’il ne peut pas s’endormir sans ce bruit-là, et il en parle depuis trois semaines comme d’un ami disparu. Après le repas, Ibrahim s’assoit par terre devant l’appareil. Il dévisse la grille, essuie les pales couvertes de poussière rousse, fait tourner l’axe du doigt : il tourne librement, la panne n’est donc pas mécanique. Il branche. Le moteur ronronne à peine, une sorte de bourdon enfermé, et les pales ne bougent pas. Il les lance d’une pichenette ; elles partent, lentement, puis tiennent leur vitesse. Ibrahim hoche la tête. Le condensateur.

Dans son sac, entre la règle et la trousse, il garde une petite boîte de pièces récupérées sur des appareils que personne ne réclamait. En la sortant, il fait tomber une demi-feuille pliée en deux, son brouillon du matin, qu’il pose sur la table basse sans la regarder. Il choisit un condensateur, lit la valeur imprimée sur le flanc, la relit, la compare à celle de l’ancien. Il raccorde les fils, isole au ruban, rebranche. Cette fois, les pales démarrent seules.

Un murmure fait le tour du salon. L’oncle remercie, sincèrement, avec cette chaleur qu’il met dans tout, et il ajoute aussitôt, assez fort pour que les cousins, les tantes et la mère d’Ibrahim l’entendent : « Bon. Et toi, tu comptes faire comme Moussa ? » Au mur, au-dessus du buffet, dans un cadre doré, le grand frère sourit en blouse blanche, le stéthoscope autour du cou comme une écharpe neuve. Moussa étudie la médecine, loin de la maison. Ibrahim a la grille du ventilateur sur les genoux et quatre vis dans le creux de la main.

\asterisme

Ce jour-là, Ibrahim s’est comparé deux fois : le matin à une copie, le soir à une photo. Il n’a choisi ni l’une ni l’autre. La copie se trouvait à quinze centimètres de son coude ; la photo est accrochée au mur depuis deux ans ; et l’oncle aime les comparaisons autant que le thé très sucré, il en sert à tout le monde.

L’oncle a lancé sa phrase comme on tend un verre, sans méchanceté, avec la conviction de rendre service. Dans sa bouche, Moussa est une preuve : la preuve que la famille peut, qu’un fils de chauffeur peut porter une blouse, et il voudrait que cette preuve serve deux fois. Seulement, une comparaison prononcée devant tout le monde tombe rarement là où on la vise. Celui qui la dit croit encourager ; celui qui la reçoit entend un palmarès, avec un premier et un second, dans un salon où personne n’avait annoncé de concours.

Se comparer est un réflexe, et il n’y a pas de quoi en rougir. En 1954, le psychologue Leon Festinger a proposé une idée qui a beaucoup servi depuis : quand il n’existe pas de mesure simple et objective de ce qu’on vaut dans un domaine, on s’évalue en grande partie en se regardant à côté des autres, et de préférence à côté de ceux qui nous ressemblent un peu. Personne ne se demande s’il court vite en se mesurant à une moto. On se compare au voisin de banc, au cousin, au frère. C’est une façon très ordinaire de savoir où l’on se trouve, et elle a ses usages : sans elle, tu ne saurais même pas qu’une demi-page, pour une dissertation de deux heures, c’est court.

On te conseillera peut-être de ne jamais te comparer à personne. Je ne te donnerai pas ce conseil, parce que je ne saurais pas le suivre moi-même et que tu n’y arriverais pas davantage : autant demander à quelqu’un de ne pas voir ce qu’on a posé sous ses yeux. La comparaison viendra, de toute façon. Reste à savoir sur quoi elle porte, et ce que tu en fais une fois qu’elle est là.

Or elle porte presque toujours sur ce qui se voit. La longueur d’une copie. Une note annoncée dans la cour. Une photo en blouse dans un cadre doré. Ce qui se voit est vrai, mais ne forme qu’une petite partie de ce qui est vrai, et rarement la partie qui explique le reste.

\asterisme

À la fin du devoir, quand le professeur a ramassé les copies, Ibrahim a aperçu le cahier de brouillon de son voisin, resté ouvert sur le banc. Quatre pages couvertes de flèches. Un plan en deux parties, raturé, puis un autre en trois parties, raturé aussi, puis un dernier, encadré deux fois. Trois débuts d’introduction barrés l’un sous l’autre. Dans la marge, des exemples suivis de points d’interrogation. Ce que la copie montrait comme une écriture tranquille avait commencé ici, dans le désordre, par des essais manqués. Le brouillon d’Ibrahim, lui, tenait sur une demi-feuille : le sujet recopié, et une phrase sur un ventilateur. Le voisin a refermé son cahier en grimaçant. « Mon père me donne un sujet chaque dimanche, a-t-il dit. Depuis la quatrième. » Il ne s’en vantait pas. Il avait plutôt l’air de penser à tous les dimanches où il aurait préféré aller au terrain.

Ibrahim avait comparé une demi-page à six pages. La comparaison juste aurait été tout autre : une demi-page écrite sans plan, contre six pages précédées de quatre pages de brouillon et de trois années de dimanches. Le calcul change du tout au tout. Ce qu’il prenait pour une différence de tête était, au moins en partie, une différence d’heures, d’entraînement, et d’un père qui sait ce qu’est une dissertation.

La photo de Moussa fonctionne de la même manière. Elle montre une blouse. Elle ne montre pas les nuits de révision, les livres empruntés, les examens où il a peut-être tremblé, ni ces dimanches où, au téléphone, il parle surtout de sommeil parce qu’il en manque. Rien de tout cela ne retire quoi que ce soit à Moussa : il a travaillé, il travaille encore, et sa blouse, il l’a gagnée. Seulement, une photo est faite pour montrer le résultat, jamais la fabrication. Mettre sa vie de tous les jours à côté de la photo de quelqu’un, c’est poser son brouillon à côté de la copie au propre d’un autre.

Les écrans font ce travail à grande échelle. Ils montrent des vitrines et presque jamais d’arrière-boutiques : la capture des résultats, sans la soirée passée à pleurer devant l’exercice. Tu le sais déjà ; il suffit de t’en souvenir au moment où tu regardes.

La comparaison cache encore autre chose : elle dépend de qui se trouve autour de toi. Le chercheur Herbert Marsh et ses collègues ont étudié, à partir des années 1980, ce qu’ils ont appelé l’effet du gros poisson dans une petite mare. Leurs travaux, menés dans beaucoup de pays, suggèrent qu’à niveau égal un élève a tendance à se juger moins bon dans une classe forte, et meilleur dans une classe plus faible. Le même élève, avec les mêmes connaissances, change d’image de lui-même en changeant de banc. Ce que tu crois valoir dépend donc en partie de la mare où l’on t’a mis, et pas seulement de toi. Chez l’oncle d’Ibrahim, la mare est un salon, et un très gros poisson en blouse blanche nage au mur.

\asterisme

Ibrahim, pourtant, sait comparer utilement. Il l’a appris ailleurs qu’en classe. Quand il a commencé à bricoler, vers douze ans, il allait regarder travailler un vieux réparateur installé sous un auvent, près du marché. De cet homme, il ne convoitait ni la boutique ni la moustache, et il ne rêvait pas de lui ressembler. Il voulait son ordre : vérifier d’abord la prise, puis le cordon, puis l’interrupteur, et seulement ensuite ouvrir le moteur. Le vieux ne lui a jamais rien expliqué. Il travaillait, Ibrahim regardait, et un jour Ibrahim a fait pareil sans y penser. De cet homme, il avait gardé la méthode, et rien d’autre.

On peut poser deux questions à quelqu’un qui réussit. « Combien tu as eu ? » te donne un nombre, que tu vas aussitôt poser à côté du tien et qui te renverra à ton rang. « Comment tu travailles ? » te donne une façon de faire à emporter chez toi, et une façon de faire se copie sans rien voler à personne. En refermant son cahier, le voisin d’Ibrahim lui a donné sans le savoir une information précieuse : un plan au brouillon avant d’écrire, et un sujet par semaine. Ibrahim ne peut pas emprunter le père, ni les trois années de dimanches. Il peut emprunter le brouillon.

Encore faut-il l’emprunter à sa taille. Trois années de dimanches, Ibrahim ne les rattrapera pas en un mois, et personne chez lui ne lui donnera de sujet. Mais avant le prochain devoir de français, il peut faire ce que fait n’importe quel réparateur devant une panne qu’il ne connaît pas : poser les pièces sur la table avant de remonter quoi que ce soit. Écrire au brouillon l’idée du ventilateur, en chercher une deuxième qui dise le contraire, les mettre dans un ordre, et seulement ensuite ouvrir la copie double. Il en sortira peut-être une page au lieu d’une demi-page, et une page qui sait où elle va. Une méthode copiée rétrécit presque toujours en route ; c’est normal, elle doit passer par ta vie avant de servir.

Il existe enfin une comparaison qui trompe beaucoup moins que les autres : celle qui te met en face de toi-même, quelques mois plus tôt. Ibrahim la pratique sans le savoir. Le premier ventilateur qu’il a réparé lui a pris un samedi entier et une brûlure dont il garde, au pouce gauche, une petite cicatrice blanche ; celui de l’oncle, ce soir, lui a demandé quarante minutes. Personne n’a eu besoin de lui dire qu’il avait progressé : il connaissait les deux durées. En français, il n’a rien de tel. Aucune trace de sa dissertation d’il y a trois mois, aucun point de départ noté quelque part. Il ne peut se mesurer qu’au voisin, seule mesure qu’il ait sous la main. Ton cahier de brouillon sert à cela : garder le point de départ, pour qu’un jour tu aies autre chose à quoi te comparer que la copie d’à côté.

On pourrait, pour finir, ranger Ibrahim dans deux cases bien commodes : nul en français, fort en technique. Lui-même le fait. Quand on l’interroge sur ses notes, il dit qu’il est nul en français, lentement, avec un sourire tranquille qui ne trompe personne, et surtout pas lui. Mais les deux cases mentent ensemble. Ce soir, devant le ventilateur, il a fait ce que demandait le sujet du matin : il a confronté ce qu’il avait lu sur les condensateurs à ce que ses mains trouvaient dans le moteur. Il a observé, supposé, vérifié, conclu. Il a la pensée ; il lui manque, devant la copie double, les outils pour la poser sur le papier : bâtir un plan, passer d’un exemple à une idée, lier deux paragraphes. Ceux du diagnostic, il les a ; ceux de la dissertation, il ne les a pas encore. Les uns comme les autres s’apprennent, et personne n’est né avec un tournevis dans la main. Où tout cela le mènera, je n’en sais rien, et je laisse la question à Ibrahim.

\soustitre{Ce soir, essaie ceci}

Choisis une personne qui réussit dans la matière où tu peines. Pas forcément la première de la classe : quelqu’un d’accessible, à qui tu peux parler demain sans cérémonie, un camarade de banc, une cousine en classe supérieure, un ancien du quartier. Ce soir, contente-toi de choisir, et d’écrire son nom dans ton cahier.

Demain, pose-lui une seule question, la même pour tout le monde, mot pour mot si tu veux : « Comment tu t’y prends, concrètement, la veille d’un devoir ? » Le mot qui compte est « concrètement ». Si la personne répond « je révise », demande-lui ce qu’elle fait de ses mains pendant qu’elle révise : si elle relit ou refait des exercices, à quelle heure elle commence, si elle travaille seule ou avec quelqu’un, combien de temps elle y passe. Les réponses vagues ne t’apprendront rien ; les gestes, si.

Pendant qu’elle parle, écoute sans comparer. C’est plus difficile qu’il n’y paraît, parce que tu vas entendre des choses que tu n’as pas (un père qui donne des sujets le dimanche, un répétiteur, une chambre calme) et qu’elles te donneront envie de cesser d’écouter. Laisse-les passer. Sa vie est peut-être plus facile que la tienne ; garde cette question pour un autre jour, et cherche un geste que tu pourrais refaire chez toi, dès demain soir.

Rentre et choisis, dans tout ce qu’elle t’a dit, une seule habitude, même si d’autres te tentent. Faire un plan au brouillon avant d’écrire. Refaire le dernier exercice du cours sans regarder la correction. Se coucher plus tôt la veille. Peu importe laquelle, pourvu qu’elle soit précise et ne coûte rien. Essaie-la pendant une semaine, chaque fois que l’occasion se présente, et note dans ton cahier les jours où tu l’as faite et ceux où tu l’as oubliée.

Au bout de la semaine, avant de décider si l’habitude « marche », ouvre ton cahier à la première page. Relis la phrase que tu y avais écrite, celle que quelqu’un a dite sur toi. Ne la corrige pas. Feuillette ensuite les pages qui suivent : la note que tu as regardée de près, l’escalier que tu as dessiné, la question que tu as réécrite. Cherche ce qui a bougé depuis, même d’un millimètre. C’est la comparaison la plus honnête que tu puisses faire, puisque tu en tiens les deux bouts.

\asterisme

Dans le salon de l’oncle, la question flotte encore au-dessus du ventilateur. Les cousins attendent une réponse ; la mère d’Ibrahim regarde ses mains. Moussa sourit toujours dans son cadre doré, et il n’y est pour rien. Ibrahim présente la grille devant le moteur, aligne les trous, engage une vis, puis une autre. Il pourrait dire qu’il est nul en français, avec son sourire. Il pourrait dire qu’il ne sait pas encore. Il ne dit rien.

Le ventilateur tourne. La fraîcheur de novembre n’en demandait pas tant, mais l’oncle ferme les yeux de contentement et oublie un instant qu’il attendait quelque chose. Ibrahim serre la dernière vis, d’un quart de tour de plus, pour être sûr. Sur la table basse, à côté des verres de thé, la demi-feuille du matin est restée pliée en deux. L’oncle, sans la voir, pose son verre dessus.

\begin{pourcahier}{Pour le cahier}
Écris le nom d’une personne à qui tu te compares souvent. Sous son nom, deux colonnes en prose : ce que tu vois d’elle, ce que tu ne vois pas. La deuxième doit être plus longue que la première. Si elle reste courte, c’est qu’il te manque une question à lui poser.
\end{pourcahier}


%% 06-la-voix-de-dedans.md
\chapitre{5}{La voix de dedans}

En janvier, à Bafoussam, le froid arrive avec la nuit. Il descend des collines dès que le soleil tombe, entre dans les maisons par les fentes des persiennes et s’installe sur les draps, si bien qu’on se couche en pull. Divine Nfor a gardé celui du handball, aux manches trop longues, qu’elle tire sur ses mains. Le dîner est fini depuis longtemps. Sa mère, rentrée tard de sa garde à l’hôpital, lui a demandé comment s’était passée la journée ; Divine a répondu « fine », puis, parce que sa mère attendait encore, « ça va ». Deux mots pour la même chose, comme on paierait deux fois le même article dans deux monnaies différentes. Son père, au bout de la table, refaisait des colonnes de chiffres sur un bloc de bureau. De la dictée, il n’a pas été question.

Elle est dans le sac, pliée en quatre, entre le cahier de mathématiques et la trousse. Un texte d’une quinzaine de lignes sur un marché du samedi, lu deux fois par le professeur de français, la seconde fois plus lentement. Divine a écrit vite, en serrant le stylo, en entendant chaque mot une demi-seconde trop tard, comme s’il lui parvenait d’une autre pièce. Ce matin, la copie est revenue avec vingt-trois petits cercles rouges. Le rouge entoure, sans un mot. Vingt-trois fois, net, rond, chaque fois au bon endroit, ce qui est presque pire. La note, en bas de la page, elle ne l’a pas regardée longtemps. Elle la connaissait avant de la lire.

La famille est arrivée de Bamenda il y a trois ans. Divine pense en anglais. Elle compte en anglais, y compris les moutons, les nuits où le sommeil traîne, et il lui arrive de rêver en pidgin. Mais elle passe ses journées en français. Elle le parle bien, elle le comprend presque toujours ; c’est à l’écrit qu’il lui tend des pièges, des pièges qu’elle ne voit qu’après, quand quelqu’un les a entourés.

Elle éteint. La maison fait ses bruits de nuit : le réfrigérateur qui démarre et s’arrête, une moto au loin qui monte la côte en peinant, la radio d’un voisin qui ne dort jamais. Et dans le noir, la voix commence.

\emph{Twenty-three. Twenty-three mistakes. I can’t even write down what someone reads to me.}

Elle parle d’abord en anglais, parce que c’est la langue la plus proche de l’os. Puis, au bout d’un moment, sans prévenir, elle passe au français, comme pour être sûre d’être comprise des deux côtés de la tête.

\emph{Vingt-trois. Même pas capable d’écrire ce qu’on me dicte. Je suis nulle. Je suis nulle, c’est tout.}

Dans les deux langues, elle dit la même chose. La voix est bilingue, elle aussi ; elle a eu trois ans pour apprendre. Divine se tourne vers le mur. Elle sait qu’en décembre elle avait seize en mathématiques, l’une des meilleures notes de la classe, et qu’elle en a été contente pendant au moins deux jours. La voix ne le mentionne pas. Elle ne lit que le rouge.

\asterisme

Cette voix, on la ramasse en chemin, morceau par morceau. Elle est faite de choses entendues : la remarque d’un professeur pressé, le soupir d’un parent fatigué, le rire d’un banc entier, un mot écrit un jour en capitales sur un cahier qui n’était pas le sien. Ces mots ont été dits une fois, souvent sans y penser, par quelqu’un qui les a oubliés le lendemain. Celui qui les a reçus les garde. Il les répète la nuit, de sa propre voix, et au bout de quelques années il ne sait plus qu’ils viennent d’ailleurs. Il les prend pour ses pensées.

Si Divine écoutait bien, elle reconnaîtrait un accent, celui d’une fille de sa première classe à Bafoussam qui, le premier mois, avait répété en l’imitant la façon dont Divine prononçait « heureusement ». Tout le banc avait ri. Divine aussi, un peu après les autres, pour ne pas être la seule à ne pas rire. Une seconde dans une matinée. La fille a changé de lycée l’année suivante ; elle ne se souviendrait sans doute pas de la scène si on la lui racontait. Mais certaines nuits, quand la voix dit \emph{nulle}, elle le dit avec cette petite pointe moqueuse qui traîne sur la fin du mot, et Divine ne comprend pas pourquoi ce mot-là lui fait plus mal en français.

Un mot écrit sur un cahier, on peut au moins le voir. Il est là, à l’encre, sur la couverture ; on peut le montrer à quelqu’un, le barrer, coller quelque chose par-dessus. La voix de dedans n’a pas de papier. Elle se réécrit chaque soir, et chaque soir de ta propre écriture, ce qui lui donne une autorité terrible : comment douter de ce que tu prononces toi-même ? À force de revenir, elle devient familière, comme une chanson qu’on finit par fredonner sans l’avoir aimée, et la familiarité se fait passer pour la vérité. Ce que tu te dis la nuit ne prouve pas ce que tu crois. La chanson a été écrite par d’autres, et tu la chantes depuis si longtemps que tu crois l’avoir composée.

Remarque aussi l’heure à laquelle elle parle. Le jour, elle a de la concurrence : le professeur au tableau, les camarades sur le banc, la file à la cantine, le ballon du mercredi qui ne laisse pas le temps de penser à autre chose. Le soir, tout cela se tait. Il ne reste qu’elle, la radio du voisin et le plafond. Elle a la salle pour elle seule, comme un orateur sans contradicteur. C’est pour cela qu’une copie supportable à midi devient, vers 22 h, un verdict sur toute une vie. Entre-temps, la copie est restée pliée dans le sac avec ses vingt-trois cercles ; seule la salle s’est vidée.

Faire taire cette voix, je ne sais pas, et je me méfie de ceux qui prétendent savoir. On peut en revanche l’écouter de plus près, comme on relit une copie avant de la rendre, en cherchant où elle se trompe.

\asterisme

Des psychologues ont beaucoup travaillé sur une question d’apparence simple : l’histoire que l’on se raconte après un échec. Ils ont observé que cette histoire pèse, parfois autant que l’échec lui-même. Une explication peut placer la cause en toi ou hors de toi. Elle peut la rendre durable ou passagère, la faire porter sur un point précis ou sur tout, et te laisser une prise ou aucune. Selon l’histoire qu’on s’est racontée le soir, on ne fait pas la même chose le lendemain matin.

Prends ce que dit la voix : \emph{je suis nulle}. Regarde-le comme on regarde une explication. La cause est en Divine, entièrement. Elle est permanente : le verbe est « être », au présent, sans date. Elle porte sur tout, Divine en entier, mathématiques comprises, alors que la copie ne parlait que d’orthographe. Et elle ne laisse aucune prise, puisqu’on ne peut rien faire contre ce que l’on est. C’est une explication parfaite, à la manière d’un couvercle bien vissé : rien ne sort, pas même la question suivante.

Essaie maintenant celle-ci : \emph{je n’ai pas encore automatisé l’accord du participe passé avec « avoir » quand le complément est placé avant}. Elle est plus longue, moins jolie, et franchement moins dramatique. Personne ne pleure dans son lit pour une phrase pareille. Mais elle dit où se trouve le problème, depuis quand à peu près il dure, et ce qu’il faudrait faire pour qu’il cesse. Elle dévisse le couvercle d’un quart de tour. La seconde phrase n’est pas plus gentille que la première. Elle est plus exacte.

Être précis avec soi coûte quelque chose. Il faut regarder les fautes une par une au lieu de se noyer dans le total, et reconnaître ce qui dépend de toi. « Je n’ai pas encore automatisé » laisse tranquilles le professeur, la langue et le déménagement, et désigne ce qui manque, à l’endroit où ça manque.

Le plus étrange, c’est que Divine sait déjà faire cela, ailleurs. Le mercredi, au handball, quand son tir passe à côté du but, elle se dit : trop à gauche, ou trop tôt, ou le bras pas assez haut. Elle le dit en anglais, le plus souvent, et elle reprend la balle pour recommencer. Le mot « nulle » n’entre jamais sur le terrain. Personne ne lui a appris à se parler ainsi ; c’est venu avec les heures. La même fille, devant une feuille, perd ce réflexe d’un coup, comme si le stylo la privait de son vocabulaire.

Dans le noir, Divine refait la dictée de mémoire. Elle se souvient de quelques cercles. « Les mangues qu’elle avait cueilli », sans e : elle entend encore la phrase dans la bouche du professeur, elle se souvient d’avoir hésité, d’avoir pensé, sans le formuler, à l’anglais, où le participe ne s’accorde avec rien ni personne, et d’avoir choisi l’anglais sans le savoir. « Addresse » avec deux d, parce que \emph{address}. « Example » avec un a, parce que \emph{example}. Ces deux fautes-là sont peut-être les seules de la page qui viennent d’un savoir en trop plutôt que d’un savoir en moins : celle qui écrit \emph{example} sait déjà écrire ce mot dans une autre langue, il ne lui reste qu’à apprendre où passe la frontière. La voix, elle, ne fait pas ce genre de tri. Pour elle, vingt-trois est un bloc. Vingt-trois, dit dans le noir, pèse comme une seule pierre.

Ce qui arrive quand cette pierre revient tous les soirs, des chercheurs l’ont étudié, avec des détours qui comptent. Dans les années 1960, Martin Seligman et Steven Maier ont observé, au cours d’expériences sur des animaux, que ceux qui avaient d’abord subi des situations désagréables sans aucun moyen de les faire cesser finissaient par ne plus essayer, même plus tard, quand une issue était à leur portée. Ils ont appelé cela l’impuissance apprise : à force de constater que rien de ce qu’on fait ne change rien, on apprendrait à ne plus rien faire. L’idée a été reprise partout, jusque dans les salles des professeurs. Puis, en 2016, près de cinquante ans plus tard, les deux mêmes chercheurs sont revenus sur leur théorie à la lumière de travaux plus récents. Leur nouvelle lecture, pour la dire simplement et avec prudence : la passivité serait plutôt la réaction de départ face à une difficulté qui dure, et ce qui s’apprend, c’est la sortie. Il faut avoir vu, au moins une fois, que ce qu’on fait change quelque chose.

J’aime qu’une théorie célèbre ait été corrigée par ceux-là mêmes qui l’avaient construite, et j’aime ce que la correction suggère pour toi. Si chaque dictée est vécue comme une sentence sur laquelle tu n’as aucune prise, il est presque normal que tu cesses d’essayer. Vu de loin, on appellerait cela de la paresse. N’importe qui, après avoir poussé assez longtemps un mur qui ne bouge pas, laisse retomber les bras. Pour retrouver de la prise, un discours pèse peu, et l’élan de motivation qui tient jusqu’au lendemain midi encore moins. Il faut une poignée. Une règle, comprise, appliquée, vérifiée sur la dictée suivante. Un cercle rouge de moins, à un endroit précis, dont tu sais pourquoi il n’y est plus.

\asterisme

Autre détail : la voix de Divine dit « je ». \emph{Je suis nulle.} Les reproches les plus durs, on se les fait souvent à la première personne, de l’intérieur, collé contre la faute, sans aucun recul, comme quand on essaie de lire une page le nez sur le papier.

Des expériences menées par une équipe de psychologues américains suggèrent que, dans les moments difficiles, se parler à soi-même en employant son prénom, ou « tu », plutôt que « je », aide à prendre de la distance : on regarde alors sa situation à peu près comme on regarderait celle d’un autre, et l’on y voit plus clair. L’effet n’a rien de magique et ne remplace rien. Mais l’essai coûte une phrase. \emph{Je suis nulle} devient : « Divine, tu as fait vingt-trois fautes, et beaucoup se ressemblent. » Cela, on pourrait le dire à une amie sans la blesser ni lui mentir. C’est peut-être une raison de plus pour laquelle ce livre te dit « tu » : il essaie, de l’extérieur, de te tenir à la bonne distance de ta propre copie.

Dans son lit, Divine n’a rien lu de tout cela. Mais il lui arrive, sans le savoir, de faire l’essai. \emph{Divine, you made twenty-three mistakes.} Cela sonne bizarrement, presque comme si sa mère le prononçait, avec ce calme d’infirmière qui a vu pire. Ce n’est pas gai pour autant, et cela laisse un peu d’air entre Divine et la copie, assez pour qu’une question passe : lesquelles, au juste.

Tu as peut-être déjà croisé, sur une affiche, dans une vidéo ou dans la bouche d’un adulte bien intentionné, la « mentalité de croissance ». L’idée vient des travaux de la psychologue Carol Dweck. Dans sa forme la plus simple, elle oppose deux manières de voir : penser que l’intelligence est une quantité fixe, distribuée une fois pour toutes, ou penser qu’elle se développe avec le travail. Sur le fond, pour ce qui compte à l’école (connaissances, méthodes, savoir-faire), la seconde vision est plus juste que la première : ces capacités-là se travaillent ; tout ce livre repose là-dessus.

Mais l’idée a connu un destin de slogan, et la recherche est plus réservée que les affiches. Des études menées à grande échelle, et des synthèses qui en rassemblent beaucoup d’autres, suggèrent que les interventions destinées à enseigner cette mentalité ont, en moyenne, des effets modestes sur les résultats scolaires. Ces effets semblent plus nets chez certains élèves en difficulté, et dans des établissements où l’entourage rend vraiment l’effort possible, où l’on a le droit de se tromper sans être moqué. Ailleurs, ils sont faibles, parfois à peine mesurables.

Le slogan peut même faire pire que rien. Quand on dit à un élève qui échoue qu’il « n’a pas le bon état d’esprit », on lui fabrique une sentence de plus. Hier il était mauvais en orthographe ; aujourd’hui il est mauvais en mentalité. La voix de dedans adore ce genre de cadeau. Elle range la nouvelle formule à côté des anciennes, bien pliée, et la ressort la nuit.

Ce livre n’a pas besoin du slogan. Il n’a pas besoin non plus que tu te répètes devant un miroir des compliments en espérant que la voix finira par y croire. Elle n’y croira pas : repérer un mensonge, c’est sa spécialité. Remplacer une phrase fausse par une autre, fausse elle aussi et plus douce, revient à lui donner raison. Ce dont tu as besoin, c’est de méthodes et de vérité. La vérité, dans le cas de Divine, tient en une ligne : vingt-trois fautes, dont beaucoup du même genre, dans une langue qu’elle écrit tous les jours depuis trois ans seulement. Cette ligne ne console personne, et on peut travailler dessus.

\soustitre{Ce soir, essaie ceci}

Prends ton cahier de brouillon et ouvre-le à une page neuve. En haut, écris la date.

Puis écris, mot pour mot, ce que ta voix intérieure t’a dit aujourd’hui. Écris-le tel quel, avec ses mots à elle, même s’ils sont durs, même si elle te l’a dit dans une autre langue que le français. Si elle t’a dit plusieurs choses, choisis celle qui est revenue le plus souvent. Écris-la, et ne la commente pas.

En dessous, réécris-la en la rendant précise. Quatre questions t’y aident : quoi, dans quelle matière, sur quel point, depuis quand. « Je suis nul en anglais » deviendra peut-être « Je confonds encore le prétérit et le \emph{present perfect}, depuis le début de l’année ». « Je ne comprends rien en physique » deviendra « Je ne sais pas quelle formule choisir dans les exercices sur la vitesse, depuis la dernière séquence ». Si tu n’arrives pas à répondre à l’une des quatre questions, note-le aussi : c’est une information, et elle te servira.

Ensuite, la troisième ligne. Pense à ton meilleur ami, ou à ta meilleure amie, qui viendrait s’asseoir à côté de toi avec la même copie, les mêmes cercles rouges, la même note. Tu ne lui mentirais pas, tu tiens trop à lui pour ça ; tu lui dirais sans doute ce qui ne va pas, et par où commencer. Écris cela, en lui disant « tu ».

Relis maintenant les trois, lentement, l’une après l’autre. Tu remarqueras peut-être que la première est la plus courte, la plus sûre d’elle, et la moins renseignée. Elle juge sans avoir enquêté. La deuxième est plus longue et plus ennuyeuse ; la troisième est souvent plus juste qu’on ne l’aurait cru. La plus vraie est rarement la plus douce, et presque jamais la première.

Demain soir, recommence avec ce que la voix t’aura dit dans la journée. Ne cherche pas à faire disparaître la première ligne : elle reviendra. Tu t’entraînes seulement à ne plus la laisser écrire seule. Dans trois jours, relis les trois pages d’un coup. Il se peut que la même ligne revienne chaque soir, avec le même flou. C’est souvent le signe qu’elle désigne une seule marche, toujours la même, et que cette marche mérite qu’on s’en occupe pour de bon.

\asterisme

Parfois la voix ne se contente plus de commenter une dictée. Elle ne s’arrête plus, ni le soir ni le matin. Le sommeil s’en va, ou au contraire on ne voudrait plus faire que dormir. Plus rien n’a de goût, ni la nourriture, ni le sport, ni les amis. Et parfois vient l’idée de disparaître, ou de se faire du mal. Si c’est ton cas, ou si cela y ressemble, même de loin, on a quitté le terrain de la méthode, et aucun exercice de ce livre ne suffira. Parles-en aujourd’hui, sans attendre la fin des examens : à un médecin, à l’infirmerie du lycée ou au centre de santé, ou à un adulte de confiance, un parent, un professeur, un responsable religieux si c’est quelqu’un à qui tu te fies. Tu n’as pas besoin de trouver les bons mots ; « ça ne va pas, et ça dure » suffit pour commencer. Si la personne à qui tu parles ne comprend pas, va en voir une autre. Demander de l’aide, dans ce cas-là, c’est faire ce qu’il faut.

\asterisme

Il est plus de 23 h quand Divine se relève. Le froid a gagné le sol de la chambre ; elle cherche ses sandales du bout du pied, ne les trouve pas, renonce. Elle prend son téléphone sur la chaise et allume la lampe torche. Le faisceau est blanc, étroit, un peu cru. Il passe sur le sac, sur la porte de l’armoire, puis sur l’étagère où, sous deux manuels, dort un carnet bleu qu’elle avait acheté au marché à la rentrée pour y recopier des paroles de chansons. Une seule page a servi. Elle la tourne.

Elle s’assoit par terre, le dos contre le lit, le téléphone coincé entre son genou et le carnet pour que la lumière tombe au bon endroit. Elle déplie la dictée. Les vingt-trois cercles sont toujours là, aussi rouges qu’à midi. Elle commence par le premier. Sur une ligne, elle recopie la phrase fautive telle qu’elle l’a écrite ; à côté, la règle. Pour la règle, elle hésite, puis l’écrit d’abord en anglais, pour elle, parce qu’en anglais elle la comprend du premier coup : \emph{the past participle agrees with the object when the object comes first}. En dessous, plus petit, la même règle en français.

Deuxième ligne. Troisième. À la quatrième, elle s’arrête : c’est la même faute que la première, sur un autre verbe. Elle trace une petite flèche de l’une à l’autre. À côté d’« addresse », elle écrit \emph{address} d’un côté, \emph{adresse} de l’autre, et entoure le second d de l’anglais, celui que le français n’a pas, au stylo bleu, sans appuyer.

La voix n’est pas partie. Elle parle encore, plus bas, comme la radio du voisin à travers le mur : on l’entend, on ne distingue plus tous les mots. À la septième ligne, le téléphone signale que la batterie faiblit. Divine éteint la lampe torche, glisse la dictée dans le carnet en guise de marque-page et pose le tout sur la chaise. Il reste seize lignes.

\begin{pourcahier}{Pour le cahier}
Recopie ta phrase intérieure du jour. En dessous, réécris-la avec un « pas encore » et un détail précis. Exemple de forme : « Je ne sais pas encore accorder les participes avec avoir. »
\end{pourcahier}


%% 07-redescendre.md
\partie{Le travail}{II}

\chapitre{6}{Redescendre}

Le samedi, la boutique de tonton Fabrice ouvre avant le marché. C’est une cabine de planches peintes en bleu, au carrefour d’Ebolowa, serrée entre une vendeuse de beignets et un étal de chaussures d’occasion. Derrière la vitre du comptoir dorment des chargeurs emmêlés et des téléphones qui attendent qu’on revienne les chercher. Sur le comptoir lui-même, une radio branchée sur une rallonge, une lampe à pince et un tapis de caoutchouc noir où chaque vis, une fois sortie, a sa place et n’en change plus. Rodrigue y passe ses samedis. Il a quatorze ans, un petit tournevis dans la poche de sa chemise, et l’habitude de se taire quand il travaille.

Ce matin de février, en pleine saison sèche, une commerçante du marché pose sur le comptoir un téléphone qu’elle tient à deux mains, comme on porte un poussin malade. « Il ne s’allume plus. Depuis hier. Rien du tout. » Elle l’a secoué, laissé en charge toute la nuit, et elle a même prié dessus, ce qu’elle raconte sans sourire. Tonton Fabrice, penché sur une soudure, lève le menton vers son neveu. Cela veut dire : à toi.

Rodrigue ne dit pas que le téléphone est mort, ni fichu, ni nul. Ces mots-là, ce sont les clients qui les emploient, et ils arrivent à la fin, quand on n’a pas cherché. Il le branche d’abord. L’écran reste noir, mais l’appareil vibre une fois, faiblement, contre la vitre du comptoir. Quelque chose, à l’intérieur, est vivant. Il passe tout de même la batterie au multimètre, parce qu’un téléphone qui vibre a probablement une batterie correcte, et que « probablement » n’a jamais rien réparé. La batterie est correcte. Il ouvre la coque avec la spatule en plastique, dévisse la plaque, la soulève. Dessous court la nappe de l’écran, ce ruban souple et plat qui relie l’affichage au reste de l’appareil, et au bout de la nappe, le connecteur. Sur le connecteur, une croûte verdâtre, à peine visible, comme celle qui ronge les vieilles pièces de monnaie.

« Il est tombé dans l’eau, maman ? » La commerçante réfléchit, puis se souvient : la bassine où elle rince ses légumes, avant-hier, une seconde à peine, elle l’a repêché tout de suite. Rodrigue frotte le connecteur avec un coton imbibé d’alcool, puis avec une vieille brosse à dents dont il a coupé la moitié des poils. Il attend que ça sèche. Il remet la nappe et appuie jusqu’au petit clic. Il ne referme pas encore : il rallume l’appareil tel quel, ouvert sur le tapis, les entrailles à l’air. L’écran s’éclaire. Le logo apparaît, blanc sur noir, puis l’image de fond, une photo d’enfants en uniforme devant une porte. Alors seulement il revisse, reclipse la coque, vérifie que l’écran répond au doigt et que la charge se fait. La commerçante compte deux fois sa monnaie et s’en va en parlant à son téléphone comme à un enfant qu’on a retrouvé au marché.

Tonton Fabrice ne dit ni « bien » ni « bravo ». Sans lever les yeux de sa soudure, il tend à son neveu le téléphone suivant, dont le haut-parleur grésille.

\asterisme

Ce que Rodrigue vient de faire, tonton Fabrice l’appelle « remonter ». On part de ce qu’on voit, l’écran noir, et on remonte vers ce qu’on ne voit pas, pièce par pièce : la charge, la batterie, la nappe, le connecteur, et en dernier l’écran lui-même, qui coûte le plus cher et qu’on ne remplace qu’après avoir innocenté toutes les autres pièces. On teste une chose à la fois. Changer trois pièces d’un coup en espérant que l’une d’elles était la bonne, c’est s’interdire de savoir pourquoi le téléphone remarche, s’il remarche. Et quand on a trouvé, on répare en bas, là où est la panne, mais on vérifie en haut, là où était le symptôme. À la fin, la seule preuve qui compte, c’est l’écran qui s’allume.

Le lundi, le même garçon est assis en classe devant un exercice de calcul littéral, et il fait tout à l’inverse. Il bloque à une ligne, essaie plusieurs choses à la fois, rature, recommence en haut de la page, tente autre chose, abandonne. Puis il conclut, comme la commerçante : rien du tout. D’un téléphone éteint, il ne dirait jamais qu’il est nul ; de lui-même, devant une copie, il le pense sans y penser. Il possède déjà la méthode. Il la range simplement dans la mauvaise poche, celle qui ne s’ouvre que le samedi.

Je trouve cela à la fois rassurant et un peu vexant, comme toutes les bonnes nouvelles qui arrivent en retard. Beaucoup d’élèves qu’on dit nuls savent déjà, ailleurs, chercher une panne avec méthode : réparer une moto, retrouver l’erreur dans les comptes d’une boutique, comprendre pourquoi la pâte des beignets n’a pas levé et ne changer qu’une seule chose le lendemain. Ils ont la démarche. Il leur manque l’idée qu’elle s’applique aussi à une page de cahier.

Si tu as dessiné l’escalier de ta matière la plus difficile et entouré la marche où tu n’es plus sûr, tu as fait la moitié du travail de Rodrigue : tu es remonté du symptôme jusqu’à la panne. L’autre moitié, la réparation, commence ici. Les mots se croisent alors de façon un peu cocasse, car ce qui s’appelle « remonter » dans la boutique s’appelle « redescendre » à l’école. Il s’agit du même geste, vu des deux bouts : pour atteindre la cause, il faut revenir en arrière dans les classes, jusqu’à la leçon qu’on aurait dû savoir il y a deux ou trois ans. Pour Rodrigue, la panne tient en une ligne qu’il n’arrive pas à franchir depuis novembre, 2/3 + 1/4. Sous cette ligne dorment les fractions de sixième, et sous elles, peut-être, celles du CM2.

C’est là que la honte s’installe, plus lourde que la difficulté. Ouvrir en troisième un manuel de CM2 ressemble à un aveu. Rodrigue a longtemps tourné autour avant de demander à sa cousine Bella de lui prêter le sien ; elle a accepté à condition qu’il ne corne pas les pages, et elle a vérifié le soir même. Le premier soir, il est tombé sur une page où un gâteau dessiné était coupé en parts égales, avec des couleurs pour celles qu’on avait mangées, et il a refermé le livre aussitôt, vexé comme si quelqu’un l’avait surpris à jouer avec les jouets de Bella. Il l’a rouvert le lendemain. L’explication du gâteau tenait en six lignes, et ces six lignes disaient clairement ce qu’aucun cours de troisième ne prend plus le temps de dire, puisque tout le monde est censé le savoir depuis longtemps : on ne peut additionner des parts que si elles ont la même taille. Ce que la honte oublie de dire, c’est qu’à la table de la cuisine, à 21 h, il n’y a aucun témoin. Le manuel du petit frère, celui de la petite cousine, le vieux cahier d’un voisin passé en seconde : ces livres ne savent pas qui les lit. Ils expliquent la marche à quiconque s’assoit devant eux, dans l’ordre où elle doit être expliquée, avec des exemples faits pour quelqu’un qui découvre. Il te faut précisément cela.

Encore faut-il redescendre au bon rythme. Une fois la marche trouvée, la tentation est de vouloir la réparer d’un coup : un dimanche entier, trois heures, le manuel de sixième de la première à la dernière page. Cela ne marche presque jamais. On ne tient pas trois heures sur une chose qu’on déteste, et ce qu’on entasse en une seule séance s’en va plus vite que ce qu’on reprend un peu chaque jour (le chapitre suivant reviendra longuement sur ce point). Un quart d’heure par jour, sur une seule marche, vaut mieux. C’est moins héroïque. Tonton Fabrice, d’ailleurs, ne s’attaque jamais à deux pannes à la fois sur le même téléphone.

\asterisme

L’après-midi, le téléphone au haut-parleur qui grésille attend toujours sur le tapis. Rodrigue ne commence pas par commander un haut-parleur neuf chez le grossiste. Il lance une chanson, écoute, coupe le son, le remet, colle l’oreille à l’appareil. Le grésillement vient du bas de l’appareil, et l’écouteur, en haut, n’y est pour rien. Il démonte, trouve la petite grille bouchée par une boue de poussière et de sueur, la nettoie avec une aiguille, remonte, relance la chanson. Le son est propre. Tonton Fabrice, qui servait un client, n’a même pas tourné la tête ; il a entendu la musique changer.

Rodrigue sait par cœur, dans la boutique, une chose dont il ne parle jamais : l’emplacement des vis. Sur les modèles qu’il voit passer le plus souvent, il sait combien il y en a, de quelle longueur, laquelle se cache sous l’autocollant. Il n’y réfléchit plus. Et c’est justement parce qu’il n’y réfléchit plus que son attention reste libre pour regarder le connecteur et remarquer la croûte verte. Le par-cœur, dans la boutique, fait de la place à l’intelligence.

À l’école, c’est la même chose. On oppose souvent comprendre et apprendre par cœur, comme deux équipes qui ne se saluent pas, alors qu’elles jouent ensemble. On retient beaucoup mieux ce qu’on a compris, parce qu’une chose comprise s’accroche à d’autres et qu’on peut la reconstruire quand elle nous échappe. À l’inverse, certaines choses gagnent à être sues sans y penser, les tables de multiplication, les identités remarquables, les accords les plus fréquents : quand elles viennent toutes seules, l’esprit garde sa place pour la leçon du jour. Le par-cœur n’est donc pas l’ennemi. Le par-cœur sans compréhension, lui, est fragile. Il tient tant que l’exercice ressemble trait pour trait à celui qu’on a appris, et il casse dès qu’on déplace un chiffre.

Rodrigue en a fait l’expérience sans la reconnaître. En quatrième, il avait retenu qu’avec les fractions « on multiplie en croix ». La formule lui avait servi pour une égalité de deux fractions, elle avait marché, il l’avait gardée comme on garde un outil qui a servi une fois. Devant 2/3 + 1/4, il l’a ressortie, il a multiplié dans tous les sens, il a obtenu des nombres qui ne voulaient rien dire, et il en a conclu que c’était lui qui ne voulait rien dire. La formule était vraie, pour un autre usage, et rien, dans sa mémoire, ne lui disait lequel, puisqu’il ne l’avait jamais su.

Comment Rodrigue a-t-il appris à ouvrir un téléphone ? En regardant, d’abord. Pendant des semaines, il a suivi tonton Fabrice du premier geste au dernier, sans toucher. Puis Fabrice lui a laissé la dernière étape, revisser. Ensuite les deux dernières. Plus tard, il ne lui a plus laissé que le diagnostic, en restant à côté, et un jour il n’est même plus resté à côté. Des chercheurs en psychologie de l’éducation, autour de John Sweller et de sa théorie de la charge cognitive, ont décrit quelque chose de très proche pour les apprentissages scolaires. Pour quelqu’un qui débute, étudier un exercice résolu pas à pas est souvent plus efficace que chercher seul, à l’aveugle, pendant une heure : chercher sans outils, c’est surtout s’épuiser à tenir trop de choses dans la tête en même temps. Ces travaux suggèrent ensuite de retirer l’aide progressivement, une étape après l’autre, jusqu’à ce que l’élève fasse tout seul. Tonton Fabrice a procédé ainsi sans avoir jamais lu Sweller, et il ne le lira sans doute jamais. Les mêmes recherches indiquent d’ailleurs que l’avantage de l’exemple résolu diminue à mesure qu’on progresse : pour quelqu’un qui maîtrise déjà la marche, chercher seul redevient le meilleur exercice. L’exemple résolu sert d’appui au débutant. On le lâche quand les jambes tiennent.

Second détail : quand Rodrigue travaille, il murmure. Des bouts de phrases pour lui-même : « Ça vibre, donc il y a du courant. Ça charge, donc la batterie, c’est bon. Alors c’est entre la carte et l’écran. » Chaque geste arrive avec sa raison. Michelene Chi et ses collègues ont étudié ce geste, qu’ils appellent l’auto-explication : des élèves qui s’expliquent à eux-mêmes, à voix haute, pourquoi chaque ligne d’un exemple résolu est là, comprennent en général mieux que ceux qui se contentent de le lire, même attentivement. Or, en classe, Rodrigue recopie les corrigés en silence, ligne après ligne, sans jamais écrire le moindre « donc ». Il a le corrigé sans les raisons, comme un téléphone remonté sans qu’on sache où était la panne : il marche, jusqu’au jour où il ne marche plus.

Une idée qui circule beaucoup fait perdre du temps aux élèves qui cherchent comment apprendre. Elle veut que chacun ait son « style d’apprentissage », visuel, auditif ou kinesthésique, et qu’il faille lui présenter les leçons dans ce style pour qu’il les retienne. L’idée est séduisante ; elle flatte, et beaucoup d’élèves s’y reconnaissent. En 2008, pourtant, une revue de la recherche menée par Harold Pashler et ses collègues n’a pas trouvé de preuve solide qu’adapter l’enseignement au style déclaré d’un élève améliore ce qu’il apprend. Tu as sans doute des préférences, comme tout le monde. Elles renseignent peu sur la façon dont tu apprends le mieux : un schéma aide pour ce qui se dessine, une phrase pour ce qui se raisonne, et cela vaut pour à peu près tous les élèves.

L’autre idée est plus ancienne et plus têtue : on n’utiliserait que dix pour cent de son cerveau. C’est faux, simplement. Et cette fausse idée a un défaut de plus, celui de laisser croire qu’une réserve dort quelque part et qu’une astuce suffirait à l’ouvrir. Aucune pièce secrète n’attend sa clé. Ce qui manque, d’ordinaire, ce sont des marches, et du temps pour les poser.

\asterisme

Le même mois, à Bafoussam, un samedi matin, monsieur Kamga donne ses répétitions dans le salon d’une famille du quartier, autour de la table basse. Ils sont six, des élèves de seconde qui ne sont pas tous de sa classe. Une fille, au bout de la table, ne comprend pas pourquoi, dans 3x + 5 = 20, le cinq « passe de l’autre côté » en devenant moins cinq. Kamga explique une première fois, avec la règle. Elle hoche la tête de cette manière qu’on a de hocher la tête quand on veut que l’explication s’arrête. Il le voit. Il se lève, demande à la maîtresse de maison une poignée de haricots secs et trois boîtes d’allumettes vides, qu’il remplit en tournant le dos. Sur la table basse, à gauche, il pose les trois boîtes fermées et cinq haricots ; à droite, vingt haricots. « Chaque boîte contient autant de haricots que les autres. Les deux côtés pèsent pareil. Combien dans une boîte ? » Il retire cinq haricots à gauche. Puis, sans un mot, cinq à droite. La fille regarde les quinze haricots qui restent, les trois boîtes, et elle fait, avant lui, trois tas de cinq. « Donc passer de l’autre côté, c’est enlever des deux côtés. » Kamga rit, et la fille sursaute.

La règle était juste, et la deuxième explication valait la première ; elle entrait seulement par un autre côté. Beaucoup d’enseignants font cela tous les jours sans que personne le remarque : ils voient un regard décrocher et recommencent autrement, avec un dessin au tableau ou un exemple pris au marché. Ils ne peuvent pas le faire pour soixante-dix élèves à chaque leçon, et aucun ne le pourrait. C’est pourquoi la deuxième explication, il faut parfois aller la chercher soi-même, dans un autre manuel, chez un camarade qui a compris par un autre bout, auprès du professeur après le cours, une question précise à la main. Quand une explication ne passe pas, la conclusion raisonnable est qu’il en faut une autre.

Rodrigue, lui, dit qu’il déteste les mathématiques. Il le dit volontiers, en riant le premier comme toujours, avec l’assurance qu’il mettrait à dire qu’il n’aime pas le gombo. Une haine de ce genre a pourtant presque toujours une histoire. Elle commence par une leçon qu’on n’a pas comprise, puis une deuxième qui s’appuyait sur la première, puis une note, puis une remarque lancée devant la classe. À force, l’incompréhension se donne un caractère : elle devient un goût, une humeur, une matière « qui n’est pas pour moi ». Je ne te demanderai pas d’aimer la matière que tu détestes ; ce serait inutile, et un peu vexant pour toi. J’observe seulement, comme beaucoup d’enseignants l’observent, que l’intérêt vient souvent après un début de compétence, plus rarement avant : on s’attache à ce qu’on commence à réussir. Avant d’en réparer un, Rodrigue aimait démonter les téléphones, ce qui est une autre affaire, et c’est en apprenant à les remonter qu’il s’est mis à les aimer.

Avec une matière détestée, une ruse modeste consiste à ne pas l’attaquer par son sommet. On commence par la marche la plus basse qu’on ne maîtrise pas, celle où l’on a une chance raisonnable de réussir trois exercices de suite dès la première semaine. Ces trois exercices réussis ne prouvent rien de grand. Ils fournissent pourtant ce qui manquait depuis des années dans cette matière-là, un souvenir récent où l’on n’a pas eu honte. La haine reste dans la pièce. Elle baisse un peu la voix, et c’est assez pour travailler à côté d’elle.

Ce mois-là, il fait chaque soir trois exercices du manuel de Bella, en murmurant les « donc » comme dans la boutique. Il s’arrête au bout de vingt minutes, même quand il pourrait continuer, ce qui l’agace. En classe, le professeur ouvre le chapitre des équations, et Rodrigue recopie dans son cahier, entre deux autres lignes, la propriété écrite au tableau : « Un produit de facteurs est nul si et seulement si l’un au moins des facteurs est nul. »

\soustitre{Ce soir, essaie ceci}

Reprends la marche que tu as entourée dans ton cahier de brouillon, celle où tu n’es plus sûr. Si tu ne l’as pas encore trouvée, prends un exercice raté récemment et demande-toi ce qu’il fallait savoir pour franchir la première ligne qui t’a bloqué. La marche est presque toujours là, un ou deux étages plus bas.

Ensuite, trouve un livre d’une classe inférieure où cette marche est expliquée : le manuel d’un petit frère ou d’une cousine, le cahier d’un ancien, celui qu’un parent a gardé dans un carton. Ce que tu lis ne regarde que toi. Cherches-y un exemple résolu, un exercice déjà fait, ligne par ligne. Lis-le une fois, lentement. Puis cache-le avec une feuille et refais-le seul, dans ton cahier. Si tu bloques, découvre une seule ligne, recache, continue.

Puis explique chaque étape à voix haute, vraiment à voix haute, au mur s’il le faut, ou au petit frère qui fait semblant de ne pas écouter. Pour chaque ligne, dis pourquoi elle est là, en commençant par « parce que ». Si tu ne sais pas finir la phrase, tu viens de trouver l’endroit où il faut relire l’exemple.

À partir de demain, fais trois exercices de cette marche par jour, pendant dix jours, et jamais plus de vingt minutes. Si tu as fini en dix minutes, arrête quand même : mieux vaut garder l’envie de revenir demain que d’user la marche ce soir. Prends les exercices dans le même livre, dans l’ordre où il les donne, des plus simples aux moins simples.

Ne passe à la marche suivante que lorsque tu peux faire un exercice de celle-ci sans regarder l’exemple et sans hésiter sur l’ordre des étapes. Alors seulement, remonte : reprends l’exercice de ta classe qui t’avait bloqué au départ, et refais-le. C’est le test de la boutique. On a réparé en bas ; on vérifie en haut.

Dix jours ne réparent qu’une marche. C’est peu. S’il en manque trois sous la tienne, il te faudra un mois, peut-être davantage, et tu en sortiras un peu plus solide sur tes appuis, guère plus. Un mois, à l’échelle d’une année scolaire, cela se trouve.

\asterisme

Le samedi soir, tonton Fabrice rabat le volet de bois et compte la recette de la journée sous la lampe à pince. Dehors, la vendeuse de beignets racle le fond de sa bassine. Rodrigue a poussé sur le côté le tapis de caoutchouc et ses vis. Il a ouvert son cahier de brouillon à la page de novembre, celle où la ligne l’attend toujours, avec en bas un seul mot qu’il avait noté sans savoir pourquoi. Voilà dix soirs qu’il fait les exercices du manuel de Bella. Ce soir, il remonte.

Il pense à sa mère, à la table de la cuisine, devant les comptes de sa tontine, et à ce qu’elle avait répondu quand il lui avait demandé comment elle faisait : « Je partage. » Alors il partage. Si la chose entière avait douze parts, les deux tiers en feraient huit, le quart en ferait trois. Il écrit 8/12 + 3/12, sans se presser, en murmurant les « donc ». Il vérifie dans l’autre sens, comme on rallume un téléphone encore ouvert avant de le refermer.

Il trouve onze douzièmes. Personne n’applaudit ; tonton Fabrice recompte une liasse, et la radio de la boutique parle d’un match. Rodrigue passe à la ligne suivante.

\begin{pourcahier}{Pour le cahier}
Écris la marche que tu répares cette semaine. En dessous, recopie de ta main un exemple résolu de cette marche et, à côté de chaque ligne, un « parce que… ». Si une ligne reste sans son « parce que », entoure-la : c’est ta question pour demain.
\end{pourcahier}


%% 08-fermer-le-cahier.md
\chapitre{7}{Fermer le cahier}

La veille de l’examen, Nadège Kenfack relit ses notes pour la quatrième fois. Il est un peu plus de minuit dans sa chambre de la mini-cité, à quelques minutes à pied du campus de Ngoa-Ekellé, et janvier fait à Yaoundé ce qu’il fait chaque année : de la poussière tout le jour, une fraîcheur brève vers le matin. Sur la table : le classeur de biologie cellulaire, un paquet de biscuits entamé, une bouteille d’eau et quatre surligneurs rangés côte à côte, capuchons tournés vers elle. Le jaune pour les définitions, le rose pour ce qui « tombe », le vert pour les exemples, le bleu pour ce que le chargé de travaux dirigés a répété deux fois. Elle a mis au point ce système en octobre, quelques semaines après être arrivée de Bertoua, et elle en est assez fière. Il a la rigueur d’un inventaire de quincaillerie, et elle sait de quoi elle parle : elle a grandi derrière un comptoir où l’on compte les boulons.

À la quatrième lecture, pourtant, le système s’est effondré sous son propre poids. Presque chaque ligne a reçu sa couleur, certaines en ont reçu deux, et les pages ressemblent maintenant à l’étal d’une vendeuse de pagnes au marché : plus de fond, rien que des motifs. Le bout de ses doigts est taché de jaune. Cela ne l’inquiète pas, parce que tout va bien. La mitose, la membrane plasmique et sa double couche de lipides, la mitochondrie qui respire pour la cellule : chaque paragraphe lui paraît familier, presque évident, comme une rue qu’on prend tous les jours. Elle tourne les pages plus vite qu’à la première lecture, et elle prend cette vitesse pour un bon signe. Vers 1 h, elle éteint l’ampoule. Au bout du couloir, quelqu’un tire encore de l’eau au robinet commun.

Le lendemain matin, dans l’amphithéâtre, le sujet arrive par rangées, de main en main, et chacun garde la tête baissée sur sa table vide jusqu’à ce que la sienne arrive. Il fait encore frais. Les bancs de bois sont froids sous les avant-bras, et quelque part au fond quelqu’un fait cliquer son stylo à intervalles réguliers, comme une horloge qu’on n’aurait pas réglée. Première question : décrire les étapes de la mitose et ce qui se passe au cours de chacune. Nadège sait où se trouve la réponse. Page onze, en haut à droite, dans un encadré jaune, avec une flèche verte qui descend vers un schéma. Elle voit l’encadré. Elle voit la flèche. Elle ne voit pas ce qu’il y a dedans. Pendant de longues minutes, elle reste devant sa copie avec une seule certitude, absurde et exacte : elle se souvient de la couleur.

Elle finit par écrire quelque chose, des morceaux. Le mot « chromosomes », le mot « fuseau », un nom de phase dont elle n’est plus sûre de la place, une phrase qui commence bien et s’arrête au milieu, faute de suite. Les questions suivantes se passent un peu mieux, d’autres plus mal. En sortant, sur les marches, elle entend deux garçons comparer leurs réponses et s’éloigne avant de savoir lequel avait raison. La veille, elle reconnaissait tout. Ce matin, elle n’a presque rien su retrouver. Quand les résultats de la session tombent, quelques semaines plus tard, ils ne lui apprennent rien que la première question ne lui ait déjà dit.

\asterisme

Nadège a travaillé. Elle a lu quatre fois un cours épais, tard, avec ordre, et elle a fait ce que fait presque tout le monde. Dans beaucoup de maisons, « réviser sa leçon » veut dire la relire, et un enfant qui relit a l’air de travailler, ce qui est vrai. Relire produit d’ailleurs une sensation bien réelle : la familiarité. Les phrases glissent, les schémas sont à leur place, on devine la fin du paragraphe avant de l’avoir lue. L’ennui, c’est que cette sensation ressemble trait pour trait à celle de savoir : deux cousines qu’on confond de loin.

Pense à quelqu’un que tu croises au marché, un samedi, entre les tas de tomates et les sacs de riz. Son visage te dit quelque chose ; tu es sûr de l’avoir connu, à l’école primaire peut-être. Tu le reconnais sans hésiter. Mais si l’on te demandait son nom, là, tout de suite, tu chercherais, et peut-être que rien ne viendrait. Reconnaître, c’est dire « je l’ai déjà vu » quand la chose est devant toi. Se rappeler, c’est la faire revenir quand elle n’est pas là. La relecture entraîne la première de ces deux opérations ; l’examen demande presque toujours la seconde. Nadège a passé quatre soirées à s’exercer pour une épreuve que personne ne lui ferait passer.

Les psychologues parlent parfois d’illusion de maîtrise. Le mot est sévère pour une impression qui, au départ, est vraie, et c’est justement pour cela qu’on s’y laisse prendre. Le même mécanisme joue avec un corrigé lu trop tôt. Tu ouvres la solution d’un exercice que tu n’as pas cherché, chaque étape s’enchaîne, tout paraît logique, presque facile, et tu refermes le livre avec le sentiment d’avoir compris un exercice que tu serais incapable de refaire seul. Une solution lue avant d’avoir cherché fabrique de la familiarité, et rien d’autre tant qu’on n’a pas fermé le livre pour refaire l’exercice.

En 2006, deux chercheurs américains, Henry Roediger et Jeffrey Karpicke, ont mené une expérience devenue classique. Ils ont donné à des étudiants de courts textes à apprendre. Les uns les relisaient plusieurs fois. Les autres, après une première lecture, devaient écrire sur une feuille, sans regarder, tout ce dont ils se souvenaient. Cinq minutes après, ceux qui avaient relu s’en sortaient un peu mieux. Une semaine plus tard, l’ordre s’était inversé : ceux qui s’étaient testés retenaient nettement davantage. Et, interrogés sur leurs chances, ceux qui avaient relu se croyaient mieux préparés que les autres.

Ils se sentaient plus sûrs, et ils savaient moins. Depuis, de nombreuses études vont dans le même sens, avec des élèves de tous âges et dans des matières très diverses. Conséquence déroutante : se tester sert à vérifier ce qu’on sait, comme on compte la monnaie avant de quitter le comptoir, et aussi, peut-être surtout, à apprendre. Chaque fois que tu vas chercher une réponse dans ta tête, sans la page sous les yeux, tu rends ce trajet plus facile pour la fois suivante. Chaque fois que tu relis, tu laisses la page faire le trajet à ta place.

En 2013, une équipe conduite par John Dunlosky a passé en revue ce que la recherche disait de dix techniques d’étude parmi les plus répandues. Deux d’entre elles se détachaient par leur utilité : se tester, et étaler ses révisions dans le temps. Relire et surligner figuraient parmi les moins utiles. Si tu as quatre surligneurs dans ta trousse, rassure-toi : surligner n’abîme rien. Le risque est qu’on s’arrête là, avec le sentiment du devoir accompli, parce qu’une page colorée ressemble à une page travaillée. Garde tes surligneurs si tu les aimes. Sers-t’en pour marquer ce sur quoi tu vas t’interroger ensuite.

\asterisme

Le dimanche qui suit les résultats, Nadège appelle sa mère, comme tous les dimanches. À Bertoua, la quincaillerie est ouverte ; derrière la voix maternelle, on entend un client qui réclame des pointes et le frottement d’un sac de ciment qu’on traîne sur le sol. Nadège a préparé ses phrases dans le taxi, en revenant du campus ; elle annonce la nouvelle d’une traite, pour ne pas avoir à la répéter. Sa mère se tait un moment, puis demande, sans reproche : « Tu n’avais pas bien lu ? » Nadège répond qu’elle a lu quatre fois. « C’est peut-être ça, le problème, maman. » Sa mère croit à une insolence et le lui dit. Nadège rit pour la première fois depuis une semaine, sans pouvoir expliquer pourquoi.

Le soir, sans trop y croire, elle fait une chose qu’elle n’a jamais faite. Elle prend une feuille blanche, pose le classeur fermé à l’autre bout de la table, et essaie d’écrire la mitose. Elle connaît maintenant ce chapitre mieux qu’aucun autre : elle l’a relu après l’examen, par rage, et deux fois encore depuis. Sur la feuille, il lui faut plusieurs minutes pour aligner les phases, et quand elle ouvre enfin le classeur pour vérifier, elle découvre qu’elle en a inversé deux et oublié ce que deviennent les chromosomes à la fin. Elle corrige au surligneur vert, faute d’autre stylo sous la main. Pour la première fois depuis octobre, une couleur sert à marquer ce qu’elle ne savait pas.

Ce qu’elle vient de vivre sur cette feuille, tout le monde le vit le jour où il essaie. La page blanche ne laisse rien passer ; elle n’a pas la politesse du classeur ouvert, qui te souffle la suite avant que tu l’aies cherchée. Mais elle a sur l’examen un avantage décisif : elle arrive avant lui, et elle ne met pas de note.

La feuille a montré à Nadège une chose plus inquiétante : même relu et relu, un chapitre s’en va. On peut le prendre mal. On peut aussi y reconnaître l’une des plus anciennes observations de la psychologie. À la fin du dix-neuvième siècle, un chercheur allemand, Hermann Ebbinghaus, a mené sur lui-même une expérience d’une patience presque comique. Il apprenait par cœur des listes de syllabes sans aucun sens, des assemblages du genre de « dax » ou « bok », pour être sûr que rien de ce qu’il savait déjà ne viendrait l’aider. Puis il mesurait le temps qu’il lui fallait pour les réapprendre, après quelques minutes, après quelques heures, après plusieurs jours. Ce qu’il a publié en 1885 a été retrouvé bien des fois depuis : on oublie très vite juste après avoir appris, puis de plus en plus lentement. Et une chose qu’on a sue se réapprend plus vite que la première fois, même quand on croit l’avoir entièrement perdue.

Oublier est le fonctionnement ordinaire de toute mémoire humaine, la tienne comme celle du premier de ta classe, qui oublie lui aussi et revient simplement plus souvent. On ne peut pas empêcher l’oubli. On peut le ralentir, et la manière la mieux établie de le faire tient en un verbe : revenir. Plutôt que de travailler une leçon pendant quatre heures la veille, la reprendre le lendemain, puis quelques jours plus tard, puis la semaine suivante, en laissant chaque fois passer davantage de temps. Chaque retour arrive au moment où la notion commençait à s’effacer, et la fixe mieux que le précédent. Une vaste synthèse de travaux, publiée en 2006, a confirmé très largement l’avantage de cet espacement sur le travail concentré en une seule séance.

Le bachotage de la dernière nuit laisse bien des traces pour le lendemain matin. Puis elles s’évaporent, souvent avant la fin du mois, et il faut tout recommencer pour l’examen suivant, comme on remplit un seau percé. Dans un emploi du temps, quatre heures en une nuit et quatre heures réparties sur trois semaines pèsent le même poids ; dans la mémoire, les secondes restent nettement plus longtemps.

\asterisme

Si se tester et espacer donnent de si bons résultats, pourquoi si peu d’élèves le font-ils ? La réponse tient en partie à une sensation : la bonne méthode paraît moins bonne que la mauvaise. Relire est doux. On avance, on reconnaît, on tourne les pages, on se couche avec le sentiment d’avoir bien employé sa soirée. Se tester est vexant : dix minutes devant une feuille blanche, et ce qui revient fait pitié. Revenir sur une leçon au bout d’une semaine est pire encore, puisqu’on y constate surtout ce qu’on a perdu. Tout, dans l’expérience immédiate, conseille de relire. Nadège l’a senti dès le deuxième soir. Après dix minutes de feuille blanche sur la membrane, elle a écrit en bas de la page, de son écriture la plus nette : « Méthode recommandée pour se sentir bête. » Puis elle a quand même corrigé ce qui manquait, parce que la feuille était déjà gâchée et qu’elle déteste gâcher.

Le chercheur américain Robert Bjork a donné à ce paradoxe un nom un peu étrange : les difficultés désirables. Ce qui coûte un effort au moment d’apprendre, comme chercher sans regarder ou attendre avant de revoir, aide souvent à retenir ; ce qui coule tout seul s’efface à peu près aussi vite qu’il est venu. L’expression a l’air d’une provocation et demande une lecture attentive, car toute difficulté n’est pas bonne à prendre : un polycopié illisible ou une chambre où deux radios se disputent le couloir compliquent le travail sans rien lui apporter. Les difficultés dont parle Bjork sont précises et voulues : l’effort de faire revenir ce qu’on a appris, le temps qu’on laisse passer exprès.

Cette idée console un peu, à condition de ne pas la prendre pour une promesse. Une séance de questions dont tu sors avec l’impression d’avoir été mauvais a peut-être été une bonne séance. La vexation est un renseignement : elle dit où sont les trous pendant qu’il est encore temps de les boucher. Le jour de l’examen, les mêmes trous seront là, et il sera trop tard pour les découvrir.

Une autre difficulté de ce genre consiste à mélanger. Quand tu fais vingt exercices sur le même modèle, le quinzième ressemble tellement au quatorzième que tu finis par appliquer la méthode sans te demander pourquoi c’est celle-là. Or, à l’examen, personne n’écrit en haut du sujet quelle méthode employer : la moitié du travail consiste à reconnaître de quel problème il s’agit. Des chercheurs, en mathématiques surtout, ont observé que des élèves qui mélangeaient plusieurs types d’exercices dans une même séance, de sorte qu’il faille chaque fois identifier le type avant de le traiter, réussissaient souvent mieux par la suite, alors même que la séance leur avait paru plus pénible. Les résultats sont prometteurs, et restent à confirmer. L’idée, en tout cas, ne coûte rien à essayer : au lieu de faire toute la série du chapitre quatre, puis toute celle du chapitre cinq, prends deux exercices de chacun et mélange-les.

Nadège l’essaie à sa façon, quelques jours plus tard. Elle écrit sur des bouts de papier des questions sur trois chapitres différents, la mitose, la membrane, les enzymes, les plie, les jette dans une boîte à chaussures vide et les tire au hasard. La première fois, elle répond à une question sur la membrane en parlant de la mitose, parce que c’est le chapitre qu’elle vient de reprendre, et elle s’en aperçoit au milieu de sa phrase. Elle note l’erreur sans commentaire. L’examen aurait trouvé cette confusion pour elle, sans prévenir.

\asterisme

Le plus difficile, peut-être, est le temps long, parce qu’il n’a rien de spectaculaire. Nadège a devant elle quelques semaines avant la session de rattrapage. Elle les compte comme elle compte son argent, sans illusion et sans arrondir. Sur une feuille, elle part de la date de l’examen et remonte vers aujourd’hui. Elle place chaque chapitre une première fois, puis prévoit pour chacun plusieurs retours, de plus en plus espacés. La feuille n’a rien d’impressionnant : des cases, des flèches, des petites croix. Elle la fixe au mur au-dessus de la table, à côté d’un calendrier offert par une pharmacie.

Le premier retour tombe un mardi, entre deux séances de travaux dirigés. Nadège s’assoit sur un muret derrière le bâtiment, à l’ombre d’un flamboyant dont les feuilles sont grises de poussière, déplie la feuille de questions qu’elle garde dans sa poche et cache les réponses avec la main. Sur cinq questions, elle en retrouve deux en entier, une à moitié, et deux pas du tout. Elle trouve le résultat humiliant et le note quand même, avec la date. Le mardi suivant, les deux questions perdues reviennent, mais une autre, qu’elle croyait acquise, s’est éloignée. La progression ressemble aux comptes de la quincaillerie pendant la saison des pluies, quand les clients paient en retard : des entrées, des sorties, et un solde qui, certaines semaines, finit quand même par monter. Nadège, qui a appris à lire un cahier de comptes avant de savoir faire une division, ne demande pas davantage à ses feuilles de questions.

Un calendrier à rebours rend visibles deux choses qu’on préfère ne pas voir : le nombre de jours qui restent, et le fait que chaque leçon a besoin de plusieurs passages. Il oblige aussi à accepter une idée peu glorieuse, celle du petit peu. Vingt minutes de questions sur la leçon de la semaine dernière, un mardi soir, ne donnent l’impression de rien et ne font pas une bonne histoire à raconter. Répétées pendant deux mois, elles font souvent davantage que les grandes séances héroïques de la dernière semaine, qui, elles, font d’excellentes histoires.

\soustitre{Ce soir, essaie ceci}

Après avoir étudié une leçon, n’importe laquelle, ferme le cahier. Pose-le loin de toi, à l’autre bout de la table ou sur le lit, assez loin pour que tes yeux ne puissent pas tricher. Prends une page de ton cahier de brouillon et écris tout ce dont tu te souviens, sans regarder, pendant dix minutes. Les définitions, les étapes, les exemples, le schéma si tu peux le refaire, même de travers. Écris aussi ce dont tu te souviens à moitié, en le signalant d’un point d’interrogation dans la marge. Les dix minutes vont te paraître longues. Ne t’arrête pas à la cinquième sous prétexte que plus rien ne vient : c’est souvent dans les dernières minutes que quelque chose remonte.

Ensuite seulement, rouvre le cahier. Corrige ta page avec un stylo d’une autre couleur : ajoute ce qui manque, rectifie ce qui est faux. Ne te juge pas à la quantité de couleur. Ces corrections dessinent la carte de ce que tu dois retravailler, une carte que la relecture ne t’aurait jamais donnée.

Puis écris trois questions sur cette leçon. Choisis des questions qui t’obligent à produire quelque chose, une explication, une étape, un schéma, plutôt qu’à répondre par oui ou par non. « Pourquoi la deuxième étape vient-elle avant la troisième ? » vaut mieux que « Est-ce que je connais les étapes ? ». Écris les réponses plus loin dans le cahier, là où tu ne les verras pas par accident. À côté de chaque question, note des dates : demain, puis dans trois ou quatre jours, puis la semaine prochaine. À chacune de ces dates, tu répondras sans tes notes, et seulement après tu vérifieras.

Les intervalles exacts importent moins que le fait de revenir. Si tu manques un rendez-vous, ne recommence pas tout : réponds le lendemain et continue. Si, demain, les trois questions te paraissent faciles, écris-en une quatrième, plus difficile. Et si elles te paraissent impossibles, garde-les telles quelles : ce sont les bonnes.

Pour les choses qu’il faut savoir par cœur, définitions, formules, dates, vocabulaire, il existe une manière d’organiser ces retours sans calendrier, sans courant, sans connexion et sans argent. On l’appelle la boîte de Leitner, du nom de l’homme qui l’a popularisée dans les années 1970. Prends un carton d’emballage, celui d’une boîte de lait en poudre ou de chaussures. Avec deux morceaux du même carton, découpe deux cloisons et glisse-les à l’intérieur : tu as trois compartiments. Dans le reste, découpe des rectangles de la taille de ta paume. Sur une face, écris une question ; sur l’autre, la réponse. Toutes les cartes neuves vont dans le premier compartiment.

Les cartes du premier compartiment se travaillent chaque jour, celles du deuxième tous les deux ou trois jours, celles du troisième une fois par semaine. Tu tires une carte, tu lis la question, tu réponds à voix haute ou dans ton cahier, puis tu retournes la carte. Si tu as trouvé, elle avance d’un compartiment. Si tu t’es trompé, elle retourne dans le premier, d’où qu’elle vienne. Ce qui est su s’éloigne peu à peu et revient de moins en moins souvent ; ce qui résiste reste devant toi jusqu’à ce qu’il cède.

La boîte fait d’elle-même le tri que tu aurais eu du mal à faire seul. Elle a ses limites : une démonstration ou un commentaire de texte ne tiennent pas sur une carte. Dans ce cas, la carte porte seulement la consigne, « Refais la démonstration vue mardi », et la réponse se fait dans le cahier, puis se vérifie dans le cours. Le carton, lui, ne tombe jamais en panne de batterie.

\asterisme

Le soir où elle a fini son calendrier, Nadège ferme son classeur. Elle le pose sur l’étagère, au-dessus du seau d’eau, hors de portée de la main. Dans le couloir de la mini-cité, une radio passe d’une chanson à un journal, des sandales traînent jusqu’au bout du couloir et reviennent. Sous le lit traîne le carton vide d’une boîte de lait en poudre qu’elle gardait sans savoir pourquoi, par habitude de commerçante qui ne jette rien. Elle le découpe aux ciseaux en rectangles inégaux, en garde deux bandes pour les cloisons, et sur le premier rectangle elle écrit : « Les étapes de la mitose, dans l’ordre, et ce qui arrive aux chromosomes à chacune. » Puis deux autres questions, sur la membrane. Les surligneurs restent dans la trousse.

Elle ne sait pas encore si ça marchera. Elle sait que relire ne marchait pas.

\begin{pourcahier}{Pour le cahier}
Écris aujourd’hui trois questions sur ta leçon du jour. Note à côté de chacune la date où tu y répondras sans regarder. Le jour venu, coche celles que tu as retrouvées, et recopie les autres un peu plus bas, avec une nouvelle date.
\end{pourcahier}


%% 09-vingt-et-une-heures.md
\chapitresn{Interlude — Vingt et une heures}

Il est 21 h. Tu es au milieu d’une leçon, à la moitié d’une page, sur une phrase que tu viens enfin de comprendre. Le courant part.

Il part d’un coup, comme toujours, sans prévenir personne. L’ampoule s’éteint, la télévision des voisins se tait au milieu d’un mot, et le frigo cesse de ronronner : tu t’aperçois seulement maintenant qu’il ronronnait. Le quartier pousse ce soupir que tu connais, le même dans toutes les maisons à la fois, moitié agacé, moitié habitué. Quelqu’un dit « Encore ! » dans la cour d’à côté. Un enfant rit, parce que pour lui le noir est encore un jeu.

Deux maisons plus loin, un groupe électrogène tousse, hésite, puis démarre. Pas chez toi. Son bruit remplit la rue, et là-bas une fenêtre s’allume en jaune, seule dans l’alignement des toits.

Tu ne vois plus ta page. Tu la sens sous ta main, le grain du papier, le léger relief que le stylo a creusé en appuyant. Les mots, eux, sont partis avec la lumière.

Tu pourrais chercher la bougie. Elle est quelque part dans le tiroir de la cuisine, avec les allumettes, si personne ne les a prises. Tu pourrais aussi allumer la lampe torche de ton téléphone et continuer comme si de rien n’était. Tu ne cherches pas tout de suite.

Tu fermes les yeux, puisqu’ils ne servent plus. Et tu récites.

Ce que tu as appris ce soir. Le titre de la leçon, d’abord : il vient tout seul. La première définition, presque mot pour mot, avec la virgule au bon endroit. Le schéma du bas de la page, tu le vois encore, avec ses flèches ; mais au bout d’une des flèches, il y avait un mot, et ce mot te manque. Tu cherches. Le groupe électrogène ronfle. Un moustique passe près de ton oreille et repart. Le mot ne vient pas. Tu le laisses où il est.

Puis les trois questions écrites hier. La première, tu y réponds à voix basse, en entier, et ta réponse a la forme exacte de la phrase du cahier. La deuxième, tu la commences bien, tu hésites au milieu, tu retrouves la fin par un autre côté. La troisième reste au bord. Tu sais qu’il y avait plusieurs étapes. Tu n’en retrouves que deux.

À Ngaoundéré, au même moment, le courant est parti aussi. Ibrahim Saïdou a entendu le claquement juste avant : c’est le transformateur du quartier, il le sait, et personne n’y touchera avant demain. Il est chez son oncle ce soir ; dans le coin du salon, le ventilateur sur pied qu’il a remis en marche en novembre ralentit, tourne encore un peu sur sa lancée, puis s’arrête, et la chaleur redescend sur la pièce. Pour s’occuper, il récite à mi-voix la loi d’Ohm : la tension aux bornes d’un conducteur ohmique est égale au produit de sa résistance par l’intensité du courant qui le traverse. Il sourit dans le noir, parce qu’il ne reste plus un seul ampère dans le quartier pour lui donner raison.

Toi aussi, tu es dans la chaleur, ou dans le bruit des moteurs, ou dans l’odeur de pétrole d’une lampe qu’on allume ailleurs. Tu ne sais pas quand le courant reviendra. Les voisins non plus.

Mais tu sais une chose que tu ne savais pas il y a dix minutes. Tu sais ce qui est à toi : ce qui est resté quand la page a disparu, et qui tiendrait encore si on te réveillait demain matin pour te le demander. La page, les couleurs et le soulignement ne sont plus là pour te faire croire que tu sais. Ne compte que ce qui revient. Le reste, le mot au bout de la flèche, la troisième étape, tu sais maintenant où le chercher, et ce sera la première chose que tu regarderas quand tu rouvriras le cahier.

Ce soir, si le courant part, ne cherche pas la bougie tout de suite.

Le groupe électrogène des voisins continue de tourner. Ta page est toujours là, sous ta main, illisible et entière. Le courant ne revient pas.

Tu continues.


%% 10-dormir-n-est-pas-tricher.md
\chapitre{8}{Dormir n’est pas tricher}

À Garoua, fin février, la nuit rafraîchit à peine les chambres ; elle se contente d’y éteindre le jour. La tôle a bu le soleil toute la journée et elle le rend par en haut, lentement, comme un fer à repasser qu’on a débranché mais qu’on n’ose pas encore ranger. Il est 2 h du matin. Aïcha Bello, Terminale D, est assise par terre, le dos contre le mur le moins chaud de la pièce, le classeur de SVT ouvert sur les genoux. Le bac blanc commence à 7 h 30. Il lui reste deux chapitres à revoir, et elle a décidé, vers minuit, qu’elle les reverrait tous les deux.

Le courant est parti avant 21 h, comme souvent cette semaine. La lampe solaire est posée sur un seau retourné. Elle a chargé toute la journée sur le toit, calée entre deux briques pour que le vent ne l’emporte pas, et depuis le début de la soirée elle a tenu son rôle. Mais à cette heure-ci, sa lumière a changé. Elle est passée du blanc au jaune, puis à une sorte de beige fatigué qui n’accroche plus le papier. Aïcha approche la page de ses yeux. Le schéma de la synapse, c’est elle qui l’a dessiné, à l’encre noire, avec des flèches plus nettes que celles du manuel ; elle le regarde maintenant comme on regarde le plan d’un quartier où l’on n’a jamais mis les pieds.

Près de la porte, sur la natte, Hadja dort. Elle a douze ans et un talent que personne ne lui a appris : trouver, dans une maison qui étouffe, l’endroit exact où passe l’air. Ici, c’est le bas de la porte, là où le bois ne touche pas tout à fait le seuil. Elle dort sur le dos, un bras en travers des yeux, le pagne repoussé jusqu’aux genoux. De temps en temps, une moto passe sur la route, et la poussière fine de l’harmattan, qui est partout en cette saison, se dépose un peu plus sur le classeur.

À côté du seau, un verre de café soluble. Aïcha y a mis trois cuillères de sucre. Elle a expliqué un jour à sa mère que le café sans sucre était une punition, et qu’elle avait déjà le Bac pour ça. Le café est tiède, il colle aux dents, il ne fait plus grand-chose. Elle le boit quand même, par principe.

Elle sait qu’elle devrait dormir. Elle le sait comme on sait qu’il faudrait boire plus d’eau ou ranger ses affaires : une vérité générale, qui ne la concerne pas ce soir. Car dormir, ce soir, lui ferait l’effet d’une tricherie. Se coucher quand il reste un chapitre, c’est quitter le terrain avant la fin du match. Dans sa tête, le travail se mesure en heures prises à la nuit, et plus la nuit est courte, plus le travail est honnête. Au mur, le calendrier où elle barre les jours un à un attend le trait d’aujourd’hui. Elle ne l’a pas encore tracé. La journée n’est pas finie, puisqu’elle ne s’est pas couchée.

Vers 3 h, elle s’allonge « cinq minutes » sur le côté, le classeur contre la poitrine. La lampe s’éteint un peu plus tard, sans que personne la voie s’éteindre.

Le lendemain matin, dans la salle du bac blanc, l’épreuve de SVT s’ouvre sur une question de trois lignes. Aïcha la lit. Arrivée au bout, elle a perdu le début. Elle la relit en suivant les mots de la pointe de son stylo, comme on suit des traces dans le sable. Le début tient, cette fois, mais c’est la fin qui s’échappe : quelque chose sur le message nerveux, sur une expérience, sur ce qu’il faut « en déduire ». Elle la lit une troisième fois. Elle connaît chaque mot. Elle a dessiné cette synapse plus souvent qu’elle ne saurait le dire. Mais les mots refusent de rester ensemble assez longtemps pour devenir une question. Autour d’elle, les stylos grattent déjà. Elle a chaud, elle a sommeil, et ce qu’elle sait le mieux au monde se trouve quelque part derrière une vitre sale. Elle finit par répondre à la deuxième question, puis à la quatrième, et revient à la première à dix minutes de la fin, pour y écrire trois lignes dont elle n’est pas fière.

À midi, devant le portail, pendant que les autres comparent leurs réponses, la réponse lui revient d’un coup, entière, avec l’expérience, la déduction et même le schéma qu’il aurait fallu faire. Elle la connaissait. Elle l’a retrouvée comme on retrouve ses clés dans sa poche une fois rentré à la maison, après avoir fouillé partout ailleurs.

\asterisme

Du nombre d’heures qu’il faut dormir, je ne donnerai pas de chiffre. Je me méfie des chiffres ronds qu’on brandit comme des règlements : le besoin de sommeil varie d’une personne à l’autre et d’un âge à l’autre, et la vie d’une élève de Terminale à Garoua ne se règle pas sur une horloge de laboratoire. Ce qu’on sait tient en quelques points, chacun avec son degré de certitude, et tu en feras ce que tu voudras.

On a longtemps imaginé le sommeil comme une pause : le corps s’arrête, la tête s’éteint comme la lampe d’Aïcha, et l’on reprend le lendemain là où l’on s’était arrêté. Les chercheurs qui étudient le sommeil décrivent des nuits plus occupées. Les travaux de Jan Born et de ses collègues, en Allemagne, et ceux de beaucoup d’autres équipes, suggèrent qu’une partie de ce qu’on a appris dans la journée est reprise pendant la nuit et en partie stabilisée, un peu comme si la mémoire repassait certains souvenirs récents pour les fixer. Personne ne sait encore bien comment cela se fait, et il faut se méfier des affirmations spectaculaires qui circulent sur le sujet. Mais l’idée générale tient, et de nombreuses études convergent : le sommeil fait partie de l’apprentissage. Il en est même, peut-être, la dernière étape, celle qu’on ne voit pas.

Attention : on n’apprend pas en dormant. Ce que tu n’as jamais étudié ne se glissera pas dans ta tête pendant la nuit, et le classeur posé sous l’oreiller n’a jamais donné à personne autre chose qu’un torticolis. La nuit travaille sur ce que tu as étudié dans la journée, et ce travail-là a besoin d’elle pour tenir.

On sait ensuite, avec plus d’assurance encore, ce que fait le manque de sommeil. Après une nuit trop courte, l’attention flanche : elle tient quelques minutes, puis elle glisse, revient, glisse encore. L’humeur se gâte ; on s’irrite plus vite, on se décourage plus tôt. Et le raisonnement, celui qui doit garder en tête le début d’une question pendant qu’on en lit la fin, devient lent et troué. Aïcha l’a vécu devant ses trois lignes : la synapse était toujours là, quelque part, et il lui manquait de quoi tenir une phrase entière.

Voilà pourquoi la nuit blanche avant un devoir est presque toujours une mauvaise affaire. On gagne quelques heures de lecture, faites dans un état où l’on retient mal, et on les paie le lendemain avec l’outil même dont on a besoin pour répondre. « Presque toujours », parce que la vie ne suit pas les manuels et qu’il y a des nuits où l’on n’a pas le choix : un petit frère malade, ou un travail du soir, comme celui de Nadège à Yaoundé. Mais quand on a le choix, et qu’on choisit de veiller par principe, on se trompe de principe. Le café peut repousser la somnolence un moment. Il ne rend pas ce que la nuit aurait fait.

Si Aïcha avait noté, cette nuit-là, ce qu’elle avait réellement fait entre minuit et 3 h, elle aurait trouvé moins qu’elle ne croit. Une heure correcte, peut-être, sur le premier chapitre. Ensuite, des pages tournées, des lignes relues sans les voir, un schéma recopié par réflexe, un quart d’heure sur son téléphone, et de longs moments à fixer la lampe en se demandant si elle tiendrait. La fatigue a ceci de traître qu’elle continue de te donner le sentiment de travailler bien après que le travail s’est arrêté.

L’idée d’Aïcha, pourtant, mérite le respect : c’est une idée de travailleuse. Elle a vu son père rentrer de ses gardes au centre de santé avec les yeux rouges, sa mère coudre tard quand une commande pressait. Elle en a conclu que les choses sérieuses coûtent, et elle a raison. Elle se trompe seulement de monnaie. Le sommeil n’est pas le prix qu’on refuse de payer pour réussir ; c’est une partie de ce qu’on achète.

\asterisme

Le bac blanc fini, Aïcha rentre à pied sous le soleil de 14 h. Elle mange à peine, s’allonge sur la natte de Hadja, près de la porte, et dort trois heures d’un sommeil lourd et moite dont elle sort plus fatiguée qu’avant, la joue marquée par les fibres de la natte. Hadja, rentrée de l’école, l’enjambe pour poser son sac et lui demande comment ça s’est passé. Aïcha répond sans ouvrir les yeux : « J’ai lu la première question trois fois. Je crois qu’elle aussi m’a lue, et qu’elle n’a pas été impressionnée. » Hadja rit, sans bien savoir de quoi. Le soir, évidemment, Aïcha n’a plus sommeil. À minuit, elle a les yeux ouverts sur le plafond et écoute la tôle craquer en refroidissant. Elle en conclut qu’elle est « faite pour la nuit ». Elle a en partie raison, pour une raison qu’elle ne soupçonne pas.

Chez les adolescents, l’horloge intérieure qui règle le sommeil tend à se décaler. Les travaux de la chercheuse américaine Mary Carskadon, et de ceux qui l’ont suivie, ont décrit ce glissement : vers la puberté, l’envie de dormir arrive plus tard le soir qu’à dix ans, et le réveil du matin tombe alors au milieu de ce qui serait encore, pour le corps, un morceau de nuit. Le glissement a été observé sous bien des latitudes. Rien à voir avec la paresse, donc, et ta mère, quand elle t’appelle à 5 h 30 d’une voix de plus en plus forte, ne corrigera pas ce décalage en haussant le ton. Elle a d’autres raisons de t’appeler, et elles sont souvent bonnes.

L’école ne s’adaptera pas à ton horloge ; il serait malhonnête de te promettre le contraire. Les cours commencent à l’heure où ils commencent, et l’aîné ou l’aînée d’une famille a souvent de l’eau à tirer ou un petit frère à préparer avant de partir. Ce que tu peux faire de cette information est plus modeste. Tu peux cesser de te traiter de paresseux à 6 h du matin, ce qui ne réveille personne. Et tu peux éviter de bousculer l’horloge plus qu’elle ne l’est déjà : on conseille souvent, si l’on dort le jour, une sieste courte plutôt que trois heures en plein après-midi, et une heure de coucher à peu près fixe d’un soir à l’autre, samedi compris. Une horloge qu’on dérange tous les jours finit par ne plus savoir à quelle heure sonner.

Vient ensuite la chaleur, et ce qui l’accompagne quand on vit nombreux dans peu de pièces : la lumière des autres, la radio du voisin qui ne se tait qu’après les informations. Les conseils qu’on lit d’habitude sur le sommeil supposent une chambre à soi et du silence. Ceux qui suivent s’en passent. Hadja a déjà trouvé le premier toute seule : dormir là où l’air passe, même si ce n’est pas sa place habituelle, et ouvrir dès que le dehors est plus frais que le dedans. Le deuxième tient dans un verre : la chaleur dessèche sans prévenir, et le mal de tête de 16 h est parfois simplement de la soif. Boire de l’eau au fil de la journée ne fait de mal à personne.

Le troisième demande de réorganiser un peu la journée. À Garoua, l’heure la plus fraîche est celle qui précède le lever du soleil, quand la maison dort encore et que la tôle a enfin rendu toute sa chaleur. C’est l’heure des choses difficiles : l’exercice qui résiste, le chapitre qu’on ne comprend pas, le schéma à refaire sans modèle. Les heures lourdes de l’après-midi, elles, conviennent aux tâches qui demandent moins de tête : recopier une leçon au propre, ranger ses fiches. Pour le bruit, on n’a que des arrangements. Demander à la maison une heure calme, toujours la même, et la demander pour de bon, en expliquant pourquoi, plutôt que de soupirer chaque soir en espérant qu’on comprenne. Aller travailler dans la cour à l’aube, ou dans la salle de classe avant l’arrivée des autres. Et quand il n’y a qu’une lampe pour deux, s’entendre sur les heures plutôt que se la disputer à chaque coucher de soleil : chez les Bello, il arrive qu’Aïcha gagne parce qu’elle a le Bac, et que Hadja gagne parce qu’elle a douze ans et une dictée pour le lendemain.

Le corps, enfin, entre en classe avec la tête. Le matin du bac blanc, Aïcha n’avait rien avalé, sauf le fond de café froid resté dans le verre. De menu, tu n’en trouveras pas ici, et tu mangeras ce qu’il y a à la maison. Mais un ventre vide à 10 h du matin pense mal, et manger quelque chose avant de partir, même peu, est un geste de travail. Quant au mouvement, beaucoup d’élèves constatent qu’une marche ou un match dans la cour les laisse mieux disposés à s’asseoir ensuite. Des études suggèrent que l’activité physique a des effets favorables, mais modestes, sur l’humeur et l’attention. Modestes : je tiens au mot. Méfie-toi de tout ce qui se vend en flacon avec la promesse d’une mémoire neuve. Le seul produit que ce chapitre te recommande est gratuit, et tu t’en sers déjà mal tous les soirs.

\asterisme

Revenons à la nuit d’avant le bac blanc, un peu après 1 h. Le téléphone d’Aïcha, posé contre le seau, s’est allumé tout seul. Quelqu’un, dans le groupe de la classe, venait d’envoyer « les sujets probables » de l’épreuve de SVT. Une quarantaine de messages ont suivi : des rires, un message vocal qu’elle a écouté deux fois, une dispute pour savoir si c’était sérieux, puis un camarade qui jurait que son cousin connaissait quelqu’un. Aïcha a tout lu. Cela a pris un quart d’heure. Quand elle est revenue à sa page, le schéma de la synapse avait l’air d’avoir été dessiné par une autre. Les « sujets probables », bien sûr, ne sont pas tombés.

Ce genre de quart d’heure a été bien étudié, et personne n’y voit une faute morale. Les psychologues qui travaillent sur ce qu’on appelle la commutation de tâches ont observé depuis longtemps que passer d’une activité à une autre a un coût : à chaque changement, il faut un moment pour se remettre dans ce qu’on faisait, et ce moment, on ne le sent pas passer. Ce qu’on appelle « faire deux choses à la fois » est le plus souvent une alternance rapide entre deux choses, et chaque aller-retour se paie un peu. Une notification est une petite interruption. Elle ne prend qu’une seconde à l’écran, mais elle te renvoie au début de la phrase que tu étais en train de comprendre.

En 2017, une équipe menée par Adrian Ward a publié une étude plus troublante, qui a beaucoup circulé : des étudiants dont le téléphone était posé sur la table, écran retourné et son coupé, réussissaient moins bien certains exercices de mémoire et de raisonnement que ceux dont le téléphone était resté dans une autre pièce. D’autres équipes ont tenté de retrouver cet effet et n’y sont pas toujours parvenues. La conclusion la plus prudente suffit : éloigner ton téléphone pendant que tu travailles ne te coûte rien, et t’aide probablement. C’est un pari où l’on ne peut presque rien perdre.

On entend aussi beaucoup parler de la lumière bleue des écrans, qui empêcherait de dormir. Les recherches sur la question existent, mais dans la vie ordinaire l’effet semble discuté et sans doute modeste. Ce qui te garde éveillé à minuit, c’est plus souvent ce que tu regardes que la couleur de ce que tu regardes : la vidéo suivante, le message auquel il faut absolument répondre ce soir. Une lampe ne te tient pas éveillé en te racontant une histoire. Un téléphone, si.

Sur les jeunes et les écrans, je n’ai pas de leçon à donner. Ce téléphone est aussi un outil de travail, et un outil cher : quand le forfait acheté pour chercher un cours s’en va en vidéos le soir de révision, tu paies deux fois, une fois en francs et une fois en heures. À Ebolowa, Rodrigue, qui sait ouvrir un téléphone et remonter de l’écran noir jusqu’au connecteur, perd des soirées entières sur le sien. Savoir comment une chose est faite ne protège pas de ce qu’elle fait.

\soustitre{Ce soir, essaie ceci}

Regarde quand tombe ton prochain devoir important : une séquence ou un bac blanc. Compte sept jours en arrière. Pour ces sept jours, choisis une heure de coucher, une seule, réaliste pour ta maison. Si la radio du salon tourne jusqu’à 22 h 30, ne choisis pas 21 h. Écris cette heure en haut d’une page de ton cahier de brouillon. Puis tiens-la, surtout les soirs où tu n’as pas fini, parce que ce sont ces soirs-là qu’elle sert à quelque chose : une heure de coucher qu’on ne respecte que les jours faciles ne protège rien. Ce qui reste à faire, note-le sur une ligne et place-le en tête du lendemain, à l’heure fraîche.

La veille du devoir, ne commence rien de nouveau après le dîner. Le chapitre que tu n’as pas vu ne sera pas compris à minuit, et il emportera avec lui une partie de ce que tu savais. Reprends plutôt les questions que tu t’es écrites les jours précédents, réponds-y sans regarder, vérifie, et éteins. Ce soir-là, aller dormir fait partie de la révision ; c’en est même la dernière séance, et la plus courte.

Pour tes séances de travail, à partir de demain : avant de t’asseoir, confie ton téléphone à quelqu’un de la maison, une sœur ou un parent, pour quarante-cinq minutes. Demande-lui de ne te le rendre qu’à l’heure dite, même si tu le réclames avant. Tu le réclameras. Si personne ne peut le garder, coupe le réseau et mets-le hors de portée, dans une autre pièce ou au fond d’un sac fermé, n’importe où sauf sur la table, même retourné. Quarante-cinq minutes, c’est court, et c’est voulu : assez long pour avancer, assez court pour que tu tiennes.

Quand le temps est écoulé, avant de reprendre ton téléphone, écris dans ton cahier ce que tu as fait pendant ces quarante-cinq minutes, réellement fait. Une ligne suffit. « Trois exercices sur les suites, le troisième faux. » « Relu la leçon sur l’immunité, rien retenu, trop sommeil. » Les deux lignes sont utiles. La seconde surtout, parce qu’elle dit ce que tu ne sais pas encore voir tout seul : l’heure à laquelle ta tête s’arrête.

\asterisme

La semaine suivante, un mardi, à 23 h, Aïcha referme son classeur au milieu d’une page. Il reste un exercice sur la régulation de la glycémie. Elle l’écrit sur une ligne, en haut de la feuille du lendemain, puis elle se lève et barre le jour sur le calendrier, d’un trait net, parce que la journée est finie.

La lampe solaire, de toute façon, clignote déjà. Elle a passé la journée sous un ciel blanc de poussière et n’a pas beaucoup chargé. Aïcha ne la secoue pas, ne la tourne pas vers elle pour gagner encore une ligne. Elle la laisse s’éteindre seule. Dans le noir, la tôle craque en rendant sa chaleur. Près de la porte, Hadja, dans son sommeil, tire le pagne à elle.

\begin{pourcahier}{Pour le cahier}
Pendant cinq soirs, note l’heure de ton coucher. Le lendemain, ajoute une phrase sur ton attention de la journée : à quel moment elle a tenu, à quel moment elle a lâché. Ne conclus rien avant le cinquième soir.
\end{pourcahier}


%% 11-l-encre-rouge.md
\chapitre{9}{L’encre rouge}

Il pleut sur Bafoussam depuis le début de la soirée, une de ces premières pluies de mars qui ne savent pas encore s’arrêter. Sur la tôle, le bruit est serré, si constant qu’au bout d’une heure Hilaire Kamga ne l’entend plus. Il est assis à la table de la salle à manger, sous la lampe de bureau dont il a calé le fil avec un dictionnaire pour qu’elle consente à rester allumée. À sa gauche, la pile des copies de la séquence : sa Seconde C a rendu plus de soixante-dix devoirs, il en a pris quarante pour ce soir. À sa droite, la pile corrigée monte lentement. Entre les deux, la boîte de stylos rouges achetée en septembre, où il n’en reste qu’un.

Il corrige vite, parce qu’il a appris à le faire, et lentement, parce qu’il n’a jamais appris à mal le faire. Exercice 2 : résoudre −2x + 6 > 0. À la quatrième copie, il entoure une ligne. L’élève a bien écrit −2x > −6, puis, très proprement, x > 3. Il fallait x < 3 : quand on divise les deux côtés par un nombre négatif, le sens de l’inégalité se retourne. Kamga écrit « sens » dans la marge et passe à la suivante. La cinquième a fait la même chose. La sixième aussi. Vers la vingtième, il retourne une enveloppe et commence à tracer des bâtons au dos, par groupes de cinq, comme on compte les sacs à la descente d’un camion.

À la quarantième copie, il y a trente bâtons.

Il pose le stylo. Dehors, la pluie redouble, puis se calme, comme si elle aussi reprenait son souffle. Trente fois la même erreur sur quarante copies, ce n’est pas trente élèves nuls, c’est une leçon qu’il faut refaire. Il se souvient très bien de cette leçon : un jeudi, la dernière heure, une classe qui pensait déjà au portail. Il avait énoncé la règle, l’avait encadrée au tableau, avait fait un exemple et était passé au suivant parce que le programme, lui, n’attend personne. La raison de la règle, il l’avait gardée pour lui, faute de temps. Trente élèves avaient retenu qu’il existait une histoire de sens qui change, sans savoir quand, ni pourquoi, et le jour du devoir, devant une ligne où rien ne clignotait, ils avaient fait ce que fait n’importe qui : la chose la plus naturelle.

Kamga n’aime pas cette découverte. Elle lui coûtera une heure de cours qu’il n’a pas, dans un trimestre déjà en retard. Il la note quand même sur l’enveloppe, sous les bâtons : \emph{refaire inéquations, lundi, avec la droite}.

\asterisme

Une erreur isolée dit quelque chose d’un élève. Trente fois la même dit quelque chose d’une leçon. Et entre les deux, il y a ta copie, celle qui te revient un mardi matin, barrée de rouge, que tu plies en quatre avant que ton voisin ne la voie.

Une copie, d’habitude, se lit en trois secondes : deux pour la note en haut à droite, une pour celle du voisin, et c’est fini. Si la note est bonne, on ne regarde pas le reste parce qu’on n’en a pas besoin ; si elle est mauvaise, parce que ça fait mal. Dans les deux cas, le travail du correcteur, ces ronds et ces mots qu’il a tracés tard le soir, ne sert à personne. C’est dommage, parce que la note est la partie la moins bavarde de la copie : elle donne un total, et les adresses sont dans les marges.

Une erreur, regardée de près, est une trace. Elle montre l’endroit exact où ta pensée a pris une autre direction que celle qu’il fallait. Une copie blanche ne montre rien ; une copie fausse montre beaucoup, à condition de la lire comme Kamga lit ses trente bâtons, en cherchant ce qui se répète plutôt qu’en comptant ce qui manque. Le rouge marque l’endroit. Trouver la raison, c’est ton travail.

Pour la chercher, il aide de savoir que les erreurs se rangent en familles. Certaines sont des erreurs de lecture : la consigne demandait de vérifier, tu as résolu ; elle demandait deux exemples, tu en as donné un ; tu as lu « augmente » là où il était écrit « diminue ». D’autres sont des erreurs d’attention : tu savais, mais tu as recopié un 3 qui était un 8, oublié un signe moins entre deux lignes, accordé un participe avec le mauvais mot parce que tu pensais déjà à la phrase suivante. D’autres encore sont des erreurs de méthode : tu connaissais la leçon, mais devant l’exercice tu ne savais pas par où commencer, ni dans quel ordre avancer. Les dernières, enfin, sont des erreurs de notion : quelque chose n’a pas été compris, et tout ce qu’on a construit dessus penche.

Chaque famille appelle un remède différent, et c’est tout l’intérêt du tri. Une erreur de lecture se soigne en soulignant les verbes de la consigne avant d’écrire une ligne. Pour l’attention, on se garde cinq minutes à la fin, et l’on relit lentement, ligne par ligne, en refaisant les calculs au lieu de les regarder. Une erreur de méthode demande des exercices résolus, étudiés puis refaits sans regarder. Une erreur de notion renvoie à la leçon, et souvent à une question qu’on ira poser. Soigner une erreur de notion en « faisant plus attention » ne sert à rien ; soigner une étourderie en réapprenant tout le chapitre est une perte de temps. Le diagnostic évite les deux. Il a aussi un avantage plus discret. « J’ai encore raté les maths » fait mal et ne mène nulle part ; « sur cinq erreurs, trois viennent de consignes mal lues » mène à un geste précis, qu’on peut faire dès le prochain devoir, avec un crayon.

Les trente élèves de Kamga ont fait une erreur de notion déguisée en étourderie : personne ne leur avait montré qu’il y avait une raison de retourner le signe.

Kamga a remarqué autre chose en corrigeant. Beaucoup de ces x > 3 étaient écrits d’une main ferme, certains soulignés deux fois. Ses élèves étaient sûrs d’eux, et ce détail est plutôt une bonne nouvelle. Des travaux menés par Janet Metcalfe et ses collègues suggèrent que les erreurs commises avec assurance, une fois corrigées, sont souvent particulièrement bien retenues, mieux que celles qu’on avait faites en hésitant. Les chercheurs appellent cela l’effet d’hypercorrection. L’explication la plus probable tient en un mot : la surprise. Quand on découvre qu’on s’est trompé sur une chose dont on était certain, l’attention se réveille d’un coup, et ce qu’on apprend à ce moment-là s’accroche. Personne ne te conseille de te tromper exprès. Mais une erreur dont tu étais fier, regardée en face, peut devenir l’une des choses que tu sauras le mieux. À condition de la regarder.

\asterisme

Le lundi, Kamga ne commence pas par la règle. Il écrit au tableau 2 < 5, et demande à la classe ce qui se passe si l’on multiplie les deux côtés par −1. Il ne donne pas la réponse. Il laisse chercher, cinq minutes, pendant que la craie attend sur le rebord. Des chercheurs ont observé que laisser des élèves se débattre avec un problème avant de leur montrer la solution peut, dans certaines conditions, les aider à mieux comprendre ensuite ; dans d’autres conditions, cela les décourage seulement. Kamga ne connaît pas ces travaux par cœur. Il connaît sa classe, et il sait que cinq minutes, c’est assez pour que la question devienne la leur, pas assez pour qu’ils la détestent.

Puis il dit : « Qui préfère devoir deux sacs de ciment à quelqu’un, et qui préfère en devoir cinq ? » On rit. Personne ne préfère devoir cinq. « Donc, quand on parle de dettes, moins deux, c’est mieux que moins cinq. Deux était plus petit que cinq. Moins deux est plus grand que moins cinq. Le signe moins a tout retourné. » Il trace une droite, place les nombres, et montre que multiplier par −1, c’est regarder la droite dans un miroir : ce qui était à gauche passe à droite. Le sens de l’inégalité, lui, ne fait que suivre les nombres.

L’heure se termine. La plupart des élèves sortent en courant, parce qu’il pleut encore et que le préau est petit. Divine Nfor reste. Elle attend que le professeur ait effacé le tableau, qu’il ait rangé ses craies, qu’il ne soit plus en train de regarder l’heure. Puis elle pose sur le bureau un carnet bleu, ouvert.

Kamga reconnaît le carnet sans l’avoir jamais vu. Les premières pages, qu’elle n’a pas pensé à cacher, sont pleines de phrases en français, une par ligne, avec à côté une règle d’accord, parfois une traduction en anglais entre parenthèses. Il ne pose aucune question sur ces pages-là. La page ouverte, elle, est nouvelle. À gauche, Divine a recopié sa propre ligne du devoir : −2x > −6, donc x > 3. Au milieu, la correction. À droite, une question, écrite au crayon puis repassée au stylo :

\emph{Ce matin j’ai compris QUE le sens change. Je ne comprends pas encore POURQUOI quand on divise. Vous avez montré la multiplication. Diviser par −2, c’est la même chose que multiplier par −1/2 ?}

Kamga lit deux fois. Il est fatigué, il a une autre classe dans dix minutes à l’autre bout de la cour, et c’est la meilleure question qu’on lui ait posée depuis la rentrée. Il le lui dit, plus sèchement qu’il ne voudrait, parce qu’il n’a pas l’habitude de dire ces choses-là. Puis il répond : oui, diviser par un nombre, c’est multiplier par son inverse, et l’inverse d’un négatif est négatif. Le miroir est toujours là. Divine écrit dans la marge, en anglais, un seul mot : \emph{mirror}. Elle referme le carnet et s’en va. Elle a mis moins de trois minutes.

Le rouge, sur une copie, est une forme de retour : quelqu’un a lu ce que tu as fait et te dit ce qu’il en pense. Les synthèses de John Hattie et Helen Timperley, qui ont rassemblé un grand nombre d’études sur le sujet, suggèrent que ce retour d’information peut compter parmi les leviers les plus puissants de l’apprentissage, et qu’il peut aussi, mal donné, ne servir à rien ou même nuire. Un bon retour dit où l’on en est, où l’on devait arriver, et ce qu’il faut faire ensuite. Une note seule ne dit que le premier point. Un mot dans la marge, « sens », dit un peu du deuxième. Le troisième, presque personne n’a le temps de l’écrire.

Kamga a plus de soixante-dix élèves dans cette classe, et d’autres classes. S’il écrivait trois lignes utiles sur chaque copie, il y passerait ses nuits, et il y passe déjà ses soirées. Un professeur qui rend des copies avec une note et quelques ronds rouges a souvent fait ce qu’il pouvait dans le temps qu’il avait ; y lire de l’indifférence serait injuste. La suite du retour, celle qui dit quoi faire ensuite, reste donc à écrire, et Divine l’a écrite elle-même. Son carnet bleu ne contient rien d’autre : la partie de la correction que le professeur n’avait pas le temps de faire, faite par l’élève. L’encre rouge ne sert à rien si la copie va directement au fond du sac. Elle commence à servir le jour où tu la recopies.

Demander de l’aide est une compétence, au même titre que résoudre une inéquation, et elle s’apprend de la même façon : en le faisant mal, puis un peu mieux. Elle a trois parties qu’on oublie souvent. À qui : le professeur, bien sûr, mais il n’est qu’une des personnes qui savent ; il y a aussi un élève d’une classe au-dessus, une grande sœur, le surveillant qui a fait des mathématiques, un camarade qui a réussi l’exercice. Quand : à un moment où la personne peut s’arrêter, à la fin d’un cours plutôt qu’entre deux portes, et surtout avant que l’erreur ne devienne une marche manquante sur laquelle on construira tout le chapitre suivant. Et avec quelle question : une question précise, qui montre où l’on est arrivé et où l’on s’est arrêté. On a vu, avec le ticket de Grâce, ce que change une question écrite à l’avance sur un bout de papier ; Divine en donne une version de plus. Sa question tient en trois lignes, elle nomme la ligne fausse, elle dit ce qu’elle a déjà compris. Un professeur pressé peut y répondre en trois minutes. À « je n’ai rien compris », il ne peut répondre qu’en refaisant le cours, et il n’a pas le temps.

\asterisme

Le vendredi après les cours, quand la pluie a enfin cessé, trois élèves se sont installés sous le manguier de la cour avec un petit tableau d’ardoise emprunté à une salle vide, appuyé contre le tronc. La terre est encore mouillée. Quelqu’un a apporté un morceau de craie trouvé dans la rigole du tableau de la Seconde B.

Divine est là, avec une fille du premier banc qui fait ses dictées presque sans fautes et ne comprend rien aux valeurs absolues, et un garçon qui comprend tout en mathématiques à condition qu’on ne lui demande pas de l’écrire. Ils ont fixé une règle, et ils la tiennent mieux qu’ils ne l’auraient cru : chacun passe au tableau à tour de rôle, explique un exercice sans ses notes, et les deux autres ont le droit de l’interrompre pour demander pourquoi. Divine refait l’inéquation, avec le miroir et les sacs de ciment. Le garçon, qui l’écoute en mangeant des arachides, lui demande pourquoi on ne retourne pas le signe quand on ajoute −6 des deux côtés ; elle commence à répondre, s’arrête, repart, et trouve en parlant une chose qu’elle n’avait pas trouvée en écrivant. Puis la fille du premier banc prend la craie et explique à Divine l’accord du participe passé avec avoir, avec une patience un peu agacée, en anglais quand il le faut, parce qu’elle s’est aperçue que ça va plus vite.

Ce qui se passe sous le manguier a été étudié, sous d’autres noms. Expliquer à quelqu’un oblige à mettre de l’ordre dans ce qu’on sait ; on découvre en parlant les trous qu’on ne voyait pas en relisant. Des chercheurs ont même observé que le simple fait de s’attendre à devoir enseigner une leçon pouvait améliorer la façon dont on l’apprend, avant d’avoir enseigné quoi que ce soit. Et le tutorat entre élèves profite souvent aux deux : à celui qui reçoit l’explication, parce qu’elle vient de quelqu’un qui vient lui-même de comprendre et se souvient de ce qui bloquait, et à celui qui la donne, parce qu’il doit la fabriquer. Divine apprend les participes ; la fille du premier banc apprend, en les expliquant, qu’elle ne savait pas très bien pourquoi elle les accordait juste. Sous le manguier, le professeur change toutes les vingt minutes.

Un groupe de travail peut aussi devenir très vite un salon, et un salon très agréable. On parle du match et de la fête de samedi, on rit beaucoup, et à la fin de l’après-midi on a l’impression d’avoir travaillé parce qu’on était ensemble avec des cahiers. Les trois élèves du manguier tiennent avec peu de choses : une heure, pas plus, mesurée à la montre du garçon ; un seul objectif, décidé avant de s’asseoir ; chacun au tableau, sans exception, même celui qui préfère regarder ; et jamais de devoir recopié. Sur ce dernier point, une phrase suffit : recopier le devoir d’un camarade, c’est emprunter une note, mais la marche que le devoir devait construire reste vide, et c’est toi qui devras monter dessus le jour où personne ne sera assis à côté.

Kamga traverse la cour à ce moment-là, son cartable sous le bras, et voit le petit tableau d’ardoise contre le manguier. Il reconnaît son miroir, dessiné à la craie, maladroitement. Il ne s’arrête pas. Mais le surveillant, qui marchait à côté de lui, l’entend rire tout seul, une fois, et se demande longtemps de quoi.

\soustitre{Ce soir, essaie ceci}

Prends ta dernière copie corrigée, celle que tu as peut-être pliée en quatre. Déplie-la. Et ouvre un carnet d’erreurs : un petit carnet si tu en as un, sinon les dernières pages de ton cahier de brouillon, en partant de la fin, pour qu’elles ne se mélangent pas avec le reste.

Pour chaque erreur de la copie, une seule, puis la suivante, écris quatre choses, chacune en une phrase. D’abord ce que tu as fait, exactement, sans adoucir : « J’ai divisé par −2 sans changer le sens. » Ensuite pourquoi tu l’as fait, honnêtement, même si la réponse est « je ne sais pas » ou « je n’ai pas lu la question jusqu’au bout ». Puis ce qu’il fallait faire, avec la règle et, si tu la connais, la raison de la règle. Enfin, un exercice semblable, que tu prends dans ton manuel ou que tu fabriques toi-même en changeant les nombres. À côté de cet exercice, écris une date : dans trois jours. Ce jour-là, tu le referas sans regarder la correction. Si tu le réussis, tu reviendras sur cette erreur une semaine plus tard ; si tu le rates, tu la remets à trois jours. C’est la façon dont Nadège révise ses cartons, appliquée à tes erreurs : le carnet se refait, à dates fixes, sans regarder.

Avant de refermer le carnet, regarde tes quatre ou cinq erreurs et note, au bout de chacune, sa famille : lecture, attention, méthode, notion. Si la même famille revient trois fois, tu tiens quelque chose de plus utile qu’une note : ton point faible du moment, et il a un remède.

Si une erreur te résiste, si même avec la correction sous les yeux tu ne comprends pas pourquoi c’était faux, ne t’acharne pas ce soir. Écris ta question sous l’erreur, en nommant la ligne, et apporte-la à quelqu’un dans la semaine.

Et cette semaine, justement, propose à deux camarades une heure de tour de tableau. Pas besoin de tableau, d’ailleurs : une grande feuille ou le dos d’un cahier suffisent. Choisissez ensemble un chapitre avant de commencer. Chacun prépare un exercice qu’il explique aux deux autres, à voix haute, sans ses notes ; les autres ont le droit, et même le devoir, de demander « pourquoi ? » à chaque étape. Quand l’heure est finie, on s’arrête, même au milieu d’un exercice. Tu remarqueras peut-être que c’est en expliquant, plus qu’en écoutant, que tu découvriras ce que tu ne savais pas. Note-le dans ton carnet.

\asterisme

Quelques soirs plus tard, il pleut de nouveau sur Bafoussam, moins fort. Kamga corrige les copies de son autre seconde, sous la même lampe, avec le même dictionnaire sur le fil. À la dix-huitième, au milieu d’un mot de la marge, le stylo rouge s’épuise. Il écrit « pres », puis plus rien qu’un sillon pâle dans le papier. Il secoue le stylo, souffle dessus, le secoue encore. C’était le dernier de la boîte de septembre. Il cherche dans le cartable, puis dans le tiroir du buffet. Il ne trouve qu’un stylo bleu, celui de la liste des courses.

Il hésite un moment, puis finit le mot au bleu : « que ». « Presque », en deux couleurs. Il continue avec le bleu jusqu’à la dernière copie. Il entoure, il écrit dans les marges les mêmes mots qu’avant, aussi courts, aussi pressés, et vers la fin il s’aperçoit que ça ne change rien à ce qu’il écrit. Une erreur entourée en bleu reste une erreur, et un endroit où quelqu’un a fait attention.

Avant d’éteindre, il ouvre le cahier vert. Il le tient depuis des années, et il lui arrive de passer des semaines sans trouver quoi y mettre pour certains élèves. Il trouve la ligne de Divine Nfor, sous ce qu’il avait écrit en novembre sur sa démonstration. La pluie s’est presque arrêtée ; on n’entend plus sur la tôle que quelques gouttes tombées du manguier du voisin. Il écrit, au bleu, puisque c’est ce qui reste : « Est venue avec sa question écrite. »

\begin{pourcahier}{Pour le cahier}
Recopie ta dernière erreur importante, telle que tu l’as faite, sans la corriger. À côté, écris de quel type elle est : lecture, attention, méthode ou notion. Dessous, la date où tu referas un exercice du même genre, sans regarder. Ce jour-là, reviens ici et écris ce qui s’est passé.
\end{pourcahier}


%% 12-la-copie-double.md
\chapitre{10}{La copie double}

À 5 h 15, le pont sur le Wouri est presque vide. Grâce Ewane le traverse dans un taxi où l’on n’est encore que quatre, et c’est la première fois qu’elle voit le fleuve sans entendre klaxonner derrière elle. Ce matin, il n’y a pas d’embouteillage, et ce vide lui paraît plus inquiétant qu’un bouchon. Sa mère a tenu à ce qu’elle parte trop tôt. Sur ses genoux, Grâce serre sa trousse fermée, comme on tient un sac au marché quand il y a foule.

Le centre d’examen est un établissement qu’elle ne connaissait que de nom, de l’autre côté de la ville, avec un portail vert et une cour de terre battue déjà balayée. Devant le portail, une file s’est formée avant le jour. Des candidats qu’elle n’a jamais vus révisent debout, une fiche à la main, les lèvres bougeant à peine ; d’autres ne révisent rien et parlent trop fort. Un garçon mange un beignet en le regardant comme s’il pouvait contenir une formule. Quelqu’un récite à mi-voix une définition d’histoire, toujours la même, et se trompe toujours au même mot.

À 7 h, un homme en chemise claire sort avec une liste et commence l’appel. L’ordre n’a rien à voir avec celui de la classe. D’ailleurs, ici, la classe n’existe plus : seulement des numéros de table, des noms qui se suivent selon une logique connue de la seule administration, et Grâce entend le sien entre deux noms inconnus, avec l’accent sur la mauvaise syllabe. Elle répond « présente » d’une voix qu’elle ne reconnaît pas. Elle tend sa pièce d’identité et sa convocation à une dame qui compare la photo et le visage avec l’application d’une douanière, puis lui fait signe d’avancer.

La salle sent la craie et la peinture fraîche. On a repeint les murs pour l’occasion, un jaune pâle qui colle encore un peu au doigt, et Grâce prend soin de ne pas s’y appuyer. Sur chaque banc, un numéro écrit à la craie ; le sien est près de la fenêtre, au troisième rang. On est seul sur son banc, ce qui, après des années à trois, donne une impression étrange, de place et d’abandon à la fois. Au tableau, d’une écriture appliquée : la date, l’épreuve (mathématiques), l’heure de début et l’heure de fin.

Les surveillants passent entre les rangs et posent sur chaque banc une copie double, blanche, à grands carreaux, avec en haut le cadre où l’on écrit son nom. Puis une feuille de brouillon. Puis, face cachée, le sujet, que personne n’a le droit de retourner avant le signal. Grâce a les mains moites ; elle les essuie sur sa jupe, une fois, deux fois, et elles redeviennent moites aussitôt, comme si elles avaient leur propre avis sur la matinée. Dans sa trousse, entre le compas et la règle, dort un bout de rouleau de tickets de caisse pris au kiosque de sa tante. Au dos, au stylo, trois lignes. Elle ne les relit pas. Elle sait qu’elles sont là.

Il reste quelques minutes avant le signal. Le sujet est posé devant elle, blanc au dos, parfaitement muet.

\asterisme

Ce qui se passe dans le corps de Grâce à cet instant, tu le connais. Le cœur qui cogne comme si l’on venait de monter quatre étages en courant. Les mains. Une envie pressante d’aller aux toilettes alors qu’on en revient. La bouche sèche, la langue qui colle. Et parfois cette blancheur dans la tête, où les choses sues la veille semblent avoir déménagé sans laisser d’adresse.

C’est normal, d’abord. Cela ne signale ni un manque de préparation ni une fragilité particulière. Presque tout le monde, dans cette salle, a le cœur qui cogne ; certains le cachent mieux, voilà tout. Le garçon au beignet aussi, probablement, et c’est peut-être pour cela qu’il mangeait si lentement. Le stress d’examen est la réaction ordinaire d’un corps qui sent qu’un moment compte. Et le Probatoire compte.

On entend souvent qu’un peu de stress aide et que trop de stress nuit, et l’on dessine pour le montrer une courbe en cloche : à gauche, l’élève trop détendu qui bâille ; au sommet, celui qui est juste assez tendu pour être vif ; à droite, celui qui l’est tant qu’il ne pense plus. Cette courbe s’appuie sur de vieilles expériences menées sur des animaux au début du siècle dernier, et on lui a fait dire depuis bien plus qu’elle ne disait. C’est une simplification. Garde l’idée, qui a du bon sens : une certaine tension te réveille, une tension excessive t’encombre. Laisse de côté la certitude, et le dessin trop propre. Personne ne sait où se trouve ton sommet à toi, pas même toi.

Pourquoi l’excès encombre, tu le sais depuis la main de Grâce sous la table, un jour de pluie : l’inquiétude occupe une partie de la place dont tu aurais besoin pour raisonner. La blancheur, dans ta tête, est une pièce encombrée, et une pièce encombrée se range.

Une équipe de chercheurs, autour de Jeremy Jamieson, a travaillé en 2010 avec des étudiants américains qui préparaient un examen d’admission pour poursuivre leurs études. Avant un test d’entraînement, on a fait lire à certains d’entre eux un court texte : les signes du stress, le cœur rapide, les mains moites, peuvent être la manière dont le corps se prépare à un effort, plutôt qu’un danger. Ceux-là ont mieux réussi la partie mathématique que les autres. Une seule étude ne fait pas une loi, et les travaux qui ont suivi sont encourageants sans être toujours aussi nets. Je n’en tire pas une recette, plutôt une façon de lire les signes : au moment où ton cœur s’emballe, tu peux te dire, si ça t’aide, que ton corps se prépare.

Ton corps ignore que l’adversaire est une feuille de papier ; il fait ce qu’il ferait avant une course, et il a raison de prendre la matinée au sérieux. Tu n’as pas besoin qu’il se taise. Tu as besoin de ne pas prendre son bruit pour un verdict.

Le souffle, lui, se travaille. Grâce chante alto dans une chorale, et la cheffe de chœur a passé des dimanches entiers à apprendre aux altos à prendre l’air sans bruit et à le rendre lentement, sur toute la longueur d’une phrase musicale. Grâce y voyait une affaire de musique ; c’est aussi, modestement, une arme contre la peur. Respirer lentement, en laissant l’expiration durer plus longtemps que l’inspiration, aide beaucoup de gens à se calmer ; tu n’as pas besoin qu’on te dessine les nerfs et les organes pour le vérifier. Inspire en comptant jusqu’à quatre, souffle en comptant jusqu’à six, sans forcer, cinq ou six fois. Personne ne le remarquera. Les têtes autour de toi sont toutes occupées par leur propre cœur.

Sur le rouleau de tickets, la première ligne dit : \emph{Respirer.} Grâce ne touche pas encore au sujet. Elle laisse partir l’air comme au début d’un cantique qu’elle connaît par cœur, et elle sent, pas tout de suite, ses épaules redescendre d’un centimètre.

\asterisme

Ce qui se joue ce matin s’est pourtant décidé en grande partie avant, les samedis et les dimanches des mois précédents. C’est plutôt une bonne nouvelle : ces après-midi-là dépendaient davantage de toi que cette matinée-ci.

La meilleure préparation à un examen, c’est l’examen : une épreuve entière, faite en vrai, dans des conditions qui ressemblent à celles du jour. Une annale, une table dégagée, un chronomètre, une copie double, et rien d’autre, ni téléphone ni cahier ouvert « au cas où ». Relire tes cours une dernière fois ou recopier les formules au propre sur une belle feuille ne t’apprendra jamais cela. Tu te souviens des cartes de Nadège et de cette découverte têtue qu’on retient mieux ce qu’on va chercher dans sa mémoire que ce qu’on relit. L’annale chronométrée, c’est la même idée en grandeur réelle. Au lieu de trois questions sur une leçon, tu affrontes toutes celles d’une épreuve, dans le temps de l’épreuve, sans personne pour t’aider, et tu vois ce que tu sais vraiment faire.

Tu apprends aussi des choses qu’aucune leçon n’enseigne. Le temps passe plus vite quand il est compté, et la main fatigue au bout d’une heure d’écriture. Peut-être as-tu l’habitude, en travaillant, de jeter un œil à ton téléphone toutes les vingt minutes ; privé de ce geste, tu ne sais plus quoi faire de tes doigts. Et il est probable que tu t’attardes sur l’exercice que tu aimes pour bâcler celui que tu crains. Tout cela, mieux vaut le découvrir un samedi chez toi que devant le sujet du Probatoire.

Grâce l’a fait, pas toutes les semaines. Le samedi, elle tient le kiosque avec sa tante, entre les transferts d’argent et les recharges de crédit, et le dimanche matin, il y a la chorale. Mais plusieurs dimanches après-midi, depuis mars, elle a dégagé la table où sa mère coud, poussé les coupons de pagne et les bobines, posé le réveil de la maison devant elle et fait une annale de mathématiques en entier. La première fois, elle n’a pas fini. Le dimanche suivant, elle a fini en bâclant la fin. À la troisième annale, elle s’est corrigée elle-même, sévèrement, avec le corrigé et le barème, et elle a compris que la moitié de ses points perdus venaient de consignes mal lues, sur des notions qu’elle savait. La découverte n’avait rien d’agréable. Elle était utile, et elle tenait au dos d’un ticket.

Une autre piste circule, plus fragile. En 2011, Gerardo Ramirez et Sian Beilock ont demandé à des élèves d’écrire, pendant une dizaine de minutes juste avant un examen, ce qui les inquiétait à propos de cet examen. Les plus anxieux de ceux qui avaient écrit ont obtenu de meilleurs résultats que les élèves aussi anxieux qui ne l’avaient pas fait. L’idée était que poser l’inquiétude sur le papier la sortait un peu de la tête, où elle encombrait. C’est séduisant. Mais quand une grande équipe a refait l’expérience, en 2018, avec beaucoup plus de participants, elle n’a pas retrouvé l’effet, et d’autres essais en classe sont restés décevants. Je te la donne donc comme une piste fragile, rien de plus. Si tu es de ceux que l’inquiétude envahit, essaie-la un jour d’entraînement, avant une annale, et regarde honnêtement si elle t’aide, toi.

Grâce l’a essayée dans la file, ce matin, debout contre le mur du portail, sur un autre bout de rouleau. Elle a écrit petit : \emph{J’ai peur de ne pas comprendre la première question, et que ma tête se vide.} Puis, après un moment : \emph{Tantine parle déjà d’infirmerie comme si c’était fait, et moi je ne sais même pas si je veux.} Elle a plié le papier en quatre et l’a glissé dans sa poche, pas dans la trousse. Celui-là, elle ne le relira pas. Il a fait son travail en étant écrit, ou il ne l’a pas fait ; elle ne le saura jamais, et ce matin cela n’a aucune importance.

La veille et le matin, enfin, ne demandent aucun génie. La veille, tu n’apprends rien de nouveau, et tu dors : on a vu pourquoi dormir n’est pas tricher. Le matin, tu manges quelque chose, même si ton ventre est noué ; Grâce a avalé debout le pain que sa mère lui tendait, sans faim, parce que sa mère ne l’aurait pas laissée partir autrement. Et tu arrives tôt. Tu ne réviseras guère dans la file, où l’on s’affole surtout les uns les autres ; tu arrives tôt pour que l’embouteillage, la convocation oubliée ou la salle introuvable ne deviennent pas une épreuve de plus avant l’épreuve. Le pont vide, à 5 h 15, était un cadeau. On ne reconnaît ce genre de cadeau qu’après coup.

\asterisme

Ce matin de juin, des salles comme celle-ci, il y en a dans tout le pays. À Garoua, une élève de Terminale a déjà composé ; elle barre sur son calendrier mural les jours qui la séparent des résultats. À Bafoussam, Hilaire Kamga surveille une salle qui n’est pas la sienne, les mains dans le dos, et il regarde les mains des candidats plus que leurs copies.

Dans la salle de Grâce, le surveillant principal lit les consignes d’une voix qui a déjà lu les mêmes l’an dernier. Les téléphones sont éteints et déposés dans un carton près de la porte ; Grâce y a laissé le sien avec les autres, et l’on dirait un petit cimetière d’écrans noirs. Aucun document, aucun échange, aucune sortie avant l’heure fixée. Quelqu’un, au fond, fait trembler sa jambe sous le banc, et le banc avec.

Quand le sujet se retournera, les choses auront un ordre, et cet ordre s’apprend.

Lis d’abord le sujet en entier, du début à la fin, avant d’écrire la moindre ligne sur la copie. Cela paraît du temps perdu quand l’horloge tourne ; ces quelques minutes t’évitent pourtant de découvrir, à un quart d’heure de la fin, que le dernier exercice était le plus facile et rapportait le plus de points. Pendant cette lecture, mets au crayon une petite marque à côté de ce que tu sais faire, et une autre, différente, à côté de ce qui te fait peur. Tu viens de dresser la carte de ta matinée.

Commence ensuite par ce que tu sais. L’ordre du sujet est celui de la personne qui l’a conçu, et rien ne t’oblige à le suivre, sauf quand une question dépend de la précédente. Réussir quelque chose dans le premier quart d’heure fait baisser le bruit dans la tête : la pièce se range un peu, et tu retrouves de quoi penser au reste.

Répartis ton temps. Regarde combien de points vaut chaque exercice et accorde-lui à peu près la part de la durée qui lui revient. Écris dans un coin de ton brouillon l’heure à laquelle tu quitteras chaque exercice, même inachevé, et tiens-toi à cette heure. Un exercice à moitié fait et un autre à moitié fait valent souvent plus qu’un exercice parfait et un autre jamais commencé.

Laisse de la place. Quand tu sautes une question, numérote-la quand même sur ta copie et garde en dessous quelques lignes blanches, pour pouvoir y revenir sans écrire en tout petit dans la marge. Un correcteur fatigué, à sa quatre-vingtième copie de la journée, sera plus juste avec une copie où il s’y retrouve. (Je pense à lui aussi. Il le mérite.)

Et garde les dix dernières minutes pour relire. Relire, ici, veut dire vérifier ce que tes dimanches d’entraînement t’ont appris à soupçonner : les signes, les unités, les numéros des questions, une consigne lue trop vite. C’est souvent là que dorment les points les plus faciles à reprendre.

Il circule, dans certaines files d’attente, une autre stratégie qu’on se chuchote : le papier roulé dans la manche, le téléphone qu’on n’a pas déposé. Je n’en parlerai qu’une fois. La fraude vole d’abord celui qui la commet : une note empruntée, encore, posée sur une marche qui reste à construire. Et si l’on est pris, elle peut coûter beaucoup plus que l’examen lui-même.

Après l’épreuve, enfin, vient la cour. Un cercle se formera autour de ceux qui sont sûrs d’eux, et l’on refera le sujet à voix haute. Quelqu’un dira que c’était facile. Quelqu’un annoncera pour l’exercice deux un résultat qui n’est pas le tien, avec l’assurance de celui qui n’a jamais douté de rien. N’entre pas dans ce cercle. Tu n’y apprendras rien que tu puisses encore changer, et tu as peut-être une autre épreuve l’après-midi ou le lendemain : ta tête te sera plus utile pour celle-là que pour défendre la précédente. Bois de l’eau, mange, va t’asseoir à l’ombre, regarde ailleurs.

Un mot plus grave. Si, chaque fois qu’un examen approche, ton cœur s’emballe au point que tu crois mourir, si tu fais des malaises à répétition ou que tu as déjà dû quitter une salle, cela dépasse le trac ordinaire, et le courage n’a rien à y voir. Parles-en à un médecin, à l’infirmerie de ton établissement, à un adulte de confiance. Et fais-le maintenant, des semaines avant la période des examens, plutôt que la veille au soir. Ce que tu vis porte un nom, et il existe des personnes dont c’est le métier d’aider à le traverser.

\soustitre{Ce soir, essaie ceci}

Prends ton calendrier, ou la dernière page de ton cahier de brouillon, et choisis un samedi matin, dans les deux ou trois semaines qui viennent. Regarde à quelle heure commence l’épreuve que tu redoutes le plus le jour de l’examen, et réserve exactement cette heure-là. Si ton samedi est déjà pris, comme celui de Grâce, prends un dimanche après-midi ; l’important est que le créneau soit entier et que personne ne compte sur toi pendant ce temps.

D’ici là, trouve une annale de cette matière, avec son corrigé. Demande à ton professeur, à un ancien élève, à un grand frère, à la bibliothèque du lycée, à un camarade qui en a une pile. Ne la regarde pas en la recevant. Range-la, face cachée, comme le sujet de Grâce.

Le jour venu, dégage une table. Pas de téléphone : donne-le à quelqu’un de la maison, ou, s’il te sert de chronomètre, mets-le hors réseau, face contre la table, à l’autre bout de la pièce, avec une alarme réglée sur l’heure de fin. Prends une copie double, ou, à défaut, quatre pages de ton cahier. Écris ton nom en haut. Puis retourne l’annale et fais l’épreuve entière, dans le temps exact, sans t’arrêter pour vérifier une formule dans ton cours. Si tu bloques, fais ce que tu ferais le jour de l’examen : passe à la suite.

Le lendemain, corrige-toi avec le corrigé et le barème, honnêtement, comme un correcteur sévère qui ne te connaît pas. Ne t’accorde pas les points de « ce que tu voulais dire ». Puis, dans ton cahier, note pour chaque point perdu d’où il vient : une notion que tu ne savais pas, une consigne mal lue, une erreur de calcul, ou le temps qui a manqué. Ce relevé-là t’en dira davantage que la note.

Et dès ce soir, sans attendre le samedi, prends un petit papier, un ticket si tu en as un, et écris tes trois gestes des dix premières minutes. Les tiens, pas ceux de Grâce. Par exemple : respirer, lire tout le sujet, choisir l’exercice par lequel commencer. Le samedi venu, tu les essaieras avant de retourner l’annale, pour voir s’ils te vont, et tu les corrigeras s’il le faut.

\asterisme

« Vous pouvez retourner le sujet. »

Le bruit est le même dans toutes les salles du pays : des dizaines de feuilles qui se retournent presque ensemble, un froissement bref, puis rien. Grâce respire, quatre temps pour prendre l’air, six pour le rendre. Deuxième ligne du ticket : \emph{Tout lire.} Elle lit tout, jusqu’au barème en bas de la dernière page. La première question, elle ne la comprend pas. Elle relit l’énoncé, ne la comprend toujours pas, met une petite croix au crayon dans la marge et laisse dix lignes blanches sur sa copie. Troisième ligne : \emph{Commencer par ce que je sais.} La deuxième question ressemble à une annale d’un dimanche d’avril, celle où elle avait perdu des points pour avoir mal lu. Elle la lit deux fois. Elle commence.

Mais avant tout cela, avant même le signal, pendant que le surveillant finissait de lire les consignes et que le banc du fond tremblait encore, elle avait fait la seule chose qui soit entièrement sûre ce matin-là. En haut de la copie double, dans le cadre prévu, en capitales, elle avait écrit son nom.

\begin{pourcahier}{Pour le cahier}
Écris tes trois gestes des dix premières minutes, chacun sur une ligne, avec un verbe au début. Puis recopie-les sur un ticket ou un coin de feuille, que tu glisseras dans ta trousse et que tu emporteras le jour de l’examen.
\end{pourcahier}


%% 13-ce-qui-ne-depend-pas-de-toi.md
\partie{Le monde}{III}

\chapitre{11}{Ce qui ne dépend pas de toi}

Les soirs de mars, quand les amphis se vident, le cybercafé de Ngoa-Ekellé se remplit. C’est une pièce étroite, plus profonde que large, avec huit postes alignés contre le mur, une imprimante qui chauffe au fond et un comptoir en contreplaqué derrière lequel Nadège Kenfack travaille trois soirs par semaine. Elle encaisse et note les heures de connexion dans un cahier à spirale. Le reste du temps, elle imprime des mémoires pour des étudiants de master qui arrivent tous à la dernière minute et parlent tous de leur directeur avec la même voix. Quand l’imprimante avale une feuille de travers, c’est elle qu’on appelle. Elle ouvre le capot, tire doucement, jamais d’un coup, et la page ressort pliée en accordéon, avec la moitié d’un titre de chapitre imprimée de biais. Ça sent le papier chaud et la poussière de ventilation. L’onduleur, sous le comptoir, émet de temps en temps un petit bip que personne n’écoute plus.

Entre deux clients, elle révise. À côté de la caisse attend un paquet de cartes tenu par un élastique, découpées dans des cartons de lait en poudre : d’un côté une question au feutre, de l’autre la réponse. Elle en tire une, lit la question, retourne la carte face contre le bois, répond à mi-voix, vérifie. Les étapes de la mitose, dans l’ordre. Le rôle des ribosomes. Une carte qu’elle rate pour la troisième fois cette semaine retourne en tête du paquet. Quand quelqu’un entre, elle glisse la carte sous la caisse, face cachée, pour ne pas se souffler la réponse à elle-même.

Nadège connaît le prix de tout, à quelques francs près, et elle ne le dit jamais fort. Le prix d’une heure de connexion, celui d’une page imprimée, plus élevé en couleur, ce qu’elle déconseille aux clients pressés. Le prix d’un forfait qui dure une semaine si on fait attention et trois jours si on regarde des vidéos. Ce qu’elle gagne en un soir, qui se mesure assez exactement en pages imprimées. À Bertoua, sa mère tient la caisse de la quincaillerie familiale ; Nadège a hérité de la manière de compter, pas de la caisse.

Vers 20 h, une étudiante entre avec un ordinateur encore dans sa housse de mousse. Nadège la reconnaît : même filière, même amphi, trois rangs devant. L’étudiante demande le code du réseau, s’installe au poste le plus proche de la prise, téléphone coincé contre l’épaule. « Non, dis à maman que le répétiteur peut passer jeudi. Mercredi j’ai TP. » Elle a besoin d’imprimer un compte rendu, douze pages, et elle demande si on peut relier. On peut.

Nadège ne la déteste pas. Elle n’en a pas le temps, et puis ce serait mal compter. Pendant que la relieuse chauffe, elle fait autre chose, qui lui ressemble davantage : le compte de ce qu’elle a. Un tabouret, un comptoir, trois soirs par semaine dans un endroit éclairé, avec de la connexion entre deux clients. Un paquet de cartes. Une voisine de chambre, à la mini-cité, avec qui elle partage un forfait. Une amie en deuxième année qui a gardé ses sujets d’examen. Sa mère le dimanche, au téléphone, qui ne connaît rien à la biologie et qui demande toujours si elle mange.

Et, depuis cet après-midi, un résultat. Les notes du rattrapage ont été affichées aux valves du département, ces panneaux grillagés où l’université épingle ce qu’elle a décidé. Une foule se pressait devant le grillage, des épaules, des téléphones levés au-dessus des têtes, un garçon qui lisait à voix haute les noms des autres sans qu’on le lui demande. Les feuilles étaient tenues par des punaises rouillées et portaient un tampon violet dans un coin. Nadège a attendu qu’un espace se libère, puis elle a lu la liste avec le doigt, comme on lit un relevé de compte. Elle a validé une partie des unités. Plus qu’en janvier ; moins qu’il n’en fallait. Elle a photographié la feuille avec son téléphone, puis elle est venue ouvrir le cybercafé, parce que c’était mercredi.

\asterisme

Certaines choses ne dépendent pas de toi, et elles pèsent. L’étudiante à l’ordinateur neuf n’a rien pris à Nadège. Elle a simplement commencé l’année avec des heures libres que Nadège passe derrière un comptoir, avec un répétiteur pour les soirs où le cours ne passe pas, avec une machine qui n’appartient à personne d’autre. Rien de tout cela n’apparaîtra sur le relevé de notes. Le relevé dira seulement qui a validé quoi.

Ce qui pèse mérite d’être nommé, sans chiffres que je n’ai pas et sans accuser personne : tu le vis déjà, et tu as droit à ce qu’on le dise à voix haute. L’argent, d’abord, qui achète des manuels, des photocopies lisibles, des annales, des répétitions. Le courant : à Garoua, la lampe solaire d’Aïcha charge sur le toit toute la journée, et quand le ciel s’est couvert, la leçon du soir est plus courte que prévu. La connexion, qui coûte, et qui manque précisément le soir où l’on en a besoin. Le temps, ensuite, distribué de travers : certains travaillent après les cours ; d’autres, garçons ou filles, gardent les petits frères pendant que la mère est au marché, puis font la cuisine pour toute la maison avant d’avoir le droit d’ouvrir un cahier. La distance : à Douala, Grâce traverse le pont sur le Wouri chaque matin, et les embouteillages lui mangent des heures qu’aucun bulletin ne lui rendra.

La santé compte aussi. Un palu la semaine d’une séquence, et la séquence n’attend pas. Des enseignants manquent, le plus souvent parce qu’une mutation n’a pas été remplacée ou qu’ils sont tombés malades à leur tour ; ceux qui sont là tiennent, devant des classes si pleines qu’on s’assoit à trois par banc et qu’on n’entend pas le fond de la salle. Quant aux concours, les places y sont rares, et l’on soupçonne parfois, à tort ou à raison, que la copie n’est pas la seule chose qui compte. Je n’ai aucune preuve à t’apporter là-dessus et je ne t’en inventerai pas. Je sais seulement que ce soupçon existe, et qu’il décourage plus sûrement qu’un mauvais sujet.

Rien de tout cela n’est neuf. Au début des années soixante, en France, deux sociologues, Pierre Bourdieu et Jean-Claude Passeron, ont regardé de près qui réussissait à l’université et pourquoi. Dans \emph{Les Héritiers}, publié en 1964, ils soutiennent une idée qui a fait date : l’école tend à prendre pour des qualités personnelles des avantages reçus en naissant, une manière de parler, des livres à la maison, une familiarité ancienne avec ce qu’elle attend, et à les récompenser comme s’il s’agissait de mérite. Leur travail a été beaucoup discuté depuis, et la France de cette époque n’est pas le Cameroun d’aujourd’hui. Mais le mécanisme qu’ils décrivent se reconnaît partout où l’on note des élèves qui ne sont pas partis du même endroit.

Pour toi, cela veut dire qu’une partie de ta note ne te mesure pas. Elle mesure la lampe, le trajet, le forfait, le répétiteur que tu n’as pas eu, la soirée où tu as gardé ta petite sœur. Une partie seulement. L’autre partie te mesure bien, et c’est d’elle que s’occupe le reste de ce livre. Mais tu as le droit de savoir que la ligne entre les deux existe, même si personne ne la trace sur le bulletin.

Cette ligne change la façon dont on entend le mot qui ouvre ce livre. Quand quelqu’un dit d’un élève qu’il est nul, il additionne tout ce qui figure sur la copie et il l’attribue à une seule personne. Il oublie de retrancher le trajet, la lampe, la semaine de fièvre. Le calcul inverse, qui retrancherait tout et ne laisserait rien à ton compte, serait la même erreur dans l’autre sens. Fais-le plutôt comme Nadège fait la caisse : ce qui est entré, ce qui est sorti, et ce qui n’était pas à toi.

Nadège, au comptoir, ne perd pas sa soirée à comparer sa vie à celle de l’étudiante à l’ordinateur ; cette comparaison-là ne lui apprendrait rien qu’elle ne sache déjà. Elle fait son inventaire, comme à la quincaillerie de Bertoua on compte les sacs de ciment avant de décider ce qu’on commande. On commande ce qui manque dans le magasin, et le stock du voisin n’entre pas dans le calcul.

\asterisme

Au comptoir, pendant que Nadège répond à sa carte sur la mitose, une autre partie de sa tête fait des additions qui ne la concernent pas, en apparence. Le loyer de la chambre, à payer avant le cinq. Le car pour Bertoua, à Pâques, si elle rentre. Le forfait, qui finit vendredi. Les photocopies du polycopié de chimie, qu’elle n’a pas encore faites parce qu’elle attend de voir si une camarade peut lui prêter le sien. Elle porte dans sa tête une caisse qui ne ferme jamais tout à fait.

Un économiste et un psychologue américains se sont intéressés à ce que le manque fait à l’esprit. Leurs travaux suggèrent que les soucis d’argent mobilisent l’attention, même quand on croit penser à autre chose, et que ce qu’ils prennent manque ensuite à la leçon. Certains de leurs résultats ont été discutés, et d’autres équipes ne les ont pas toujours retrouvés avec la même netteté. L’idée générale me semble pourtant juste, et surtout utile : quand Nadège relit trois fois la même carte sans la retenir, la caisse, dans sa tête, a peut-être simplement pris une partie du comptoir. Cela ne dit rien de ce qu’elle vaut en biologie.

Être épuisé par le manque n’est pas être paresseux. La confusion entre les deux se fait vite, et le plus souvent dans ta propre tête. Le soir où tu n’arrives à rien après une journée de trajet, de cours et de corvées, la voix qui te traite de fainéant se trompe d’adresse.

On a d’ailleurs inventé, pour les gens dans la situation de Nadège, un mot qu’on leur épingle volontiers sur la poitrine : la résilience. Je me méfie de ce mot quand il sert à admirer la charge plutôt qu’à se demander pourquoi elle pèse sur certaines épaules et pas sur d’autres. Nadège se passerait bien des applaudissements. Ce qui l’aiderait, c’est qu’on ne prenne pas sa fatigue pour un défaut de volonté, à commencer par elle-même.

Depuis février, elle a pris une habitude modeste qu’elle n’a lue nulle part. Avant d’ouvrir le paquet de cartes, elle écrit à la dernière page du cahier à spirale les additions du jour, le loyer, le car, le forfait, avec la date à laquelle elle s’en occupera. Puis elle referme la page. Ça ne paie pas le loyer. Ça lui rend un peu de comptoir. Je ne sais pas si cela marcherait pour toi ; il me semble qu’une inquiétude écrite avec une date a moins besoin d’être répétée qu’une inquiétude qu’on garde en tête pour ne pas l’oublier.

\asterisme

Voilà presque deux mille ans, un homme qui avait été esclave avant de devenir professeur de philosophie, Épictète, faisait commencer son enseignement par une distinction très simple. Un de ses élèves l’a placée en tête du petit livre qu’on appelle le \emph{Manuel}. Parmi les choses, disait-il en substance, certaines dépendent de nous et d’autres non : nos jugements, nos élans, nos désirs et nos refus sont de notre ressort ; notre corps, nos biens, notre réputation, les charges qu’on nous donne ou qu’on nous refuse ne le sont pas. Celui qui confond les deux, ajoutait-il, se rend malheureux pour rien.

C’est une idée solide, et un ancien esclave savait mieux que personne de quoi il parlait. La version honnête, plus que celle qu’on trouve parfois peinte sur les murs, y ajoute une précision : la frontière est floue, et elle bouge.

Le courant échappe à ta main ; l’heure à laquelle tu révises, un peu moins : si le courant part souvent à 21 h dans ton quartier, la leçon difficile peut passer avant. Le prix du forfait se décide ailleurs, mais le partager avec ta voisine de chambre, comme Nadège, se décide à deux. Tes parents n’ont pas pu acheter le manuel ; celui de la grande sœur d’un camarade, qui dort dans un carton depuis deux ans, dépend d’une question que tu peux poser demain. Le répétiteur, tu ne l’as pas. Le groupe de travail qui se réunit le samedi à la bibliothèque du lycée, ou de la faculté, tu peux peut-être le rejoindre, et s’il n’existe pas, le fonder à trois.

La frontière se déplace surtout quand on s’associe. Nadège ne peut rien contre le prix des annales ; une amie de deuxième année, qui les tient elle-même d’un ancien, les lui a passées. Personne ne fera ouvrir la bibliothèque plus tard pour toi, mais tu peux apprendre à quelles heures elle est calme, et y aller à ces heures-là. Un chargé de travaux dirigés n’aura pas plus de temps parce que tu le souhaites ; à trois, avec une question écrite, on lui prend dix minutes qu’il donne volontiers, parce qu’une question précise lui fait gagner du temps à lui aussi. Les grands frères, les grandes sœurs, les cousins qui ont passé l’examen l’an dernier sont des ressources que personne ne met sur la liste des fournitures. Ils n’en sont pas moins réels.

Ce soir-là, d’ailleurs, la frontière bouge sous les yeux de Nadège sans qu’elle l’ait cherché. Pendant que le compte rendu refroidit dans la relieuse, l’étudiante à l’ordinateur regarde le paquet de cartes posé près de la caisse. « Tu fais ça toi-même ? » Nadège dit que le carton est gratuit et que le feutre lui a coûté moins qu’une page en couleur. L’autre rit, puis demande, plus bas, si par hasard Nadège n’aurait pas les sujets de l’an dernier : son répétiteur fait des exercices, dit-elle, mais pas ceux de la faculté. Nadège les a, grâce à son amie de deuxième année. Elle propose d’abord de les photocopier sur place, au prix de la page. Puis elles trouvent mieux, un échange : les sujets contre les corrigés que le répétiteur rédige chaque jeudi.

Personne, dans cette scène, n’a tout : l’une a la machine et les heures, l’autre les sujets et la méthode des cartes. L’échange reste inégal, et l’étudiante rentrera ce soir dans une chambre où personne ne lui demande de fermer une grille à 22 h. Il a eu lieu quand même. Des échanges de ce genre existent même dans une situation inégale, et on les rate souvent par fierté ou par rancune.

Un manque se comble, disais-je au début de ce livre. La formule demande une précision : certains se comblent en travaillant seul, d’autres à plusieurs, et quelques-uns ne se comblent pas du tout à ton échelle. Savoir les distinguer, c’est déjà cesser de t’en vouloir pour les derniers.

Nadège tient les deux bouts à la fois, et c’est plus difficile qu’il n’y paraît. Elle sait que le comptoir, le car et le loyer pèsent sur ses notes, et elle ne s’en accuse pas. Elle sait aussi que la carte ratée trois fois retournera ce soir en tête du paquet, et que cela, personne ne le fera pour elle. Si tu ne gardes que la première moitié, tu finis par expliquer chaque mauvaise note par le forfait et le courant, et tu as souvent raison, et tu n’avances pas. Si tu ne gardes que la seconde, tu te reproches des heures de trajet comme s’il s’agissait d’un manque de caractère.

Rien de ce que tu feras dans cette marge ne garantit quoi que ce soit. L’effort laisse l’inégalité en place ; il déplace des probabilités, un peu. Une heure gagnée ici, une annale empruntée là, une question posée au lieu d’être gardée. Ce sont de petites marges. Mais un examen se joue souvent sur de petites marges, et ce sont les seules sur lesquelles tu aies la main.

\soustitre{Ce soir, essaie ceci}

Ouvre ton cahier de brouillon à une page neuve et dresse la carte de ce que tu as. Ce qui te manque, tu le connais par cœur, tu pourrais le réciter dans le noir ; cette page-là est réservée à l’autre liste.

Commence par les heures. Écris, jour par jour, les moments de la semaine où tu peux travailler, même vingt minutes. Le trajet en car compte, si tu peux y relire trois cartes. L’heure avant que la maison se réveille compte. La récréation, parfois. Note seulement les heures que tu as vraiment, quitte à ce que la liste soit courte.

Ensuite les endroits. Où, autour de toi, trouves-tu de la lumière et un peu de calme, et à quelle heure ? La bibliothèque du lycée, une salle vide après les cours, l’église ou la mosquée en dehors des offices si l’on t’y laisse t’asseoir, la cour d’un voisin, le comptoir d’une boutique quand le patron est d’accord. Écris le nom du lieu et l’heure à côté.

Puis les personnes, et d’abord celles qui savent quelque chose que tu ne sais pas. Ceux qui t’aiment forment une autre liste, même quand les noms se recoupent. Le camarade fort en physique. La cousine qui a eu le Probatoire l’an dernier. Le voisin qui enseigne au primaire. Le professeur qui reste dans sa salle dix minutes après la sonnerie. À côté de chaque nom, écris ce qu’il sait.

Enfin les objets. Un manuel, même ancien. Des annales, même incomplètes. Le cahier d’un ancien élève, qui dort quelque part dans une maison que tu connais. Une lampe. Un téléphone partagé.

Regarde ta carte. Elle sera peut-être plus pauvre que celle de l’étudiante à l’ordinateur neuf ; elle sera sans doute plus riche que tu ne le pensais en commençant. Puis choisis un seul levier dans cette carte, un seul, et utilise-le cette semaine. Demain, demande le cahier de l’ancien. Jeudi, va à la bibliothèque à l’heure calme. Samedi, pose ta question au camarade de physique. Une seule de ces choses, pas les trois. Dans sept jours, note à côté de ce levier ce qu’il t’a donné, même si c’est peu.

\asterisme

À 22 h, Nadège ferme. Elle éteint les postes un par un, du fond vers la porte, vérifie que l’imprimante n’a gardé aucune feuille dans le ventre, recopie dans le cahier à spirale les heures de la soirée. L’onduleur émet son dernier bip. Elle compte la caisse deux fois, la remet au patron, tire la grille. Ses cartes sont dans la poche de son pantalon, l’élastique autour.

Elle descend la rue vers la mini-cité. Il fait encore chaud ; les bars laissent sortir de la musique et de la lumière par morceaux. Au carrefour, sous le lampadaire, trois élèves lisent debout, des cahiers ouverts sur le capot d’une voiture garée. L’un tient une feuille grise qu’il tourne vers l’ampoule ; les deux autres se disputent à voix basse sur la fin d’un exercice. L’ampoule est faible, d’un jaune qui fatigue les yeux. La plus petite des trois a sorti son téléphone et en braque la lumière sur la page, le bras tendu, comme on tient un parapluie au-dessus de quelqu’un.

Nadège ralentit un peu, pas assez pour qu’ils la remarquent. Elle reconnaît le geste, la page tendue vers la lumière, la voix qui vérifie. Elle pense aux valves, à la photo dans son téléphone, aux unités qui restent. Elle pense aussi qu’elle appellera sa mère dimanche, et qu’elle lui dira le résultat en entier, la partie validée et l’autre.

Elle passe. Derrière elle, sous le lampadaire, les trois continuent, avec ce qu’ils ont.

\begin{pourcahier}{Pour le cahier}
Trace deux colonnes et écris-les en phrases, pas en mots isolés. À gauche : « Ce qui ne dépend pas de moi cette semaine ». À droite : « Ce qui dépend un peu de moi cette semaine ». La colonne de droite doit contenir au moins une chose que tu feras demain, avec l’heure.
\end{pourcahier}


%% 14-ce-qu-on-attend-de-toi.md
\chapitre{12}{Ce qu’on attend de toi}

À Ebolowa, la fiche d’orientation est arrivée avec les pluies de mai. Ce soir, elle est posée sur la table de la cuisine des Essomba, la seule pièce éclairée de la maison, entre le seau d’arachides à trier pour le marché du lendemain et la boîte de sucre que Bella a encore mal refermée. Une seule feuille, imprimée recto verso, pleine de cases. Il faut la rendre au lycée avant la fin du mois, signée par le parent ou le tuteur, et le parent, ce soir, c’est Mama Odile : le père est au port de Kribi, il rentrera peut-être samedi, peut-être pas, et la fiche n’attendra pas les week-ends du port.

Rodrigue a déjà rempli sa partie. Dans la case des vœux, de l’écriture appliquée qu’il réserve aux jours où il fait attention, il a écrit : lycée technique, série électronique. Il l’a d’abord écrit au crayon, puis repassé au stylo, puis il a soufflé sur l’encre comme si elle risquait de s’envoler. Le petit tournevis est dans la poche de sa chemise, comme toujours. Il le touche à travers le tissu sans le sortir.

Mama Odile lit la fiche avec ses lunettes de marché, celles qu’elle met pour compter la monnaie. Elle lit lentement, deux fois, comme on relit un papier de la tontine quand une somme ne tombe pas juste. Puis elle pose la feuille, la lisse du plat de la main, et dit qu’elle voulait autre chose pour lui. Un bureau. Elle ne précise pas lequel ; elle n’en a pas besoin, elle le voit. Une chaise, un ventilateur au plafond, une chemise qui reste propre jusqu’au soir, un salaire qui tombe à la fin du mois, que les clients soient venus ou non. Un métier où l’on ne se salit pas les mains. Un métier où l’on ne dépend pas d’une boutique au carrefour.

Rodrigue explique. Il avait préparé ses phrases dans la journée, sur le banc de la cour, et elles sortent maintenant dans le désordre : que l’électronique ne se résume pas à dévisser des coques, qu’on y fait des calculs, des schémas, des lois avec des noms de savants ; que même les gens des bureaux, quand leur téléphone est gâté, viennent le déposer au carrefour ; qu’il aime comprendre pourquoi une chose ne marche pas. Sur « aime », il s’arrête, un peu gêné, parce qu’on n’emploie pas souvent ce verbe dans cette cuisine pour autre chose que le ndolé.

C’est à ce moment que la phrase sort. « Tu veux finir comme qui ? »

Mama Odile la regrette dès le deuxième mot. Mais une phrase lancée va jusqu’au bout, comme une pièce de monnaie qui roule sous le banc, et celle-ci arrive exactement là où elle devait arriver : au carrefour, à la boutique de tonton Fabrice, au comptoir où Rodrigue passe ses samedis à remonter des pannes. Le nom de l’oncle n’est pas prononcé. Il flotte déjà au-dessus de la table.

Pour une fois, Rodrigue ne rit pas le premier. Il regarde la fiche. Au bout de la table, la tête posée sur ses bras croisés, la petite cousine Bella a les yeux fermés et cette respiration trop régulière de ceux qui ne dorment pas. Elle est en CM2. Elle a déjà compris qu’on apprend davantage, dans une maison, en faisant semblant de dormir qu’en posant des questions.

\asterisme

Plus tôt ce jour-là, en fin d’après-midi, à Ngaoundéré, la séance de travaux pratiques des Premières F3 déborde sur l’heure suivante, dans l’atelier du lycée technique, et personne ne s’en plaint. Ibrahim Saïdou tient un mètre ruban d’une main, une gaine électrique de l’autre. Il mesure. Il marque d’un trait de crayon. Puis il mesure encore, depuis le même bout, avec la même lenteur, et seulement alors il coupe. Un camarade lui fait remarquer qu’il perd du temps. Ibrahim répond sans lever la tête qu’une gaine coupée trop court, on ne la recolle pas. Le professeur d’atelier passe derrière lui, regarde le trait de crayon, ne dit rien, ce qui, dans cet atelier-là, compte comme un compliment.

À Douala, le dimanche suivant, chez les Ewane, la machine à coudre s’est tue pour le repas. La tante du kiosque est venue, une voisine aussi, et la conversation arrive, comme chaque fois, sur l’infirmerie. Grâce ferait une bonne infirmière, dit la tante : elle est patiente, elle a la main douce, et des malades, il y en aura toujours. Tout le monde approuve, Grâce aussi, parce que hocher la tête coûte moins cher que d’expliquer. Elle a seize ans, elle chante alto dans une chorale, elle aime l’espagnol au point de se réciter des conjugaisons dans le taxi qui traverse le pont, et elle ne sait pas ce qu’elle veut faire. Elle ne sait même pas si elle devrait déjà le savoir. Autour de la table, personne ne lui pose la question sous cette forme-là. Après le repas, elle aide sa mère à reprendre un ourlet. Sa mère, des épingles entre les lèvres, ne dit rien de l’infirmerie ; elle demande seulement, au moment de couper le fil, si la robe tombe droit. Grâce répond que oui. C’est la seule question de la journée à laquelle elle soit sûre de savoir répondre.

Ces maisons ne se connaissent pas. Elles se ressemblent pourtant sur un point : quelqu’un, autour d’une table, a une idée de ce que tu devrais devenir, et cette idée est souvent arrivée avant la tienne.

\asterisme

La phrase de Mama Odile fait mal, et il serait facile de s’arrêter à la douleur. Une mère qui demande à son fils « tu veux finir comme qui ? » a l’air de mépriser un métier et, avec lui, un carrefour entier. Regarde-la mieux. Elle vend du plantain et des arachides au marché depuis des années. Elle sait ce que c’est qu’une journée sans clients, ou une pluie qui vide les allées dès dix heures. Quand elle dit « un bureau », elle pense à un salaire que la pluie ne peut pas emporter.

Les attentes des parents mêlent presque toujours plusieurs choses, si bien serrées qu’on ne les distingue plus : de l’amour, qui veut pour toi une vie moins dure que la sienne ; de la peur, qui connaît les mauvaises années et voudrait t’en protéger en choisissant à ta place ; un calcul, enfin, souvent raisonnable dans ses données de départ. Ton père, ta mère ont vu des gens s’en sortir et d’autres tomber, ils en ont tiré des règles, et ces règles datent de leur jeunesse à eux. Dans les histoires qu’on racontait autour de Mama Odile quand elle avait ton âge, celui qui obtenait un diplôme entrait dans un bureau et n’en sortait plus, et celui qui travaillait de ses mains restait au soleil toute sa vie. Ces histoires étaient vraies pour certains, à une certaine époque, et elles ont continué d’être racontées bien après que le monde a changé autour d’elles. Tout cela tient ensemble dans une seule phrase prononcée trop vite, et c’est pour cela qu’une discussion d’orientation, à la maison, ressemble si rarement à une discussion.

Ce que personne ne dit à voix haute pèse aussi. Les frais de scolarité, les cahiers, la tenue, les fournitures de chaque rentrée, payés année après année avec l’argent des arachides. Une attente, c’est parfois une dette dont on ne parle jamais, et qu’on porte des deux côtés de la table. Le parent a peur que l’investissement soit perdu ; l’enfant a peur de ne pas le rembourser. Chacun se tait sur sa propre peur et entend surtout celle de l’autre comme un reproche.

Une phrase dure est souvent une peur qui n’a pas trouvé d’autres habits. Cela ne la rend pas juste, et tu n’es pas obligé de l’avaler. Mais cela t’aide à entendre ce qu’elle dit vraiment, sous le ton. Celle de Mama Odile parle d’elle bien plus que de tonton Fabrice, dont elle sait qu’il n’a pas raté sa vie, puisque la moitié du quartier lui confie ses téléphones. Elle a peur que son fils dépende un jour, comme elle, des clients qui passent ou ne passent pas. Si Rodrigue répond au ton, il se défend. S’il répond à la peur, il peut parler.

Ce soir-là, dans la cuisine, il ne le sait pas encore. Il sait seulement que sa mère a baissé les yeux juste après avoir parlé, et qu’elle n’a pas repris la fiche. Elle trie maintenant les arachides, une à une, en écartant les mauvaises du bout de l’ongle, avec cette précision qu’elle met dans les comptes de la tontine. Lui garde la main dans sa poche. Ce silence-là ressemble à deux personnes qui font un calcul chacune de son côté, sans savoir encore qu’il s’agit du même.

\asterisme

Le lycée technique n’est pas l’endroit où l’on range ceux qui n’ont pas pu faire autrement, quoi qu’on en dise dans les familles et dans les cours de récréation. Et les métiers des mains n’ont rien d’une consolation offerte à ceux que la dissertation a refusés.

Quand un voisin l’appelle parce que tout disjoncte dès qu’on branche le fer à repasser, Ibrahim ne remplace pas les fils au hasard. Il écoute la description de la panne, il formule une hypothèse, il la teste, l’écarte, en essaie une autre. Il remonte du symptôme vers la cause, comme Rodrigue au comptoir de la boutique. Les cours de mathématiques appellent cela raisonner. Ici, une erreur de raisonnement peut brûler un câble, ou une main : si Ibrahim mesure deux fois, c’est qu’il a compris qu’un trait de crayon mal placé ne se rature pas.

Les filières techniques et professionnelles sont exigeantes. On y fait des mathématiques et de la physique, du dessin technique où un millimètre compte, de la technologie, des normes de sécurité qu’on n’a pas le droit d’à peu près savoir. On y écrit aussi : un technicien rédige des rapports, des devis, des comptes rendus d’intervention, et un client comprend mieux une panne bien expliquée. Ibrahim, qui se dit nul en français, toujours avec le même sourire, remplit chaque semaine une fiche où il doit décrire en phrases complètes ce qu’il a monté, dans quel ordre et pourquoi. Il la trouve plus difficile que le montage. Il la remplit quand même, en relisant deux fois, parce que dans l’atelier on ne rend rien sans l’avoir vérifié, pas même des phrases. La dernière lui est revenue avec un seul mot du professeur dans la marge : « Clair. » La couturière qui prend des mesures sur un corps vivant, le mécanicien qui entend au bruit d’un moteur ce qui ne va pas, le maçon qui monte un mur d’aplomb sans laser : tous font, chaque jour, un travail d’intelligence qu’on a pris l’habitude de ne pas appeler ainsi parce qu’il laisse de la poussière sur les vêtements.

Le technique ne mène pas à coup sûr à un emploi, et aucune autre voie non plus ; je me méfie de ceux qui promettent le contraire. Des diplômés de l’enseignement général cherchent du travail, des techniciens aussi, et le bureau que Mama Odile imagine n’attend nulle part, tout prêt, avec sa chaise et son ventilateur. Ce qu’on peut dire honnêtement est plus modeste. Aucune voie n’est une garantie, aucune n’est une condamnation, et ce qu’on a appris sérieusement, dans quelque filière que ce soit, on l’emporte avec soi quand les plans changent.

La plupart des décisions d’orientation se prennent sur des images : une blouse, un bureau, une boutique, une photo au mur. Une image ne dit rien d’une journée. La tante de Grâce parle de l’infirmerie depuis des années ; il n’est pas sûr que quelqu’un, dans la famille, ait jamais demandé à une infirmière ce qu’elle fait à 3 h du matin, ou ce qu’elle aurait aimé savoir avant de commencer. Pour choisir, ou pour discuter un choix, il faut parler à des gens qui exercent vraiment le métier, et pas seulement à ceux qui en parlent. Il faut aller voir : demander à visiter un lycée technique un jour de travaux pratiques, passer un samedi dans une boutique ou un garage, demander poliment, avec un adulte, à voir de près un service hospitalier, regarder des mains au travail plutôt que des photos. Le conseiller d’orientation de ton établissement existe pour cela aussi ; on oublie souvent qu’on a le droit d’aller frapper à sa porte avant que la fiche soit remplie.

\asterisme

Deux psychologues américains, Edward Deci et Richard Ryan, ont développé, à partir des années 1970, ce qu’ils ont appelé la théorie de l’autodétermination. Leurs travaux, et ceux de nombreux chercheurs qui s’en sont inspirés, suggèrent qu’on s’engage plus durablement dans une activité quand on s’y sent compétent, quand on y garde une part de choix, et quand on se sent relié aux autres en la faisant. Les détails en restent débattus ; l’idée générale est assez solide pour qu’on s’en serve.

Elle ne te donne pas pour autant le droit de faire ce que tu veux, quand tu veux, contre l’avis de tout le monde. Le troisième besoin, celui qu’on oublie, est le lien. On travaille mieux quand on n’est pas seul, quand ceux qui comptent pour nous savent ce qu’on fait et le respectent, même sans tout comprendre. Une orientation arrachée contre sa famille prive de ce lien autant qu’une orientation imposée prive du choix. Si ces recherches t’apportent quelque chose, c’est une raison de parler avec tes parents plutôt que contre eux.

On entend souvent dire, aux jeunes gens, qu’il faut trouver sa passion et la suivre, et que le reste viendra. Le conseil est aimable ; il laisse beaucoup de monde de côté. À seize ans, on peut n’avoir aucune passion, sans défaut de fabrication pour autant. On peut en avoir une qui ne nourrit personne, et le savoir. Le goût de Rodrigue pour les pannes s’est formé samedi après samedi, à mesure qu’il en trouvait, et c’est souvent ainsi que le goût vient : avec la compétence, à mesure qu’elle grandit. Grâce ne sait pas ce qu’elle veut faire, et rien n’est plus ordinaire à son âge. Cet état demande seulement qu’on aille voir, qu’on essaie, qu’on se trompe un peu.

La conversation elle-même, celle qui a mal commencé dans la cuisine d’Ebolowa et qui commence mal dans beaucoup de cuisines, a plus de chances d’aboutir si tu écoutes d’abord la peur, avant de défendre ton envie. Elle avance mieux avec un projet concret qu’avec un désir : quel établissement, quelle série, ce qu’on y apprend, ce que font ceux qui en sortent, ce que tu ferais si cela ne marchait pas. Elle demande enfin d’accepter de ne pas tout obtenir le premier soir. Propose aussi, si c’est possible, que vous alliez voir ensemble. Un parent qui a vu de ses yeux une salle de travaux pratiques, des élèves en blouse de travail penchés sur un circuit, un professeur qui explique au tableau une loi d’électricité avec la même craie qu’en mathématiques, ne parle plus tout à fait du même lycée que celui de ses inquiétudes. Il ne sera pas forcément convaincu. Mais vous parlerez enfin de la même chose. Un parent qui dit « je vais réfléchir » demande du temps pour changer d’image, et une image qu’on porte depuis vingt ans ne se décroche pas en une soirée.

\asterisme

Il arrive que ce qu’on attend de toi ne soit plus une attente. Quand on te menace, quand on te frappe, quand on veut te marier contre ta volonté ou te retirer de l’école pour cela, il ne s’agit plus d’une discussion d’orientation, et ce livre n’a pas de méthode à te proposer. Il faut en parler, sans attendre, à un adulte de confiance en dehors de la famille : un enseignant, le conseiller d’orientation, le surveillant général, un responsable religieux ou communautaire en qui tu as confiance, un travailleur social, ou un soignant dans un centre de santé. Tu n’as pas à régler cela seul. Et tu n’as pas à attendre que ce soit pire pour avoir le droit d’en parler.

\soustitre{Ce soir, essaie ceci}

Choisis un adulte qui exerce un métier qui t’intéresse, ou le métier que tes parents imaginent pour toi. L’idéal serait d’en trouver un de chaque. Cherche autour de toi avant de chercher loin : un voisin, un parent éloigné, l’infirmier du centre de santé, l’électricien qui est venu refaire un branchement l’an dernier, la couturière du quartier, le comptable d’une entreprise où travaille quelqu’un de ta famille. Écris son nom dans ton cahier de brouillon, et en dessous le métier, tel que tu l’imagines, en une phrase.

Demain, va le voir et demande-lui dix minutes. Pas une heure ni un après-midi : dix minutes, c’est une demande qu’on accorde souvent, et c’est assez pour apprendre beaucoup. Pose-lui quatre questions, à peu près dans ces termes. Qu’est-ce que vous faites vraiment, dans une journée ordinaire, du matin au soir ? Qu’est-ce que vous avez dû apprendre pour le faire, à l’école et après l’école ? Qu’est-ce qui est dur dans ce métier, que les gens ne voient pas ? Qu’est-ce que vous n’aviez pas imaginé avant de commencer ? Pendant qu’il répond, n’essaie pas de défendre ton idée ni celle de tes parents. Écoute. Demande un exemple si la réponse reste vague.

Le soir même, note ses réponses dans ton cahier, avec ses mots à lui autant que possible. Dans trois jours, tu aurais gardé l’impression générale et perdu les détails, et ce sont les détails qui t’apprendront quelque chose. Relis ensuite la phrase que tu avais écrite la veille, celle du métier tel que tu l’imaginais. Souligne ce qui était juste. Barre ce qui ne l’était pas, sans le rendre illisible.

Enfin, avant la fin de la semaine, prépare une seule phrase pour tes parents. Une phrase d’ouverture, pas un discours. Elle doit commencer par leur inquiétude ; la tienne viendra ensuite. Quelque chose comme : « Je sais que tu as peur que je dépende un jour des clients du carrefour. J’ai demandé à quelqu’un qui fait ce métier comment ça se passe vraiment, et je voudrais te dire ce qu’il m’a répondu. » Écris-la, lis-la à voix basse, corrige ce qui sonne faux. Tu n’es pas obligé de la prononcer ce week-end. Mais le jour où la conversation reviendra, et elle reviendra, tu ne commenceras pas par te défendre.

\asterisme

Dans la cuisine d’Ebolowa, les arachides sont triées. Mama Odile s’essuie les doigts sur son pagne et tend à son fils, sans un mot, son vieux téléphone, celui dont la touche du milieu reste enfoncée depuis Pâques. Elle ne pose aucune question. Rodrigue sort le tournevis de sa poche. Il ouvre la coque sur la table, sous l’ampoule, et pose chaque minuscule vis sur une fleur de la nappe en plastique, une vis par fleur, pour n’en perdre aucune. Sa mère le regarde faire. Elle regarde surtout ses mains, qui ne tremblent pas, et l’ordre dans lequel il range les pièces, de gauche à droite, comme on aligne les cotisations d’un mois.

Il nettoie le contact, remonte, revisse dans l’ordre inverse. La touche du milieu se relève avec un petit clic. Mama Odile appuie dessus trois fois, pour vérifier, comme elle recompte une somme qu’elle sait déjà juste.

Elle ne signe pas la fiche ce soir-là. Elle dit : « On va aller voir ce lycée. » Puis elle plie la feuille en deux et la range sous la nappe en plastique, avec les papiers importants, ceux qu’on ne perd pas. Au bout de la table, Bella, qui ne dort pas, compte les fleurs de la nappe.

\begin{pourcahier}{Pour le cahier}
Écris ce que ta famille attend de toi, en une phrase. Puis ce que tu attends de toi, en une phrase. Puis ce que tu ne sais pas encore. La troisième phrase a le droit d’être la plus longue.
\end{pourcahier}


%% 15-la-lampe.md
\chapitre{13}{La lampe}

Le lampadaire du carrefour découpe dans la nuit un rond de goudron à peine plus grand qu’une salle de classe. Au-delà, tout est rendu au noir : les étals bâchés du marché, les motos rangées en épi, une boutique dont le rideau de fer est baissé aux trois quarts. Dans le rond, une voiture que personne n’a déplacée depuis l’après-midi et, sur son capot encore tiède, trois cahiers ouverts. Il est un peu plus de 22 h. Autour de l’ampoule, les insectes tournent avec une obstination qu’on aimerait avoir pour la physique.

On est déjà passé par ici. Une étudiante, en rentrant du cybercafé, a vu ces trois élèves et a continué de descendre la rue ; elle avait ses propres cartes à réviser et la fatigue d’une soirée debout derrière un comptoir. Elle a eu raison. Cette fois, pourtant, on s’arrête. Les trois ne se ressemblent pas, et ils n’ont pas de prénom : on en trouve sous beaucoup de lampadaires, et leurs prénoms changent tous les soirs.

Le premier tient un polycopié. C’est la photocopie d’une photocopie d’une photocopie, et chaque passage dans la machine a pris quelque chose. Les lettres sont devenues grises, avec des bords qui bavent, et le schéma du milieu, qui devait représenter un montage électrique, se réduit à une pile, deux traits et une flèche qui ne mène plus nulle part. La légende a disparu dans le pli. Il incline la feuille vers l’ampoule, puis dans l’autre sens, comme si l’angle allait rendre ce que la machine a mangé. Il a payé ces pages. Il ne peut pas les lire.

La deuxième cherche sur son téléphone. Elle a tapé le titre de la leçon exactement comme le professeur l’a écrit au tableau, avec ses majuscules et son numéro. Le téléphone lui a répondu avec générosité : des cours, beaucoup de cours, bien présentés, parfois excellents. Mais le premier suit le programme d’un autre pays et commence par une notion qu’elle n’a jamais vue ; le suivant résout l’exercice avec une méthode que son professeur n’utilise pas et qu’il risque de ne pas accepter sur une copie. Elle fait défiler l’écran avec le pouce. En haut, l’icône de la batterie est passée au rouge, et chaque page chargée mord un peu plus dans le forfait de la semaine.

Le troisième n’a besoin de rien de tout cela, en apparence. L’explication est dans son cahier, copiée du tableau ligne par ligne, propre, soulignée à la règle. Il l’a relue dix fois. Chaque ligne, prise seule, est juste, et il pourrait la réciter les yeux fermés. Ce qui lui échappe, c’est le passage de la troisième à la quatrième : ce qui s’est produit dans la tête du professeur pendant la seconde où la craie a quitté une ligne pour commencer l’autre. Ce passage-là ne s’est jamais écrit au tableau. Il est resté dans la salle, le matin, et ce matin-là le troisième regardait peut-être par la fenêtre.

« C’est pas ça qu’il a fait », dit la fille en tournant l’écran vers les autres. Le garçon au polycopié se penche, plisse les yeux, renonce. « Ici même on voit rien. » Elle allume la lampe torche du téléphone et la tient au-dessus du cahier du troisième, parce que l’ampoule ne suffit plus pour les petites lettres. Trois têtes penchées dans un cercle de lumière blanche. Un moment, plus un mot.

Le travail se fait à cette heure-là, et c’est justement l’heure où personne ne peut l’accompagner. Le professeur est chez lui, à corriger les copies d’une autre classe, ou endormi, et il a le droit de l’être. Le grand frère qui savait la physique travaille dans une autre ville. Le répétiteur coûte ce qu’il coûte. Les parents sont à quelques rues, ils veilleraient volontiers, mais ils ne connaissent pas cette leçon-là, et ils le savent mieux que personne. Au carrefour, ce soir, il n’y a personne pour expliquer autrement.

\asterisme

Je raconte ce carrefour comme je raconte les autres scènes de ce livre : il est composé, et je ne prétends pas m’y être tenu. Mais aucun de ses morceaux n’est rare. Ce qui m’arrête, dans cette image, c’est la place vide à côté d’eux. Ces trois-là font exactement ce que les chapitres précédents demandaient : ils travaillent le soir, ensemble, avec ce qu’ils ont. Personne, en les regardant, ne pourrait soutenir qu’ils sont nuls. Il leur manque une page lisible, un cours qui corresponde à leur classe, une deuxième explication : des manques précis, dont aucun n’est une sentence. C’est à cette scène-là, et à toutes celles qui lui ressemblent, que j’ai voulu répondre, et c’est de là qu’est venue l’idée d’OnBuch, une application que je construis pour les élèves et les étudiants.

Une idée de ce genre ne vaut rien tant qu’on ne l’a pas mise au travail. Avant de construire quoi que ce soit, il fallait regarder le manque de près, comme on regarde une copie avant d’y mettre une note : ce qui manque exactement, à qui, à quelle heure. « Il n’y a pas de cours » ne veut presque rien dire. Le manque a plusieurs formes, et chacune appelle une réponse différente.

\asterisme

La plus simple à nommer est le cours introuvable. La leçon existe : elle est au programme, le professeur l’a faite, la classe l’a copiée. Mais tu étais absent ce jour-là, malade ou retenu ailleurs, ou c’est le professeur qui a manqué plusieurs semaines et qui a dû courir ensuite pour rattraper. Il reste le cahier d’un camarade, à l’écriture qu’on déchiffre mal, avec des trous aux endroits où lui-même a décroché. Il reste aussi le manuel, quand on l’a et quand il suit l’ordre de la classe. Quant au téléphone, il offre tout, sauf ta leçon.

Le polycopié illisible, lui, perd ses lettres dans un ordre qui mérite qu’on s’y attarde. À chaque reproduction, ce qui s’efface en premier, ce sont les détails fins : les indices, les exposants, les flèches, les pointillés d’un schéma, le signe moins devant une parenthèse. Autrement dit, ce qui porte le sens. Un texte grisé reste à peu près lisible ; un schéma grisé devient un dessin d’enfant. Un exposant effacé change un résultat sans prévenir personne, et l’élève qui révise sur cette page apprend parfois une erreur que personne n’a commise.

L’explication qui ne passe pas est le manque le plus difficile à voir, parce que tout semble en place. Le cours est là, complet, lisible. Seulement, une explication est fabriquée par une tête pour des têtes qui lui ressemblent un peu, et il en reste toujours quelques-unes pour lesquelles elle ne fonctionne pas, sans que ce soit la faute de qui que ce soit. Les bons professeurs le savent. Ils gardent une deuxième explication en réserve, un dessin, un exemple pris au marché. Encore faut-il avoir le temps de la sortir, devant une classe de plusieurs dizaines d’élèves et une sonnerie qui n’attend pas.

Par-dessus le reste vient la solitude du soir. Les autres manques seraient légers si quelqu’un était là, à 21 h, pour dire « montre-moi où tu bloques ». La journée d’école est pleine de gens ; la soirée de travail est vide. C’est le soir que la marche manquante apparaît, sous un lampadaire ou sous l’ampoule de la cuisine, et c’est le soir qu’on se retrouve seul devant elle, à se demander si le problème vient de la marche ou de soi.

Prends le cahier d’une élève absente trois jours, la semaine d’une leçon décisive : elle recopie celui d’une voisine et trouve, au milieu de la page, une ligne qui s’arrête sur « etc. », à l’endroit exact où il lui aurait fallu la suite. Elle ne se dit pas qu’il manque une ligne. Elle se dit qu’elle ne comprend pas.

Je n’en fais pas un réquisitoire. Les professeurs font leurs cours, souvent très bien, avec ce que la journée leur laisse ; les parents payent des photocopies qu’ils ne sauraient pas expliquer ; les camarades prêtent leurs cahiers sans compter. Le manque n’est la faute de personne en particulier. Il est simplement là, tous les soirs, à la même heure, et il prend le plus souvent la forme d’un élève qui conclut qu’il est nul alors qu’il lui manque une page.

Au carrefour, le garçon au polycopié a fini par demander au troisième de lui expliquer le passage. Le troisième essaie. Il récite sa ligne trois, puis sa ligne quatre, avec conviction, et entre les deux il glisse une phrase de sa fabrication qui a l’allure d’une explication. La fille fronce les sourcils : sur son écran, le cours venu d’ailleurs dit autre chose. Ils se disputent un moment, à voix basse, chacun sûr de sa lumière. Aucun des trois ne le sait encore, mais ils viennent de trouver la question qu’il faudra poser demain.

\asterisme

Quand on veut répondre à ce manque, la première décision n’a rien de spectaculaire, et pourtant elle commande toutes les autres : sur quel programme travailler. Le téléphone de la fille lui a montré des cours justes qui n’étaient pas les siens. Un cours peut être exact et ne servir presque à rien à un élève donné, s’il ne suit pas l’ordre de sa classe ou s’il prépare un autre examen. J’ai donc choisi de suivre le programme officiel du MINESEC, celui à partir duquel ton professeur construit ses séquences. Classe par classe, matière par matière, du collège à la terminale, dans l’enseignement général comme dans le technique. Le but est modeste et précis : que ce que tu trouves le soir corresponde à ce que tu as fait le matin.

Ce choix a ses limites, et je les connais. Un programme dit ce qu’il faut enseigner ; il dit peu de la manière. Ton professeur peut traiter les leçons dans un autre ordre, préférer une notation qui lui est chère, insister là où le programme passe vite. Les programmes eux-mêmes changent, et un cours écrit sur l’ancien vieillit sans faire de bruit. Quand ce que tu lis le soir ne ressemble pas à ce que tu as entendu le matin, c’est ton professeur qui fait foi pour ta classe, puisque c’est lui qui pose le devoir.

Concrètement, OnBuch propose des cours structurés leçon par leçon, des fiches de travaux dirigés à trois niveaux, des corrigés et des fiches méthode. Je le dis en une phrase, parce qu’il n’en faut pas davantage : ce sont les objets ordinaires du travail scolaire, ceux qu’un bon manuel ou un bon professeur fournit aussi. Tout cela est relu, et il y reste sans doute des erreurs que je n’ai pas encore trouvées. Comment ce travail se fait, avec quels outils, et ce qu’il coûte de le vérifier, c’est l’affaire du chapitre suivant.

\asterisme

J’ai appelé ce chapitre « La lampe » parce que je n’ai pas trouvé d’image plus juste pour dire ce qu’un tel outil peut faire, et surtout ce qu’il ne peut pas faire. Une lampe ne monte pas l’escalier à ta place. Elle éclaire la marche suivante, celle où tu vas poser le pied, et un peu celle d’après ; le reste demeure dans l’ombre, ce qui n’est pas grave, puisqu’on ne monte jamais qu’une marche à la fois. Une lampe ne remplace pas le soleil : la classe, le professeur qui voit ton visage au moment où tu décroches, le camarade qui sait où tu bloques parce qu’il y a bloqué la semaine dernière. Les jambes, elles aussi, restent les tiennes.

Un outil sur téléphone suppose un téléphone. Il suppose de la batterie, donc du courant pour la recharger, et souvent de la connexion, donc un forfait. Or les élèves qui manquent le plus de cours peuvent manquer aussi de cela. La fille du carrefour avait une batterie dans le rouge ; le garçon au polycopié n’avait peut-être pas de téléphone du tout. Il y a là une ironie que je ne peux pas effacer, et qui ne m’arrange pas : la lampe que je construis a besoin d’une prise, et le manque qu’elle voudrait éclairer se trouve souvent là où les prises sont rares. Je n’ai pas de réponse élégante à cela. J’ai seulement celle-ci : je préfère le dire que le taire, et je préfère que tu le saches avant de compter sur elle pour quoi que ce soit.

Sous le lampadaire, un seul téléphone éclairait trois cahiers. On peut lire une page à plusieurs sur un même écran ; on peut recopier une méthode dans son cahier de brouillon pour la retrouver le jour où la batterie lâche. Ce sont de petits arrangements, et ils ne règlent rien du fond. Ils disent seulement qu’une lampe se partage mieux qu’on ne le croit.

Alors à quoi bon, si ce n’est pas un miracle ? À donner une chance honnête, et j’entends par là quelque chose de précis. Que l’élève seul à 21 h trouve une page lisible, qui parle de sa leçon dans les mots de sa classe, avec un exercice à sa hauteur et un corrigé pour vérifier, ensuite, ce qu’il a cherché. Rien de plus. Une chance honnête ne garantit aucune note. Elle retire un des manques qui faisaient croire à l’élève qu’il était nul, et elle lui laisse tout le reste, qui est considérable : comprendre, s’exercer, se tromper, recommencer le lendemain.

Cette lampe, je la range parmi les autres, et pas au premier rang : ton professeur, qui te voit et connaît ta classe, compte davantage, et un camarade patient compte autant. Si elle peut t’éclairer, elle est là ; si elle ne le peut pas, le reste de ce livre tient sans elle.

\soustitre{Ce soir, essaie ceci}

Prends la leçon qui te résiste le plus cette semaine et choisis pour elle deux lampes, pas davantage. La carte de ce que tu as, si tu l’as dressée, en contient sans doute plus qu’il n’en faut : ton manuel ou celui d’un voisin de banc, ton professeur aux moments où l’on peut l’attraper, un camarade qui a compris, un grand frère ou une grande sœur qui a fait la même classe, les annales, le cahier d’un ancien, une application si ton téléphone le permet. Prends-les de nature différente si possible : un texte et une personne, ou deux livres qui ne viennent pas du même endroit. Écris leurs noms en haut d’une page de ton cahier de brouillon, et la date.

Demain, lis ou écoute ce que la première dit de la leçon, et résume-le en quelques lignes, avec tes mots à toi. Si c’est une personne, demande-lui de t’expliquer un seul passage, celui qui bloque, et note ce qu’elle a dit dès que tu l’as quittée. Après-demain, fais de même avec la seconde. Puis pose les deux résumés côte à côte et compare-les, ligne par ligne, comme on compare deux copies.

Si les deux lampes disent la même chose, tu peux avancer : tu tiens sur une base plus solide que si tu n’en avais consulté qu’une. Si elles disent la même chose avec des mots différents, tant mieux, car tu viens de voir la même idée sous deux angles, et c’est souvent ainsi qu’une explication finit par passer. Si elles se contredisent, ne choisis pas au hasard, et surtout pas la plus confortable. Tu as trouvé ta question pour le professeur. Écris-la en nommant l’endroit exact où les deux lampes divergent : « Le manuel écrit ceci à la deuxième étape, mon camarade fait cela ; lequel vaut pour notre classe ? » Pose-la à la fin du prochain cours, ou dépose-la sur son bureau en sortant.

Note la date dans la marge. Dans trois jours, sans rien relire, refais un exercice de cette leçon. Tu sauras alors si les lampes ont éclairé quelque chose en toi, ou seulement ta page.

\asterisme

Au carrefour, ils ont fini par se taire. Le troisième a écrit au crayon, en bas de sa page, une phrase qui se termine par un point d’interrogation, et la fille a éteint la lampe torche pour sauver ce qui reste de batterie. Le lampadaire grésille. Sa lumière tremble, vire à l’orange, revient, hésite encore, puis s’éteint d’un coup, et le rond de goudron rejoint le reste de la nuit. Les trois ne rentrent pas. Ils ramassent leurs cahiers, marchent jusqu’au lampadaire suivant, vingt mètres plus loin, devant la pharmacie fermée, et les rouvrent sur le muret.

\begin{pourcahier}{Pour le cahier}
Écris le nom de la leçon qui te résiste le plus. Sous elle, écris les deux lampes que tu vas utiliser, en les nommant précisément : pas « internet » ni « quelqu’un », mais quel manuel, quelle personne. Puis la date où tu les compareras. Le jour venu, ajoute une ligne : ce qu’elles disent toutes les deux.
\end{pourcahier}


%% 17-lettre-a-aicha.md
\chapitre{14}{Lettre à Aïcha}

Aïcha,

Je n’étais pas à Garoua ce matin-là, et je n’ai aucune envie de faire semblant d’y avoir été. Je sais seulement comment ces matins se passent, parce qu’ils se ressemblent d’une ville à l’autre, et je te demande la permission de l’imaginer. Si je me trompe sur un détail, corrige-moi dans ta tête : tu as toujours aimé corriger les schémas des autres.

La veille, je crois, le réseau n’a rien voulu charger. La page des résultats tournait sur l’écran du téléphone, la petite roue grise recommençait son tour, et chaque fois que tu relançais, quelqu’un dans la cour demandait « Alors ? ». Tu as fini par poser le téléphone face contre la natte. Tu as dit que tu irais voir au portail du lycée, le lendemain, comme on faisait avant qu’il y ait des pages.

Le matin, il pleuvait. Une pluie de juillet, lourde, qui ne s’excuse pas. Hadja a voulu venir. Elle a pris le grand parapluie de ton père, celui dont une baleine est tordue, et elle l’a tenu au-dessus de toi pendant tout le trajet, le bras levé le plus haut possible, parce qu’elle a douze ans et que tu la dépasses d’une tête. Tu avais de l’eau dans les sandales. Elle avait l’épaule gauche trempée et se taisait.

Devant le portail, la foule était déjà là : des parapluies, des sacs plastiques sur les têtes, des mères qui avaient laissé leur étal au marché. Les feuilles étaient punaisées sous l’auvent du gardien, et la pluie frappait la tôle si fort qu’on n’entendait plus que les prénoms criés par ceux qui se trouvaient. Tu t’es glissée jusqu’à la lettre B. Tu as lu la colonne de haut en bas, lentement, puis de bas en haut, plus lentement encore, comme si ton nom avait pu se cacher entre deux autres en te voyant arriver. Tu l’as lue une troisième fois. Hadja, maintenant, tenait le parapluie au-dessus de la feuille, pour que l’eau ne brouille pas l’encre.

Ton nom n’y était pas.

J’ignore ce que tu as dit. Rien, peut-être. Ou une de ces phrases sèches dont tu as le secret, sur le gardien, sur la pluie, sur la police de caractères de la liste, pour que Hadja ne regarde pas ton visage. Je sais seulement que vous êtes rentrées, et que la pluie sur la tôle de la maison a continué tout l’après-midi, avec la même indifférence que le matin.

À la maison, je suppose que ta mère a compris avant que tu parles, à la façon dont Hadja refermait le parapluie sans le secouer. Elle a arrêté la pédale de sa machine à coudre. Le tissu est resté sous l’aiguille, à moitié cousu, et personne n’a pensé à le retirer de la journée. Elle a sans doute dit une phrase maladroite, parce que le silence lui faisait plus peur que les mots.

Tu n’existes pas sous ce nom, Aïcha. Tu es faite de beaucoup d’élèves et de beaucoup de matins pareils, dans toutes les villes où l’on affiche des listes, à Douala et à Yaoundé comme à Garoua. Et pourtant c’est à toi que j’écris, à toi seule, parce qu’on n’écrit pas une lettre à une foule.

Je commence par une chose que personne ne t’a dite depuis hier, parce qu’autour de toi on cherche déjà des mots pour l’adoucir : c’est un échec. Je ne le repeindrai pas en chance déguisée ou en leçon de vie. Tu as passé le Bac D en mai, tu n’es pas admise en juillet. C’est arrivé, et ça fait mal.

Mais regarde-le bien : un échec a des contours. Il a une date, une session, une série, des coefficients, et bientôt un relevé de notes qui dira, matière par matière, ce qui s’est passé pendant quelques heures dans des salles surveillées. Il dit quelque chose de vrai : cette fois, à cette session, le total n’a pas atteint le seuil. Je ne te demande pas de trouver que c’est peu.

Ce qu’il tait est beaucoup plus large que ce qu’il dit. La liste ignore comment tu dessines une coupe de feuille, ou combien de fois tu as refait le schéma de la méiose pour que les chromosomes ne se chevauchent pas. Elle ne connaît ni tes nuits, ni ta façon de comprendre une chose en la dessinant. Ce que tu as lu au portail ne prouve pas ce que tu crois.

Tu es trop lucide pour qu’on te cache la suite. Tu as fait beaucoup de choses justes cette année. Tu as appris à dormir avant les épreuves au lieu de veiller jusqu’à 2 h. Tu as travaillé. Et ça n’a pas suffi. Les bonnes méthodes augmentent les chances ; elles ne les rendent pas certaines. Je n’ai jamais promis le contraire, et je ne commencerai pas aujourd’hui.

Alors, d’abord, tu as le droit d’être triste. Pas un peu, pas avec élégance : à la mesure de ce que tu as mis dans cette année. Sur ton calendrier mural, tu as barré les jours depuis septembre, un trait par soir, jusqu’à la date de la première épreuve. Après, les cases sont restées blanches. Il est normal qu’elles te regardent.

Je n’ai pas de phrase pour enlever ce poids, et méfie-toi de ceux qui en ont une. Ta tristesse a la taille exacte de ce que tu voulais ; tu n’as pas à la cacher à Hadja. Elle a vu combien de soirs tu as tenu la lampe allumée.

Pour cette semaine, une seule règle, et je te la donne comme on conseille la prudence sur une route mouillée : ne décide rien. Ni que tu arrêtes, ni que tu changes tout. Les décisions prises le lendemain d’une liste parlent avec la voix de la liste. Laisse passer sept jours, et sécher les sandales sur le seuil.

On va te poser des questions. Certains le feront par affection, d’autres pour comparer avec leurs propres enfants. Ton téléphone va vibrer de messages de félicitations adressés à d’autres, dans des groupes où tu es encore. Tu as le droit de couper le son. Tu as aussi le droit de répondre « je ne sais pas encore » à quiconque te demande ce que tu vas faire : la réponse est complète, elle n’a pas à s’excuser.

Au portail, tu as entendu crier des amies. Certaines t’ont cherchée du regard, puis ont baissé la voix en comprenant. On peut être contente pour elles et malheureuse pour soi dans la même minute ; l’un n’annule pas l’autre, et tu n’as pas à choisir celui qui fait meilleure figure.

Pendant ces sept jours, occupe-toi de ton corps comme d’un animal qu’on aime et qui a eu peur. Dors, même l’après-midi si la nuit refuse. Mange quelque chose de chaud, même sans faim. Sors : accompagne ta mère au marché, ou marche jusqu’au bord de la Bénoué à l’heure où la chaleur retombe, avec une amie qui ne te parlera pas de moyenne. Tu as appris cette année que dormir n’était pas tricher. C’est plus vrai encore maintenant qu’il n’y a plus rien à réviser.

Il arrive qu’après un échec la tristesse change de nature : la nuit ne finit plus, et une idée vient, celle de disparaître, de se faire du mal, de ne plus être là. Si cette idée vient, même une fois, même chassée aussitôt en se retournant sur la natte, alors ce soir même, sans attendre demain, tu en parles à quelqu’un. À ta mère, à ton père, à un adulte de confiance, à un médecin, à un centre de santé. Tu n’as pas besoin de trouver les bons mots ; « j’ai des pensées qui me font peur » suffit.

Et si tu ne sais pas vers qui aller, tu peux m’écrire ou m’appeler : +237 678 43 85 57, 689 72 19 23 ou ludovic@onbuch.site. Je ne suis ni médecin ni psychologue, et je ne remplace aucun de ceux-là. Mais tu n’auras pas à chercher seule : je te mettrai en contact avec quelqu’un de qualifié.

À ce point, la méthode et le courage n’ont plus rien à faire. C’est l’affaire du soin, et le soin se confie à d’autres que soi. Ton père passe ses journées dans un centre de santé : il sait mieux que moi que ces pensées-là se soignent, et qu’on n’en a pas honte.

Dans trois semaines, pas avant, il y aura un autre travail. Le relevé sera là, ou tu iras le chercher. Tu le poseras sur la table et tu le liras comme on lit un bulletin quand on a cessé d’y chercher un verdict : une ligne à la fois, avec son coefficient. Une note basse dans une matière lourde ne raconte pas la même histoire qu’une note basse dans une matière légère. Tu verras peut-être qu’il manquait peu, et où. Ou qu’il manquait beaucoup, dans une matière que tu croyais tenir. Les deux sont des informations, et tu sauras quoi en faire.

Ensuite seulement, tu regarderas les portes. Il y en a plusieurs, et je ne les classerai pas.

Tu peux redoubler la Terminale, dans ton lycée ou dans un autre, et reprendre l’année avec ce que tu sais maintenant de toi-même. Redoubler, c’est retrouver les mêmes murs avec des camarades plus jeunes, et un regard à soutenir les premières semaines ; c’est aussi connaître déjà le programme, et pouvoir enfin reprendre les marches qui ont cédé.

Tu peux te présenter en candidate libre, travailler de ton côté, en faisant autre chose en même temps pour respirer ou pour gagner un peu. Cela demande de se tenir soi-même, sans cloche ni appel, ce qui est plus dur qu’il n’y paraît, mais laisse du temps à soi.

Il existe aussi d’autres voies de formation, qui demandent d’autres choses et en donnent d’autres. Chaque porte a son prix, et ce prix varie d’une personne à l’autre. Chacune a ses exigences, et des gens qui la passent la tête haute.

Ces portes, ne les regarde pas seule. Parles-en avec ta mère, à côté de la machine à coudre, pendant qu’elle a les mains occupées et qu’il est plus facile de parler sans se regarder. Parles-en avec ton père, le soir, quand il rentre et qu’il enlève ses chaussures. Va voir aussi un professeur qui te connaît ; il y en a sans doute un, celui qui écrivait des remarques dans la marge de tes schémas. Tes parents ont eu peur, eux aussi, devant la liste, et la peur des parents prend parfois la forme d’un reproche. Écoute ce qu’il y a dessous avant de répondre à ce qu’il y a dessus.

On va te demander, dans les semaines qui viennent, ce que tu vas « devenir ». Le mot est lourd, comme si tu n’étais rien en attendant.

Tu dessines les schémas de SVT mieux que le manuel. Les manuels dessinent les cellules comme des briques bien rangées ; toi, tu leur laisses leur forme, leurs bords un peu mous, et on comprend en regardant ta page ce qu’on ne comprenait pas en regardant l’imprimé. C’est toi aussi, l’aînée, qui as appris à Hadja à lire l’heure sur l’horloge de la pièce, en lui expliquant que la petite aiguille est paresseuse et que la grande fait tout le travail. Elle le répète encore, je parie, à qui veut l’entendre. Et tu as pour chaque situation une phrase drôle et un peu méchante, ce qui est une forme d’intelligence que le Bac ne note pas.

Une vocation de rechange, je ne t’en trouverai pas. Ce serait facile, et ce serait faux : je ne te connais pas assez pour savoir ce qui te conviendrait, et personne ne le sait à ta place cette semaine. Tu te relèveras, ou pas tout de suite, à ton rythme, qui n’est celui de personne d’autre. Certains repartent le lundi suivant, d’autres ont besoin d’un mois de silence. Je ne mets pas de note à la vitesse.

Le calendrier mural, ne le jette pas tout de suite, et ne le garde pas non plus comme une relique au-dessus de ton lit. Un jour, tu en accrocheras un autre, ou tu n’en accrocheras plus. Les cases blanches de juillet restent des cases, sans rien annoncer.

Tu voulais faire pharmacie. Tu le veux peut-être encore. Prétendre que le diplôme compte peu serait un mensonge confortable, et tu le flairerais tout de suite. Il compte. Il ouvre des portes que rien d’autre n’ouvre, et il te donnerait le droit d’entrer là où tu voulais entrer. Tu as le droit de le vouloir fort, et de le vouloir encore. Personne n’a à te demander de « relativiser » la semaine où tu ne l’as pas.

Et pourtant, une vie ne se lit pas sur une seule liste. Je me méfie du mot « réussir » quand on le prononce au singulier, comme s’il désignait une case qu’on coche une fois pour toutes. Ce qui ressemble à une réussite, quand on regarde une vie de près, tient plutôt d’une suite de petites choses tenues, à des dates différentes, dont aucune n’a été affichée au portail. Je ne te demande pas d’y croire aujourd’hui. Je te demande seulement de ne pas laisser une feuille punaisée sous un auvent décider seule de ce que vaut ta vie.

Quelqu’un, dans cette histoire, ne quitte ni ma pensée ni la tienne : Hadja. Elle tenait le parapluie, elle t’a vue lire la liste trois fois. Elle t’a regardée tomber, et je crois qu’elle te regarde encore, depuis la porte, depuis la natte, avec ces yeux de petite sœur de douze ans qui notent tout et ne disent rien.

Elle apprendra quelque chose de ce mois de juillet, et pas forcément ce que tu crains. Elle apprendra autant de la façon dont tu te relèves qu’elle aurait appris de la façon dont tu aurais réussi. Davantage, peut-être : une réussite se montre de loin, sur une photo ; la manière de se relever se voit dans la cuisine, à l’heure des repas. Se relever, devant elle, peut vouloir dire pleurer, puis manger, puis dormir, et un jour rouvrir un cahier. Elle verra que c’est possible, et elle s’en souviendra le jour où ce sera son tour d’attendre devant un portail.

Dans ce livre, à tous ceux qui le lisent, je propose chaque fois quelque chose à faire le soir même. À toi, je dis le contraire. Ce soir, n’essaie rien. N’ouvre pas le cahier. Ne refais pas le calcul des coefficients pour savoir de combien de points tu l’as manqué, ne relis pas les sujets, ne dresse aucune liste de résolutions. Mange ce qu’il y a dans la marmite, puis couche-toi tôt, même si le sommeil tarde, et laisse-le venir quand il voudra. C’est la seule consigne de ce livre qui consiste à ne pas travailler, et je la donne aussi sérieusement que les autres.

La vraie consigne, je la repousse. Dans trois semaines, quand le relevé sera sur la table, ouvre-le, et ouvre aussi ton cahier de brouillon à une page neuve. Pour chaque matière, écris deux phrases : ce que la note dit, et ce qu’elle ne dit pas. La première doit être honnête, même si elle pique. La seconde doit être précise, même si elle te paraît dérisoire. Par exemple : « Ma note de physique dit que, ce jour-là, je n’ai pas su mener l’exercice sur les circuits jusqu’au bout. Elle ne dit pas que je ne comprends rien à la physique. » Puis ferme le cahier, et va dire à quelqu’un ce que tu as trouvé.

Si l’année prochaine sera la bonne, je n’en sais rien, et toi non plus. C’est d’ailleurs pour cela que je t’écris maintenant plutôt que dans un an : ce que j’ai à te dire doit tenir quoi qu’il arrive ensuite.

Demain matin, s’il ne pleut pas, quelqu’un montera poser la lampe solaire sur la tôle, comme tous les jours, et elle chargera jusqu’au soir. À la tombée de la nuit, Hadja viendra te demander si elle peut l’emprunter pour faire ses devoirs. Tu diras oui.

Ludovic


%% 18-epilogue-le-mot-encore.md
\chapitresn{Épilogue — Le mot, encore}

Ebolowa, juillet, le soir. Il a plu dans l’après-midi et la cour sent encore la terre mouillée. Rodrigue a eu son BEPC, sans mention, de peu : il l’a lu d’abord sur le téléphone d’un voisin qui avait encore du forfait, puis le lendemain sur la liste du portail, où il a posé le doigt sous son nom comme pour l’empêcher de glisser. Il n’y a pas eu de fête. Mama Odile a dit « bon », ce qui chez elle pèse lourd, et elle est partie au marché comme les autres jours. À la rentrée, il entrera au lycée technique, en électronique ; elle a signé la fiche le soir où ils sont revenus de la visite, et elle n’a plus reparlé de bureaux.

Ce soir, elle est à la réunion de sa tontine. Le courant est parti à 21 h, avec le soupir habituel du quartier. Sur la table de la cuisine, le téléphone de Rodrigue tient debout contre un verre d’eau, lampe torche allumée vers le plafond, et la lumière retombe en rond sur deux cahiers. Celui de Bella est un cahier de vacances acheté au marché, ouvert à la page des fractions. Elle a dix ans, des tresses qu’on lui a faites dimanche, et une patience qui finit toujours au même endroit. Elle raye, recommence, raye encore. Puis elle pose son crayon à plat, le pousse du bout des doigts jusqu’au milieu de la table, et dit : « Je suis nulle. »

Rodrigue ne répond pas tout de suite.

Il entend autre chose. Une salle de cinquième, la chaleur sous la tôle, une voix qui lit des notes très vite parce qu’il lui reste deux classes. Il entend surtout le rire, le sien, celui qui partait le premier pour que les autres tombent dessus. Le rire remonte, tout prêt, la vieille manière de faire passer le mot pour une blague. Rodrigue le laisse redescendre, sort le stylo de la poche de sa chemise, à côté du petit tournevis, et ouvre son cahier de brouillon.

\asterisme

Il écrit, en haut d’une page propre : nul. Bella regarde le mot, puis son cousin. « Tu vois, toi aussi tu le dis. » Il répond « attends », et trace à côté du mot un signe égal, puis un zéro. « En maths, nul, ça veut dire zéro. Un nombre, comme les autres, qu’on met dans les calculs. » Elle hausse les épaules, comme on les hausse devant un grand qui fait le malin.

Il écrit en dessous : 7 × 5 × 0. « Combien ? » Elle calcule à mi-voix, trente-cinq, s’arrête, et dit « zéro » sur le ton de quelqu’un qui flaire un piège. « Et si je mets mille à la place du sept ? » Zéro encore. « Voilà. » Il tapote la ligne avec le stylo. « Il suffit d’un seul zéro dans la multiplication, et tout le résultat a l’air nul. Mais le sept, il est toujours là. Le mille aussi. Ils ne sont pas devenus zéro, eux. » En février, en classe, il a recopié cette propriété sans y voir autre chose qu’une ligne de plus dans la leçon. Il la retrouve ce soir dans sa propre bouche, avec des mots qui ne sont pas ceux du professeur et qui boitent un peu. « Alors quand ça donne zéro, on ne jette pas tout le calcul. On cherche lequel est à zéro. »

Bella réfléchit. « Et moi, c’est lequel ? »

Il compte sur ses doigts, sérieux comme tonton Fabrice devant un écran noir. « Tu as mangé ? » Elle a mangé. « Il fait noir. » Il fait noir, elle l’accorde. « Et personne ne t’a montré les parts. Ça fait déjà deux. » Il ajoute, sans la regarder : « Moi aussi j’en avais, des zéros dedans. Plusieurs. On ne les a pas encore tous trouvés. »

Il se lève, prend dans le sac de Mama Odile une poignée d’arachides, celles de demain pour le marché, et en compte douze sur la table, dans le rond de lumière. « Tata Odile ne pose jamais d’opération. Elle partage. » Il fait deux tas de six. « La moitié, c’est combien ? » Six. Il défait un tas en deux. « Le quart ? » Trois. « Alors une moitié plus un quart ? » Bella pousse les arachides avec l’index, un tas, puis l’autre, et compte. Neuf. « Neuf sur douze », dit Rodrigue, et il lui montre que neuf sur douze, en regroupant par trois, fait trois quarts, ce que son cahier de vacances réclame depuis le début de la soirée. Elle vérifie deux fois. Puis elle mange une arachide.

« Onze sur douze », dit Rodrigue. Elle rit, la bouche pleine. Ensuite, elle lève la main, bien droite, au-dessus de la table de la cuisine. Il lui rappelle qu’on n’est pas en classe. Elle garde la main en l’air jusqu’à ce qu’il dise, avec la voix d’un censeur : « Oui, Bella ? » Et elle demande ce qui se passe si on partage en zéro parts. Rodrigue ouvre la bouche, la referme. « Ça, c’est une autre leçon. » Il n’est pas sûr de savoir la faire.

\asterisme

Je ne crois pas qu’on guérisse d’un mot. Celui-là entre trop vite, trop jeune, par une porte qu’on n’avait pas eu le temps de fermer, et ensuite il connaît la maison. Mais on peut le déplacer. Le décoller de l’endroit où quelqu’un l’a mis, sur un nom, et le reposer là où il a un sens : sur un résultat, sur une copie, sur une marche, sur un soir sans lumière. Ce n’est pas une morale, c’est une opération, de celles qu’on fait au brouillon quand on s’aperçoit qu’on a écrit la bonne chose à la mauvaise ligne. Un quatre en cinquième, un « je suis nulle » à dix ans devant un cahier de vacances, ne prouvent pas ce que tu crois. Ils disent qu’un facteur, quelque part, est à zéro ce soir-là. On peut le chercher. On ne le trouve pas toujours du premier coup.

\asterisme

Le courant revient d’un seul coup. Le cri monte de toutes les maisons à la fois, ce cri que personne n’a appris et que tout le quartier pousse ensemble, puis le frigo du voisin repart en toussant, une radio reprend au milieu d’une chanson. Par la fenêtre, dans la lumière retrouvée, l’escalier de béton d’à côté a une rampe depuis le mois dernier, une barre de fer peinte en vert, et les deux marches manquent toujours. Rodrigue éteint la lampe torche. Bella reprend son crayon.

Plus tard, quand elle dort, il fouille le carton sous son lit, celui des vieux cahiers que Mama Odile garde parce qu’on ne jette pas le papier. Le cahier de brouillon de cinquième est en dessous, la couverture gondolée par une ancienne pluie. L’étiquette n’a pas bougé. À la ligne « Nom : », son prénom et son nom barrés d’un trait, et par-dessus, au stylo à bille, en capitales, le mot. Il le regarde un moment sans rien en penser de particulier. Puis il glisse la pointe du tournevis sous un coin et décolle l’étiquette, lentement, pour qu’elle ne se déchire pas. Il la garde, et la recolle à l’intérieur, sur la dernière page, au milieu de ses vieux calculs, où elle n’est plus qu’un mot parmi des nombres.

Il lui faut une étiquette neuve. Il n’y a que la planche de Bella, celle de la rentrée dernière, avec un papillon rose imprimé dans le coin. Il en prend une, la colle bien droite sur la couverture, l’aplatit du pouce. Il prend le stylo de Bella, celui qui a un pompon. À la ligne « Nom : », avec ce qu’il a, il écrit son nom.

\begin{pourcahier}{Pour le cahier}
Retourne à la première page. Relis la phrase que tu y as écrite au début de ce livre. Ne la rature pas. En dessous, réécris-la à sa juste place : ce qu’elle disait d’un moment, d’une copie, d’une marche. Pas de toi.
\end{pourcahier}


%% 19-note-de-l-auteur.md
\chapitresn{Note de l’auteur}

Les prénoms et les parcours de ce livre sont composés à partir de situations réelles et typiques ; aucun n’est une personne précise. Rodrigue, Grâce, Divine, Ibrahim, Nadège et Aïcha, monsieur Kamga et Mama Odile sont faits de morceaux : une scène qui revient dans beaucoup de classes, un geste qu’on retrouve dans beaucoup de maisons, une phrase que des parents différents disent avec les mêmes mots. Je leur ai donné une ville, un âge, une famille et un détail, pour qu’ils tiennent debout ; je ne leur ai donné aucun modèle. Si tu crois te reconnaître dans l’un d’eux, c’est que la situation était typique, et c’est justement pour cela qu’elle méritait d’être écrite.

Ce livre s’appuie aussi sur des travaux de psychologie cognitive et de sciences de l’éducation. J’ai essayé de les rapporter avec prudence, en disant chaque fois si un résultat me paraît solide, discuté ou encore fragile, et, quand je n’étais pas sûr d’une référence, en gardant l’idée sans le nom. La recherche avance et se corrige : certaines choses dites ici seront nuancées demain, et c’est plutôt une bonne nouvelle, puisque c’est la preuve qu’on continue de vérifier.

Si tu veux aller plus loin, voici les travaux cités en chemin, regroupés par thème, à lire de préférence avec quelqu’un qui peut t’aider à les situer.

Sur les notes et le regard des autres : Henri Piéron, qui a forgé le mot « docimologie » dès les années 1920 et a fait le bilan de ces recherches dans \emph{Examens et docimologie} (1963) ; Robert Rosenthal et Lenore Jacobson, \emph{Pygmalion in the Classroom} (1968), traduit en français sous le titre \emph{Pygmalion à l’école}, à lire avec la synthèse critique de Lee Jussim et Kent Harber (2005) ; Thomas Gilovich, Victoria Medvec et Kenneth Savitsky (2000) pour l’effet projecteur ; Leon Festinger (1954) pour la comparaison sociale ; Herbert Marsh et ses collègues, à partir de 1984, pour l’effet du gros poisson dans une petite mare.

Sur les bases et la peur : Keith Stanovich (1986), qui a appliqué à la lecture l’« effet Matthieu » nommé par le sociologue Robert Merton (1968) ; Mark Ashcraft et Elizabeth Kirk (2001) sur l’anxiété mathématique ; Sian Beilock et Thomas Carr (2005) sur la performance sous pression.

Sur la voix de dedans : Martin Seligman et Steven Maier (1967) pour l’impuissance apprise, et leur révision de la théorie, \emph{Learned Helplessness at Fifty} (2016) ; Ethan Kross et ses collègues (2014) sur le fait de se parler à soi-même par son prénom ou en « tu » ; les travaux de Carol Dweck sur la mentalité de croissance, puis, pour leurs limites, la méta-analyse de Victoria Sisk et de ses collègues (2018) et la grande étude de David Yeager et de ses collègues (2019).

Sur la manière d’apprendre : John Sweller et la théorie de la charge cognitive, pour les exemples résolus ; Michelene Chi et ses collègues (1989) pour l’auto-explication ; et, pour en finir avec les styles d’apprentissage, l’article de Harold Pashler et de ses collègues (2008). Puis Hermann Ebbinghaus (1885) pour la courbe de l’oubli ; Henry Roediger et Jeffrey Karpicke (2006) pour l’effet de test ; la revue de John Dunlosky et de ses collègues (2013), qui compare les techniques d’étude ; la méta-analyse de Nicholas Cepeda et de ses collègues (2006) sur l’espacement des révisions ; les travaux de Robert Bjork sur les difficultés désirables ; et Sebastian Leitner, \emph{So lernt man lernen} (1972), pour la boîte qui porte son nom.

Sur le sommeil et l’attention : la revue de Susanne Diekelmann et Jan Born (2010), les travaux de Mary Carskadon sur les adolescents, et l’étude d’Adrian Ward et de ses collègues (2017) sur la présence du téléphone, dont l’effet n’a pas toujours été retrouvé depuis.

Sur l’erreur et l’examen : Brady Butterfield et Janet Metcalfe (2001) pour l’effet d’hypercorrection ; John Hattie et Helen Timperley (2007) sur le retour d’information ; Jeremy Jamieson et ses collègues (2010) sur la façon de lire les signes du stress ; Gerardo Ramirez et Sian Beilock (2011) sur l’écriture avant un examen, dont l’effet n’a pas été retrouvé lors d’une grande tentative de réplication en 2018 (Colin Camerer et ses collègues).

Sur le monde autour de l’école, enfin : \emph{Les Héritiers} de Pierre Bourdieu et Jean-Claude Passeron (1964), les travaux de Sendhil Mullainathan et Eldar Shafir sur la rareté (2013), la théorie de l’autodétermination d’Edward Deci et Richard Ryan (1985), et le début du \emph{Manuel} d’Épictète, qui est court et se lit encore très bien.

Aucune de ces lectures n’est nécessaire pour faire ce que propose ce livre. Un cahier de brouillon suffit.

Je suis par ailleurs le créateur d’OnBuch (OnBuch+), une application de cours construite sur les programmes officiels du MINESEC ; le chapitre « La lampe » dit pourquoi je la fais, et ce qu’elle ne fait pas.

J’ai écrit ce livre parce que mon cœur n’en pouvait plus : il me rongeait d’envie de faire quelque chose pour tous les élèves que je dois aider. Je ne peux pas être dans chaque classe ni à côté de chaque lampe ; un livre, lui, peut y être.

Pour m’écrire à propos de ce livre ou d’OnBuch : +237 678 43 85 57, 689 72 19 23 ou ludovic@onbuch.site.


%% 20-remerciements.md
\chapitresn{Remerciements}

Merci aux enseignants qui tiennent : à ceux qui corrigent la dernière copie de la pile avec le même soin que la première, et qui, devant un regard qui ne comprend pas, cherchent une deuxième façon d’expliquer. Ils sont fatigués, et ils reviennent le lendemain.

Je pense aussi aux parents qui veillent sans toujours savoir comment aider. Ils s’assoient à côté de la lampe sans connaître la leçon, ils demandent si on a mangé, et leur présence, même silencieuse, compte plus qu’ils ne le croient.

Et aux élèves qui ont posé leurs questions, y compris celles qu’ils trouvaient bêtes, à voix haute, sur un bout de papier ou à la sortie du cours : chacune de ces questions a servi à quelqu’un d’autre qu’eux.

Merci à la ministre des Enseignements secondaires, et à son ministère : les programmes officiels qu’ils publient sont le cadre sur lequel OnBuch est construit.

Merci enfin à l’équipe OnBuch, qui construit avec moi, jour après jour, ce que ce livre ne fait que raconter.


%% 21-postface-parents-enseignants.md
\chapitresn{Postface — Pour les parents et les enseignants}

À Ebolowa, une nuit de novembre, une femme est assise à la table de la cuisine, à côté de son fils. Il bute sur une addition de fractions. Elle ne saurait pas la faire, elle le sait, elle ne fait pas semblant. Elle a devant elle les comptes de sa tontine, qu’elle tient de tête, sans poser une opération, et elle reste là jusqu’à ce qu’ils soient justes. Elle ne lui demande pas s’il a fini. Elle ne regarde pas l’heure. Quand il lui demande comment elle fait pour ne jamais se tromper dans les parts, elle hausse les épaules et répond qu’elle partage. Elle n’a rien enseigné ce soir-là, sinon qu’elle était là, et le garçon travaille autrement quand quelqu’un est assis à côté de lui.

À Bafoussam, un homme d’un peu plus de cinquante ans calcule des moyennes sous une lampe de bureau dont il faut tenir le fil d’une certaine façon pour qu’elle reste allumée. Plus de soixante-dix noms dans une seule classe. Il connaît la leçon, lui, mieux que personne dans la maison ; ce qui lui manque, c’est le temps. Pourtant, dans un petit cahier à couverture verte, il écrit pour chaque élève une chose réussie dans le trimestre. On ne lui a rien demandé. Aucun bulletin ne le reproduira.

Ces deux personnes ne se connaissent pas. L’une ne peut pas aider sur la leçon, l’autre ne peut pas s’asseoir à côté de chacun. Elles font pourtant, chacune de son côté, une chose qu’aucune note ne mesure : elles regardent un enfant comme quelqu’un qui est en train d’apprendre.

Ces quelques pages sont pour vous qui êtes à cette table ou derrière ce bureau. Le livre que vous tenez tutoie un élève ; ici, je vous vouvoie, parce que vous êtes plusieurs et que je n’ai pas à vous apprendre votre métier de parent ou d’enseignant. Je vous confie seulement quelques gestes, qui ne coûtent rien, et que vous faites peut-être déjà.

\asterisme

Le livre commence par un mot dit en une seconde dans une salle de classe, et qu’un élève garde ensuite pendant des années. Ce mot, presque tous les adultes l’ont prononcé un jour, sous une forme ou une autre, sans y penser, un soir de fatigue. Je n’en fais le procès de personne. Je remarque seulement l’écart entre celui qui parle et celui qui reçoit : l’un a dit une phrase, l’autre a entendu une définition.

Que dire à la place ? Quelque chose de plus exact, qui peut rester ferme. Nommer la tâche. « Cet exercice n’est pas encore acquis » dit la vérité sur la copie et laisse la suite ouverte ; « tu es nul » prétend dire la vérité sur l’enfant et ferme tout. « Tu as raté les accords du participe passé » est une phrase dure, et elle est utile, parce qu’elle désigne un endroit où travailler. On peut être sévère avec un exercice ; c’est même là que la sévérité sert.

Les questions qu’on pose comptent aussi. « Tu as appris ta leçon ? » appelle un oui, et le oui ne prouve rien : l’enfant a souvent relu, il croit savoir, il est de bonne foi. Essayez plutôt : « Explique-moi ce que tu as compris. » Cette question a un avantage que je trouve précieux pour les parents qui n’ont pas fait de longues études ou qui ont oublié la chimie de leur jeunesse : vous n’avez pas besoin de connaître la leçon pour l’entendre. Vous entendez très bien quand une explication tient debout et quand elle s’effondre au milieu d’une phrase. Mama Odile ne savait pas additionner deux tiers et un quart à la manière de l’école ; elle aurait parfaitement su entendre son fils hésiter.

\asterisme

Aux parents, encore, quelques propositions.

Protégez le sommeil. Je sais combien c’est difficile dans une maison où l’on vit nombreux, où il fait chaud, où le courant part à l’heure où l’on voudrait travailler. Je sais aussi qu’un enfant qui veille jusqu’à 2 h du matin avant un examen ressemble à un enfant sérieux, et qu’on a envie de l’en féliciter. Le livre lui explique, avec prudence, pourquoi la nuit blanche coûte en général plus qu’elle ne rapporte. Il l’écoutera mieux si, à la maison, quelqu’un lui dit qu’il a le droit d’éteindre, et que dormir n’est pas tricher.

Quand le bulletin arrive, lisez-le avec votre enfant, ligne par ligne, avant d’arriver au rang en bas de la page. Le rang dépend autant des autres élèves que de lui ; il peut cacher un seize en mathématiques au milieu d’une moyenne ordinaire, ou un quatre devenu huit, qui est un vrai progrès et que personne ne voit. Demandez-lui, matière par matière, ce que la note résume et ce qui s’est passé. Dix minutes de cette conversation vous en apprendront plus qu’une semaine de reproches.

Il arrive qu’on dise une phrase dure parce qu’on a peur : peur de l’avenir, des frais, du regard de la famille, peur que l’enfant ne finisse comme on redoute. Ces phrases-là sortent avant qu’on ait fini de les penser. On peut y revenir le lendemain. Un parent qui dit « hier, j’ai mal parlé, j’avais peur pour toi » garde toute son autorité. Il a appris à son enfant une chose que l’école enseigne rarement : une parole se répare.

Enfin, et je le dis sans vous accuser de rien, beaucoup d’entre vous travaillent tard, au marché ou sur la route, et rentrent quand les devoirs sont finis ou abandonnés. Personne ne peut s’asseoir chaque soir à la table de la cuisine. Prenez ce que je décris comme une offre, à saisir quand vous le pouvez. Une fois par semaine, une question bien posée, un bulletin lu ensemble : c’est déjà beaucoup.

\asterisme

Aux enseignants, je dois d’abord ce que le livre essaie de dire tout du long, parce qu’on le dit trop peu : je sais dans quelles conditions vous travaillez. Les classes où l’on s’assoit à trois par banc, les piles de copies qui ne diminuent pas, les programmes à finir, les salaires qui obligent beaucoup d’entre vous à donner des répétitions le samedi, les lampes qui s’éteignent quand on touche le fil. Ceux qui tiennent dans ces conditions, et ils sont nombreux, méritent mieux que des leçons de pédagogie venues de l’extérieur de leur classe. Je n’en donnerai pas.

Je me permets seulement deux remarques, que vous connaissez mieux que moi.

Un élève lit rarement une correction entière ; il lit la note, puis, s’il y en a un, le commentaire. Un commentaire précis, même d’une ligne, change ce que la copie devient au fond du sac. « Raisonnement juste jusqu’à la ligne quatre ; revois le changement de signe » donne à l’élève un endroit où poser le pied, quand « insuffisant » le laisse dans le vide. Je sais que la ligne précise prend une minute de plus et que soixante-dix minutes, c’est une soirée. Il arrive aussi qu’une même erreur revienne sur une trentaine de copies ; elle dit alors moins quelque chose de trente élèves que d’une leçon à reprendre, et vous êtes les premiers à le voir.

Dans une classe pléthorique, être appelé par son nom est déjà une forme d’attention. Beaucoup d’élèves de ce livre ont peur d’être vus ; ils ont plus peur encore de ne l’être jamais. Je ne vous demande pas de tenir un cahier vert comme monsieur Kamga. Le sien montre seulement qu’une chose réussie, notée quelque part par un adulte, peut exister même quand le bulletin ne la reproduit pas. Si un jour vous la dites à l’élève concerné, elle pèsera autant que bien des notes.

\asterisme

Une dernière chose sort du domaine de la méthode.

Il arrive qu’un enfant ou un jeune change sans bruit. Il ne dort plus, ou dort tout le temps. Il ne mange plus, s’isole, laisse tomber ce qu’il aimait, ne rit plus de ce qui le faisait rire. Il parle de lui avec une dureté nouvelle, ou évoque l’idée de disparaître, de ne plus être là, parfois sous forme de plaisanterie. On prend facilement ces signes pour de la paresse ou un caprice. Ils demandent qu’on s’arrête, qu’on pose la question franchement, et qu’on écoute la réponse sans juger ni minimiser. De nombreuses études l’indiquent : parler de ces idées avec lui ne les lui donne pas, et le silence le laisse seul avec elles.

Ni un parent ni un enseignant ne peut porter cela seul. Il faut orienter, et vite : vers un médecin, un centre de santé, un professionnel de la santé mentale quand il en existe un à portée, l’infirmerie ou le conseiller d’orientation de l’établissement, un responsable religieux ou communautaire de confiance si c’est vers lui que l’enfant se tourne. Le même devoir s’impose quand un élève laisse paraître qu’il subit des coups, des humiliations répétées, des menaces, à l’école ou ailleurs. Un conseil de révision ou un sermon ne règlent rien ici. On confie ces situations à ceux dont c’est le métier, et on reste présent pendant ce temps. La dernière page de ce livre, « Si ça ne va pas », s’adresse directement aux élèves ; elle est faite pour être montrée.

\asterisme

À Ebolowa, la nuit avance. Le garçon a fini, ou il a abandonné, on ne sait pas. Sa mère a rangé les comptes de la tontine sous la nappe en plastique, avec les papiers importants. Elle n’a pas fait l’exercice. Elle n’a pas expliqué les fractions. Elle est restée assise jusqu’à la fin de ses comptes, à côté de lui, et c’était déjà beaucoup.

C’est peut-être la place des adultes, dans ce livre et ailleurs : celle qui éclaire la marche, pas celle qui la monte.


%% 22-si-ca-ne-va-pas.md
\chapitresn{Si ça ne va pas}

Cette page est la dernière du livre, et elle n’a rien à t’apprendre sur les notes.

Il arrive que la fatigue ou la tristesse dure plus longtemps qu’une mauvaise semaine. Le sommeil ne vient plus, ou il ne te lâche plus. Tu ne manges plus vraiment. Ce que tu aimais ne te fait plus rien. Ta voix de dedans ne s’arrête plus, même quand le cahier est fermé. Il arrive aussi qu’une idée vienne, celle de disparaître, de te faire du mal, de ne plus être là. Ou que quelqu’un te frappe, te menace, t’humilie, à l’école ou ailleurs.

Si c’est ton cas, la méthode ne suffit plus, et tu n’as pas à t’en sortir sans aide. Parles-en à quelqu’un, aujourd’hui si tu peux, et en tout cas sans attendre que ça passe. À un parent, à un frère ou une sœur aînée, à un adulte de confiance : un professeur, le surveillant général, l’infirmière ou le conseiller d’orientation de ton établissement. À un médecin, ou dans un centre de santé, où l’on sait que ces choses se soignent et où l’on n’a pas à en avoir honte.

Tu n’as pas besoin de trouver les bons mots. « Ça ne va pas, et je ne sais pas quoi faire » suffit. « J’ai des pensées qui me font peur » suffit aussi. Si la première personne à qui tu parles ne comprend pas, n’en conclus rien sur les autres : va voir la suivante.

Et si le danger est immédiat, si tu sens que tu pourrais te faire du mal ce soir, fais-toi accompagner tout de suite : réveille quelqu’un dans la maison, ou va vers le centre de santé le plus proche.

Tu peux aussi joindre l’auteur de ce livre : +237 678 43 85 57, 689 72 19 23 ou ludovic@onbuch.site. Ce n’est ni un service d’urgence ni un service de soin : si le danger est immédiat, va d’abord vers l’adulte ou le centre de santé le plus proche. Mais si tu ne sais pas vers qui aller, écris ou appelle : je te mettrai en contact avec quelqu’un de qualifié.

Tu peux montrer cette page à quelqu’un, plutôt que de tout expliquer. C’est aussi à cela qu’elle sert.

\end{document}
